\section{Concordance}\label{app:conc}

\subsection*{Preface}

\paragraph{What is enumerated.} This appendix is a concordance of EXI-01-a in the corpus of its discipline. Its object is every word type that occurs in the title and the eleven sentences of EXI-01-a, \cnTypes{} types in all, of which \cnContent{} are content words and the rest function words. Its corpus is the source text of the \cnOtherMss{} other manuscripts of the discipline (EXI-01-c to EXI-20-b; the present work is excluded), \cnCorpusWords{} words after the removal of markup. It has three parts. Table~A.1 lists every word type of EXI-01-a with the lemmata in which it occurs, its number of occurrences in the corpus, the number of manuscripts in which it occurs, and, for content words, its number of occurrences in each manuscript. Table~A.2 counts fifteen term families, defined by the patterns listed below, in every manuscript including EXI-01-a. Section~A.3 is a keyword-in-context concordance: every occurrence in the corpus of every content word of EXI-01-a, \cnKwic{} lines in all, with forty characters of context on each side, the manuscript and the section in which it occurs.

\paragraph{The program and the encoding.} The program \texttt{gen/generate.py} (Python~3, standard library only, deterministic) reads each manuscript's \texttt{main.tex} and \texttt{main.aux}. It keeps the text between the beginning and end of the document, removes the bibliography, code listings, comments and mathematics, resolves cross-references to their numbers using the manuscript's own \texttt{.aux} file and citations to author and year using its own bibliography, keeps the arguments of text commands such as \texttt{\textbackslash emph}, and removes all other commands. It divides the text into sections, numbered as in the manuscript (\enquote{Abs} for the abstract, letters for appendices, 0 for material before the first section), and the sections into sentences. A word is a maximal string of letters, possibly with internal apostrophes or hyphens; identifiers of manuscripts count as single words. Words are compared after conversion to lower case. For EXI-01-a, the program uses the lemmata of Section~\ref{s:text}, with citations removed, so that the occurrences in EXI-01-a are recorded by lemma number. The function words, listed in the program, are given totals only; five words that name statements of EXI-01-a or the present work (\emph{definition}, \emph{assumption}, \emph{theorem}, \emph{commentary} and the identifier of this work) are indexed in Table~A.1 but have no lines in A.3, since their occurrences in the corpus are references to statements rather than uses of terms. The families of Table~A.2 are defined by these patterns, matched against whole words: \textit{mean-} (\emph{mean}, \emph{means}, \emph{meaning}, \emph{meanings}, \emph{meaningful}, \emph{meaningfully}, \emph{meaningfulness}, \emph{meaningless}, \emph{meaningly}), \textit{life-} (\emph{life}, \emph{lives}, \emph{lived}, \emph{living}), \textit{engag-} (every word beginning \emph{engag}), \textit{activ-} (\emph{active}, \emph{actively}, \emph{activity}, \emph{activities}), \textit{project-} (\emph{project}, \emph{projects}), \textit{worth-} (\emph{worth}, \emph{worthwhile}, \emph{worthy}, \emph{worthiness}), \textit{objectiv-} (\emph{objective}, \emph{objectively}, \emph{objectivity}), \textit{valu-} (every word beginning \emph{valu}), \textit{attitud-} (\emph{attitude}, \emph{attitudes}), \textit{agent-} (\emph{agent}, \emph{agents}, \emph{agency}), \textit{availab-} (every word beginning \emph{availab}), \textit{degree-} (\emph{degree}, \emph{degrees}, \emph{extent}), \textit{matter-} (\emph{matter}, \emph{matters}, \emph{mattered}, \emph{mattering}), \textit{absurd-} (every word beginning \emph{absurd}) and \textit{sceptic-} (every word beginning \emph{sceptic} or \emph{skeptic}).

\paragraph{Checks.} Before writing the tables, the program checks that every lemma in its list occurs verbatim in EXI-01-a (the title against the title), that every phrase lemma occurs in its sentence lemma, that the sentence lemmata occur in the order of the text, that every lemma quoted in the chapters of this volume is quoted exactly and under its correct number, that every lemma in the list is quoted at least once, and that every listing of Lean code in Section~\ref{s:lean} occurs in the module. It halts if any check fails. The counts are exhaustive for the corpus as defined: every word of every sentence of every manuscript of the discipline other than this one is examined. Two limits should be stated. Text in generated tables included by a manuscript from separate files, which occurs only in EXI-10-b, is not part of the corpus, since it summarizes the work of other authors rather than the claims of the discipline. And the segmentation into sentences is mechanical, so that a line of A.3 may occasionally begin inside a sentence that contains an abbreviation.

\paragraph{Summary.} Of the \cnContent{} content words of EXI-01-a, \cnUnshared{} do not occur in the corpus: \cnUnsharedList. Three of these are the phrase \emph{sub specie aeternitatis}, which no other manuscript of the discipline uses, though EXI-04-a analyzes the standpoint it names under the name \enquote{detached standpoint}. The content words with the most concordance lines are \cnTopKwic. In the corpus, the family \textit{mean-} occurs \cnFammean{} times, \textit{engag-} \cnFamengag{} times, \textit{worth-} \cnFamworth{} times and \textit{objectiv-} only \cnFamobjectiv{} times.

\paragraph{Entries discussed.} Five entries bear on the commentary. (i) The entry \emph{objective} in Table~A.1 records \cnFamobjectiv{} occurrences of the family in the corpus against \cnFamworth{} of \textit{worth-}: the discipline speaks of the objective component of Definition~1 mainly through \emph{worth} and \emph{worthwhile}, the vocabulary of EXI-17-a's Definition~2.2, and rarely through the word \emph{objective}. The lines for \emph{objective} in A.3 show that most of its occurrences are in EXI-01-c, EXI-01-d and EXI-04-a, the manuscripts that discuss Assumption~2 directly. (ii) The entry \emph{project} records its occurrences as concentrated in EXI-01-c, EXI-03-a and EXI-03-b and in EXI-09-a and EXI-09-b. The lines in A.3 for EXI-09-a and EXI-09-b show the word in the sense of the defining project, which the comment on \Lm{2.6} distinguished from the projects of Definition~1; the lines for EXI-03-a show it in the sense of the projects for which categorical desires are desires, which is the sense of Definition~1. (iii) The entry \emph{matters} has fourteen of its occurrences in EXI-03-a, where the lines in A.3 show the phrase \enquote{what matters in survival}, a sense unrelated to the \enquote{nothing matters} of \Lm{11.3}; the concordance thus separates two senses that a count of the word would conflate. (iv) The entry \emph{sceptical} records two occurrences: one in EXI-04-a, in the sentence that quotes Remark~5 (see Appendix~\ref{app:test}), and one in EXI-08-a. The second concerns arguments about authenticity and is unrelated to the challenge of \Lm{11}. (v) The entry \emph{everyone} has five of its occurrences in EXI-01-d and the rest in EXI-02-a, EXI-07-a, EXI-11-a, EXI-17-a and EXI-20-a; the lines in A.3 for the last three are statements of the availability claim of EXI-01, and they show that the discipline uses the family conclusion's word, not Theorem~4's \enquote{every such agent}, when it reports the claim.

\subsection*{A.1 Index of the words of EXI-01-a}

{\footnotesize Columns: the word; the sentence lemmata of EXI-01-a in which it occurs (with multiplicity); its total occurrences in the other \cnOtherMss{} manuscripts; the number of those manuscripts in which it occurs; its occurrences by manuscript, with the prefix \enquote{EXI-} omitted.\par}

{\scriptsize% Generated by gen/generate.py; do not edit.
\begin{longtable}{@{}>{\raggedright}p{0.13\textwidth}>{\raggedright}p{0.1\textwidth}r r>{\raggedright\arraybackslash}p{0.55\textwidth}@{}}
\toprule Word & Lemmata & Occ. & Mss & Occurrences in the other EXI manuscripts \\ \midrule \endhead
\bottomrule \endfoot
\textbf{EXI-01-b} & 12 & 0 & 0 & --- \\
\textbf{a} & 2, 5, 6, 7\,(\(\times\)2), 8, 12\,(\(\times\)2) & 1692 & 36 & \emph{function word: totals only} \\
\textbf{able} & 4 & 2 & 1 & 01-c\,(2) \\
\textbf{absurd} & 12 & 46 & 5 & 01-d\,(1); 04-a\,(31); 04-b\,(12); 06-a\,(1); 06-b\,(1) \\
\textbf{action} & 4 & 19 & 9 & 01-c\,(6); 01-d\,(2); 04-a\,(2); 07-a\,(1); 08-a\,(2); 10-a\,(2); 14-a\,(2); 15-b\,(1); 16-b\,(1) \\
\textbf{active} & 2, 7 & 21 & 11 & 01-c\,(2); 01-d\,(1); 04-a\,(1); 07-a\,(1); 10-a\,(1); 10-b\,(2); 11-a\,(2); 11-b\,(1); 17-a\,(7); 19-b\,(1); 20-a\,(2) \\
\textbf{aeternitatis} & 11 & 0 & 0 & --- \\
\textbf{agent} & 4, 5, 6 & 104 & 19 & 01-c\,(15); 01-d\,(4); 04-a\,(15); 04-b\,(3); 06-a\,(2); 06-b\,(1); 07-a\,(15); 07-b\,(5); 08-a\,(13); 09-a\,(1); 10-a\,(5); 10-b\,(3); 11-a\,(1); 15-a\,(4); 15-b\,(7); 17-a\,(2); 18-a\,(6); 18-b\,(1); 19-a\,(1) \\
\textbf{and} & 9, 12 & 1444 & 36 & \emph{function word: totals only} \\
\textbf{answer} & 10, 11 & 34 & 20 & 01-d\,(1); 02-a\,(2); 02-b\,(1); 03-a\,(1); 04-a\,(2); 07-a\,(3); 08-a\,(1); 10-a\,(1); 10-b\,(2); 11-a\,(6); 11-b\,(1); 12-a\,(2); 14-a\,(1); 15-a\,(1); 15-b\,(1); 16-a\,(2); 16-b\,(2); 18-b\,(1); 19-a\,(1); 20-a\,(2) \\
\textbf{answers} & 9 & 32 & 14 & 02-a\,(1); 02-b\,(2); 03-a\,(1); 04-a\,(2); 07-a\,(1); 10-b\,(1); 11-a\,(1); 12-a\,(1); 13-a\,(2); 13-b\,(2); 16-a\,(7); 16-b\,(8); 19-a\,(2); 20-a\,(1) \\
\textbf{any} & 10 & 70 & 27 & \emph{function word: totals only} \\
\textbf{are} & 12 & 346 & 36 & \emph{function word: totals only} \\
\textbf{assumption} & 6, 11, 12 & 67 & 16 & 01-c\,(22); 01-d\,(2); 02-a\,(2); 02-b\,(1); 03-a\,(11); 03-b\,(4); 04-a\,(1); 06-a\,(3); 06-b\,(1); 10-a\,(2); 10-b\,(4); 13-a\,(5); 13-b\,(4); 17-a\,(1); 19-a\,(3); 19-b\,(1) \\
\textbf{at} & 4 & 197 & 31 & \emph{function word: totals only} \\
\textbf{attitudes} & 3 & 31 & 8 & 01-c\,(3); 01-d\,(1); 04-a\,(4); 06-a\,(4); 06-b\,(9); 07-a\,(5); 17-a\,(1); 18-a\,(4) \\
\textbf{available} & 5, 9 & 61 & 18 & 01-c\,(2); 01-d\,(8); 02-a\,(1); 03-a\,(1); 04-a\,(5); 04-b\,(2); 05-a\,(1); 07-a\,(4); 08-a\,(1); 10-a\,(4); 10-b\,(2); 11-a\,(8); 11-b\,(8); 14-a\,(9); 17-a\,(1); 18-a\,(1); 19-a\,(1); 20-a\,(2) \\
\textbf{by} & 6, 8 & 460 & 36 & \emph{function word: totals only} \\
\textbf{can} & 6 & 188 & 33 & \emph{function word: totals only} \\
\textbf{capable} & 4 & 15 & 4 & 01-c\,(6); 01-d\,(2); 04-a\,(2); 09-a\,(5) \\
\textbf{challenge} & 11\,(\(\times\)2) & 10 & 4 & 01-d\,(1); 04-a\,(1); 10-a\,(1); 10-b\,(7) \\
\textbf{commentary} & 12 & 0 & 0 & --- \\
\textbf{consists} & 2, 7 & 18 & 10 & 01-c\,(2); 04-a\,(2); 06-b\,(2); 07-a\,(4); 11-a\,(1); 15-a\,(2); 15-b\,(1); 16-a\,(1); 16-b\,(2); 18-a\,(1) \\
\textbf{defence} & 12 & 1 & 1 & 08-a\,(1) \\
\textbf{definition} & 8 & 153 & 31 & 01-c\,(9); 01-d\,(3); 02-a\,(6); 03-a\,(3); 03-b\,(1); 04-a\,(9); 04-b\,(7); 05-a\,(1); 06-a\,(4); 06-b\,(3); 07-a\,(6); 07-b\,(4); 08-a\,(5); 09-a\,(15); 09-b\,(13); 10-a\,(6); 10-b\,(10); 11-a\,(9); 11-b\,(3); 12-a\,(2); 13-a\,(2); 14-a\,(3); 15-a\,(5); 15-b\,(8); 16-a\,(5); 16-b\,(1); 17-a\,(5); 18-a\,(1); 18-b\,(1); 19-a\,(2); 19-b\,(1) \\
\textbf{degree} & 2, 8 & 31 & 10 & 01-c\,(2); 01-d\,(8); 02-a\,(1); 03-a\,(1); 05-a\,(2); 11-b\,(5); 15-a\,(5); 16-a\,(4); 16-b\,(2); 17-a\,(1) \\
\textbf{denies} & 11 & 14 & 9 & 01-d\,(1); 04-a\,(4); 06-a\,(1); 06-b\,(2); 09-a\,(2); 11-a\,(1); 12-a\,(1); 12-b\,(1); 16-a\,(1) \\
\textbf{depend} & 3 & 36 & 18 & 01-c\,(4); 01-d\,(3); 05-a\,(2); 06-b\,(4); 07-a\,(3); 08-a\,(1); 09-a\,(2); 09-b\,(1); 10-a\,(1); 11-a\,(3); 11-b\,(3); 15-b\,(1); 17-a\,(1); 18-a\,(2); 19-a\,(2); 19-b\,(1); 20-a\,(1); 20-b\,(1) \\
\textbf{depends} & 10 & 53 & 19 & 01-c\,(1); 01-d\,(3); 02-a\,(3); 03-a\,(2); 06-a\,(3); 07-a\,(7); 07-b\,(3); 08-a\,(3); 09-a\,(2); 09-b\,(1); 11-b\,(1); 15-a\,(2); 16-b\,(1); 17-a\,(3); 18-a\,(4); 19-a\,(4); 19-b\,(4); 20-a\,(4); 20-b\,(2) \\
\textbf{discussion} & 12 & 2 & 2 & 02-a\,(1); 04-a\,(1) \\
\textbf{does} & 3, 10, 11 & 306 & 35 & \emph{function word: totals only} \\
\textbf{done} & 10 & 11 & 7 & 01-d\,(3); 03-a\,(2); 03-b\,(1); 04-b\,(1); 13-a\,(1); 14-a\,(1); 17-a\,(2) \\
\textbf{engage} & 4, 6 & 27 & 9 & 01-c\,(2); 03-a\,(5); 06-b\,(2); 10-a\,(7); 10-b\,(4); 11-a\,(1); 11-b\,(1); 19-b\,(4); 20-a\,(1) \\
\textbf{engaged} & 3 & 34 & 14 & 01-c\,(1); 01-d\,(1); 03-a\,(1); 04-a\,(3); 06-b\,(2); 07-a\,(1); 07-b\,(1); 10-a\,(3); 10-b\,(1); 11-a\,(4); 17-a\,(7); 19-b\,(1); 20-a\,(5); 20-b\,(3) \\
\textbf{engagement} & 2, 7\,(\(\times\)2) & 148 & 17 & 01-c\,(14); 01-d\,(10); 03-a\,(3); 04-a\,(3); 04-b\,(8); 06-b\,(4); 07-a\,(14); 07-b\,(4); 10-a\,(5); 10-b\,(11); 11-a\,(13); 11-b\,(9); 14-a\,(1); 17-a\,(24); 19-b\,(11); 20-a\,(9); 20-b\,(5) \\
\textbf{every} & 4, 5 & 71 & 26 & \emph{function word: totals only} \\
\textbf{everyone} & 9 & 11 & 6 & 01-d\,(5); 02-a\,(2); 07-a\,(1); 11-a\,(1); 17-a\,(1); 20-a\,(1) \\
\textbf{exist} & 3 & 29 & 10 & 02-a\,(17); 02-b\,(1); 04-a\,(1); 09-a\,(1); 11-b\,(1); 12-a\,(1); 12-b\,(1); 13-a\,(3); 14-a\,(2); 19-a\,(1) \\
\textbf{extent} & 2, 7 & 29 & 17 & 01-c\,(1); 01-d\,(2); 02-a\,(1); 03-a\,(1); 04-a\,(2); 06-b\,(1); 07-a\,(4); 09-a\,(1); 11-a\,(4); 11-b\,(1); 16-b\,(1); 17-a\,(3); 18-a\,(1); 18-b\,(1); 19-b\,(1); 20-a\,(3); 20-b\,(1) \\
\textbf{for} & 4, 5 & 499 & 36 & \emph{function word: totals only} \\
\textbf{given} & 12 & 25 & 17 & 01-c\,(1); 01-d\,(2); 03-a\,(2); 03-b\,(1); 05-a\,(1); 06-a\,(1); 07-b\,(1); 09-a\,(1); 10-b\,(2); 11-a\,(1); 11-b\,(1); 12-a\,(2); 15-b\,(2); 17-a\,(1); 18-b\,(2); 19-a\,(2); 19-b\,(2) \\
\textbf{if} & 2 & 226 & 34 & \emph{function word: totals only} \\
\textbf{in} & 2, 3, 4, 6, 7, 12 & 752 & 36 & \emph{function word: totals only} \\
\textbf{includes} & 7 & 15 & 9 & 01-c\,(1); 04-a\,(1); 07-a\,(2); 09-b\,(2); 10-a\,(2); 10-b\,(2); 15-b\,(3); 17-a\,(1); 18-a\,(1) \\
\textbf{intentional} & 4 & 18 & 6 & 01-c\,(10); 01-d\,(2); 04-a\,(2); 09-a\,(2); 10-a\,(1); 10-b\,(1) \\
\textbf{is} & 2\,(\(\times\)2), 4\,(\(\times\)2), 5, 6, 8, 9\,(\(\times\)2), 10\,(\(\times\)2) & 1896 & 36 & \emph{function word: totals only} \\
\textbf{it} & 8, 10\,(\(\times\)2), 11\,(\(\times\)2) & 1065 & 36 & \emph{function word: totals only} \\
\textbf{least} & 4 & 14 & 8 & 01-c\,(2); 02-b\,(1); 09-a\,(1); 10-b\,(1); 15-a\,(1); 16-a\,(5); 16-b\,(2); 17-a\,(1) \\
\textbf{life} & 1, 2, 5, 7, 9, 10 & 450 & 33 & 01-c\,(7); 01-d\,(14); 02-a\,(40); 02-b\,(2); 03-a\,(62); 03-b\,(24); 04-a\,(19); 04-b\,(12); 05-a\,(4); 06-a\,(3); 06-b\,(3); 07-a\,(24); 07-b\,(2); 08-a\,(1); 09-a\,(13); 09-b\,(9); 10-a\,(2); 10-b\,(5); 11-a\,(43); 11-b\,(4); 12-a\,(25); 12-b\,(6); 13-a\,(5); 13-b\,(1); 14-a\,(1); 15-b\,(1); 17-a\,(35); 18-a\,(4); 18-b\,(1); 19-a\,(47); 19-b\,(24); 20-a\,(6); 20-b\,(1) \\
\textbf{matters} & 11 & 37 & 15 & 01-c\,(1); 01-d\,(2); 02-a\,(2); 02-b\,(1); 03-a\,(14); 03-b\,(2); 06-a\,(1); 06-b\,(1); 07-a\,(2); 10-a\,(1); 10-b\,(3); 12-a\,(1); 13-a\,(2); 17-a\,(3); 19-b\,(1) \\
\textbf{meaning} & 1, 5, 9\,(\(\times\)2) & 229 & 17 & 01-c\,(9); 01-d\,(28); 03-a\,(2); 04-a\,(5); 04-b\,(4); 06-b\,(16); 07-a\,(3); 10-a\,(8); 10-b\,(13); 11-a\,(49); 11-b\,(22); 14-a\,(1); 17-a\,(31); 18-a\,(1); 19-a\,(13); 19-b\,(21); 20-a\,(3) \\
\textbf{meaningful} & 2, 8, 10 & 92 & 16 & 01-c\,(3); 01-d\,(9); 03-a\,(1); 04-a\,(3); 04-b\,(2); 06-b\,(2); 09-a\,(1); 09-b\,(1); 10-b\,(2); 11-a\,(27); 11-b\,(3); 14-a\,(1); 17-a\,(23); 19-b\,(8); 20-a\,(5); 20-b\,(1) \\
\textbf{nor} & 11 & 23 & 14 & \emph{function word: totals only} \\
\textbf{not} & 3, 10 & 809 & 36 & \emph{function word: totals only} \\
\textbf{nothing} & 11 & 79 & 25 & 01-c\,(2); 01-d\,(3); 02-a\,(7); 02-b\,(4); 03-b\,(1); 04-a\,(5); 04-b\,(6); 05-a\,(2); 06-b\,(5); 07-a\,(4); 08-a\,(1); 09-a\,(2); 10-a\,(2); 10-b\,(1); 11-a\,(5); 12-a\,(1); 13-a\,(4); 13-b\,(2); 14-a\,(3); 15-a\,(1); 15-b\,(6); 16-a\,(1); 17-a\,(7); 19-a\,(1); 19-b\,(3) \\
\textbf{objective} & 2, 3, 4, 6, 7 & 25 & 9 & 01-c\,(7); 01-d\,(3); 03-a\,(1); 04-a\,(6); 06-b\,(2); 09-a\,(1); 11-a\,(2); 17-a\,(2); 20-a\,(1) \\
\textbf{of} & 1, 2, 3\,(\(\times\)2), 4\,(\(\times\)2), 5, 6, 7, 9\,(\(\times\)3), 10, 11, 12\,(\(\times\)2) & 1794 & 36 & \emph{function word: totals only} \\
\textbf{on} & 1, 3, 10 & 418 & 36 & \emph{function word: totals only} \\
\textbf{one} & 4 & 279 & 34 & \emph{function word: totals only} \\
\textbf{particular} & 10 & 27 & 15 & 01-d\,(2); 02-a\,(1); 03-b\,(1); 04-a\,(1); 05-a\,(3); 06-b\,(1); 07-a\,(3); 09-a\,(1); 11-a\,(1); 11-b\,(2); 13-a\,(2); 16-a\,(4); 19-a\,(3); 19-b\,(1); 20-b\,(1) \\
\textbf{positive} & 5, 7, 8 & 23 & 9 & 01-c\,(4); 01-d\,(3); 02-a\,(4); 07-a\,(1); 10-a\,(1); 14-a\,(2); 15-b\,(1); 18-a\,(3); 18-b\,(4) \\
\textbf{premise} & 11 & 18 & 10 & 01-c\,(3); 01-d\,(1); 02-a\,(5); 02-b\,(2); 03-a\,(1); 04-a\,(1); 05-a\,(2); 06-a\,(1); 10-a\,(1); 11-a\,(1) \\
\textbf{project} & 4, 6, 7 & 77 & 10 & 01-c\,(13); 01-d\,(2); 03-a\,(12); 03-b\,(8); 04-a\,(1); 06-b\,(1); 09-a\,(21); 09-b\,(17); 17-a\,(1); 18-a\,(1) \\
\textbf{projects} & 2, 3\,(\(\times\)2) & 61 & 15 & 01-c\,(10); 01-d\,(9); 03-a\,(8); 03-b\,(4); 04-a\,(5); 04-b\,(3); 06-b\,(1); 09-a\,(3); 11-a\,(2); 11-b\,(4); 13-a\,(4); 18-a\,(1); 19-a\,(3); 19-b\,(3); 20-a\,(1) \\
\textbf{question} & 9, 10 & 80 & 25 & 02-a\,(2); 03-a\,(2); 04-a\,(14); 04-b\,(2); 05-a\,(1); 06-a\,(3); 06-b\,(3); 07-a\,(1); 07-b\,(1); 08-a\,(10); 10-a\,(5); 10-b\,(2); 11-a\,(10); 12-a\,(2); 13-a\,(1); 13-b\,(2); 14-a\,(2); 15-a\,(1); 16-a\,(1); 17-a\,(2); 18-a\,(1); 18-b\,(1); 19-a\,(1); 20-a\,(7); 20-b\,(3) \\
\textbf{rather} & 11 & 39 & 20 & 01-c\,(2); 01-d\,(2); 04-a\,(1); 04-b\,(1); 06-b\,(1); 07-b\,(1); 08-a\,(6); 09-a\,(1); 10-a\,(2); 10-b\,(5); 11-a\,(1); 13-a\,(2); 14-a\,(1); 15-b\,(2); 16-a\,(2); 16-b\,(2); 17-a\,(2); 18-a\,(1); 19-a\,(2); 20-a\,(2) \\
\textbf{refuting} & 11 & 2 & 2 & 01-d\,(1); 04-a\,(1) \\
\textbf{sceptical} & 11 & 2 & 2 & 04-a\,(1); 08-a\,(1) \\
\textbf{she} & 4 & 193 & 27 & \emph{function word: totals only} \\
\textbf{some} & 7 & 59 & 20 & \emph{function word: totals only} \\
\textbf{specie} & 11 & 0 & 0 & --- \\
\textbf{sub} & 11 & 0 & 0 & --- \\
\textbf{such} & 5 & 92 & 26 & \emph{function word: totals only} \\
\textbf{than} & 11 & 139 & 30 & \emph{function word: totals only} \\
\textbf{that} & 7, 11\,(\(\times\)2) & 1327 & 36 & \emph{function word: totals only} \\
\textbf{the} & 1, 2, 3, 6, 9\,(\(\times\)3), 10, 11\,(\(\times\)2), 12\,(\(\times\)2) & 3437 & 36 & \emph{function word: totals only} \\
\textbf{them} & 3 & 122 & 31 & \emph{function word: totals only} \\
\textbf{theorem} & 9 & 81 & 33 & 01-c\,(2); 01-d\,(4); 02-a\,(2); 02-b\,(1); 03-a\,(2); 03-b\,(4); 04-a\,(4); 04-b\,(3); 05-a\,(2); 06-a\,(5); 06-b\,(3); 07-a\,(6); 07-b\,(3); 08-a\,(1); 09-a\,(1); 10-a\,(1); 10-b\,(4); 11-a\,(5); 11-b\,(2); 12-a\,(4); 12-b\,(2); 13-a\,(2); 13-b\,(1); 14-a\,(2); 15-a\,(1); 15-b\,(1); 16-a\,(3); 16-b\,(1); 17-a\,(2); 18-a\,(1); 18-b\,(1); 19-a\,(4); 19-b\,(1) \\
\textbf{there} & 3, 4, 6 & 72 & 24 & \emph{function word: totals only} \\
\textbf{those} & 3 & 55 & 24 & \emph{function word: totals only} \\
\textbf{to} & 2\,(\(\times\)2), 4, 7, 8, 9 & 1050 & 36 & \emph{function word: totals only} \\
\textbf{value} & 3 & 93 & 26 & 01-c\,(6); 01-d\,(2); 02-a\,(11); 03-a\,(2); 03-b\,(1); 04-a\,(3); 04-b\,(2); 05-a\,(2); 06-a\,(2); 06-b\,(2); 07-a\,(8); 09-b\,(2); 10-a\,(1); 10-b\,(3); 11-a\,(3); 12-b\,(1); 13-a\,(2); 14-a\,(1); 15-a\,(5); 15-b\,(4); 16-a\,(7); 16-b\,(3); 17-a\,(7); 18-a\,(1); 19-a\,(11); 19-b\,(1) \\
\textbf{what} & 9, 10 & 311 & 36 & \emph{function word: totals only} \\
\textbf{whether} & 9, 10 & 102 & 26 & \emph{function word: totals only} \\
\textbf{which} & 2, 4, 6, 10 & 345 & 35 & \emph{function word: totals only} \\
\textbf{whose} & 3 & 63 & 26 & \emph{function word: totals only} \\
\textbf{with} & 2, 7\,(\(\times\)2), 10 & 348 & 36 & \emph{function word: totals only} \\
\textbf{worth} & 2, 3, 4, 6, 7 & 82 & 17 & 01-c\,(7); 01-d\,(7); 02-a\,(2); 03-a\,(2); 03-b\,(2); 04-a\,(7); 04-b\,(3); 06-b\,(9); 09-a\,(2); 10-a\,(2); 11-a\,(16); 11-b\,(8); 15-a\,(1); 17-a\,(3); 19-b\,(7); 20-a\,(2); 20-b\,(2) \\
\end{longtable}
}

\subsection*{A.2 Term families by manuscript}

{\footnotesize Occurrences of the fifteen families defined in the preface, by manuscript; \enquote{Words} is the number of words of the manuscript's text after the removal of markup (for EXI-01-a, of its eleven sentences). A dot stands for zero. The last row totals the other manuscripts.\par}

{\footnotesize% Generated by gen/generate.py; do not edit.
\begin{longtable}{@{}l r r r r r r r r r@{}}
\toprule Ms & Words & \textit{mean-} & \textit{life-} & \textit{engag-} & \textit{activ-} & \textit{project-} & \textit{worth-} & \textit{objectiv-} & \textit{valu-} \\ \midrule \endhead
\bottomrule \endfoot
EXI-01-a & 200 & 6 & 5 & 6 & 2 & 6 & 5 & 5 & 1 \\
EXI-01-c & 1130 & 12 & 11 & 18 & 10 & 23 & 16 & 7 & 6 \\
EXI-01-d & 1031 & 38 & 18 & 15 & 1 & 11 & 13 & 3 & 3 \\
EXI-02-a & 2968 & . & 56 & . & . & . & 2 & . & 13 \\
EXI-02-b & 651 & . & 4 & . & . & . & . & . & . \\
EXI-03-a & 2359 & 4 & 73 & 11 & 2 & 20 & 9 & 1 & 4 \\
EXI-03-b & 772 & . & 24 & . & . & 12 & 4 & . & 1 \\
EXI-04-a & 2421 & 10 & 23 & 7 & 1 & 6 & 15 & 6 & 3 \\
EXI-04-b & 1182 & 6 & 15 & 8 & 2 & 3 & 8 & . & 2 \\
EXI-05-a & 1189 & 1 & 4 & . & . & . & . & . & 2 \\
EXI-06-a & 1197 & . & 3 & . & . & . & . & . & 2 \\
EXI-06-b & 845 & 18 & 10 & 9 & . & 2 & 10 & 2 & 3 \\
EXI-07-a & 2302 & 8 & 27 & 15 & 12 & . & 1 & . & 20 \\
EXI-07-b & 747 & 2 & 2 & 5 & 1 & . & . & . & 5 \\
EXI-08-a & 1505 & 5 & 2 & . & . & . & . & . & . \\
EXI-09-a & 1869 & 3 & 19 & . & 3 & 24 & 2 & 1 & 2 \\
EXI-09-b & 1013 & 1 & 11 & 1 & 1 & 17 & 1 & . & 2 \\
EXI-10-a & 1734 & 9 & 2 & 16 & 29 & . & 3 & . & 1 \\
EXI-10-b & 4261 & 15 & 5 & 17 & 19 & . & 1 & . & 3 \\
EXI-11-a & 2382 & 77 & 45 & 23 & 11 & 2 & 33 & 2 & 5 \\
EXI-11-b & 726 & 25 & 9 & 12 & 1 & 4 & 12 & . & . \\
EXI-12-a & 1312 & . & 49 & . & . & . & . & . & . \\
EXI-12-b & 758 & . & 8 & . & . & . & . & . & 1 \\
EXI-13-a & 2107 & . & 22 & . & . & 4 & . & . & 2 \\
EXI-13-b & 746 & . & 2 & . & . & . & . & . & . \\
EXI-14-a & 1210 & 2 & 1 & 1 & . & . & . & . & 1 \\
EXI-15-a & 1514 & . & . & . & . & . & 1 & . & 6 \\
EXI-15-b & 1158 & . & 1 & . & . & . & . & . & 4 \\
EXI-16-a & 2503 & . & . & . & . & . & . & . & 12 \\
EXI-16-b & 827 & . & . & . & . & . & . & . & 3 \\
EXI-17-a & 2018 & 57 & 46 & 32 & 28 & 1 & 20 & 2 & 8 \\
EXI-18-a & 1181 & 1 & 6 & . & . & 2 & 1 & . & 2 \\
EXI-18-b & 300 & . & 1 & . & . & . & . & . & . \\
EXI-19-a & 1470 & 15 & 48 & . & 1 & 3 & . & . & 15 \\
EXI-19-b & 861 & 29 & 24 & 19 & 2 & 3 & 10 & . & 4 \\
EXI-20-a & 735 & 9 & 6 & 15 & 11 & 1 & 11 & 1 & . \\
EXI-20-b & 412 & 1 & 1 & 8 & . & . & 4 & . & . \\
\midrule
Others & 51396 & 348 & 578 & 232 & 135 & 138 & 177 & 25 & 135 \\
\end{longtable}
}

{\footnotesize% Generated by gen/generate.py; do not edit.
\begin{longtable}{@{}l r r r r r r r r@{}}
\toprule Ms & Words & \textit{attitud-} & \textit{agent-} & \textit{availab-} & \textit{degree-} & \textit{matter-} & \textit{absurd-} & \textit{sceptic-} \\ \midrule \endhead
\bottomrule \endfoot
EXI-01-a & 200 & 1 & 3 & 2 & 4 & 1 & 1 & 1 \\
EXI-01-c & 1130 & 4 & 18 & 7 & 3 & 1 & . & . \\
EXI-01-d & 1031 & 1 & 5 & 22 & 11 & 3 & 1 & . \\
EXI-02-a & 2968 & 1 & . & 1 & 2 & 4 & . & . \\
EXI-02-b & 651 & 1 & . & . & . & 2 & . & . \\
EXI-03-a & 2359 & . & . & 1 & 2 & 14 & . & . \\
EXI-03-b & 772 & . & . & . & . & 2 & . & . \\
EXI-04-a & 2421 & 14 & 15 & 6 & 2 & 3 & 37 & 1 \\
EXI-04-b & 1182 & 3 & 3 & 2 & . & . & 15 & . \\
EXI-05-a & 1189 & 5 & . & 1 & 2 & . & . & . \\
EXI-06-a & 1197 & 8 & 3 & . & . & 2 & 1 & . \\
EXI-06-b & 845 & 9 & 1 & . & 1 & 2 & 1 & . \\
EXI-07-a & 2302 & 6 & 16 & 4 & 4 & 3 & . & . \\
EXI-07-b & 747 & . & 5 & . & . & 1 & . & . \\
EXI-08-a & 1505 & . & 15 & 1 & . & 1 & . & 3 \\
EXI-09-a & 1869 & . & 2 & . & 1 & . & . & . \\
EXI-09-b & 1013 & . & . & . & . & . & . & . \\
EXI-10-a & 1734 & . & 6 & 4 & . & 2 & . & . \\
EXI-10-b & 4261 & . & 6 & 2 & . & 6 & . & . \\
EXI-11-a & 2382 & . & 2 & 8 & 6 & 4 & . & . \\
EXI-11-b & 726 & . & . & 9 & 7 & . & . & . \\
EXI-12-a & 1312 & . & . & . & . & 1 & . & . \\
EXI-12-b & 758 & 1 & . & . & . & . & . & . \\
EXI-13-a & 2107 & . & . & . & . & 3 & . & . \\
EXI-13-b & 746 & . & . & . & . & . & . & . \\
EXI-14-a & 1210 & 1 & . & 9 & . & . & . & . \\
EXI-15-a & 1514 & . & 12 & . & 5 & . & . & . \\
EXI-15-b & 1158 & 1 & 7 & . & . & . & . & . \\
EXI-16-a & 2503 & 3 & 1 & . & 4 & 1 & . & . \\
EXI-16-b & 827 & . & . & . & 3 & . & . & . \\
EXI-17-a & 2018 & 1 & 2 & 2 & 4 & 5 & . & . \\
EXI-18-a & 1181 & 4 & 8 & 1 & 1 & 2 & . & . \\
EXI-18-b & 300 & . & 2 & . & 1 & . & . & . \\
EXI-19-a & 1470 & . & 1 & 1 & . & . & . & . \\
EXI-19-b & 861 & . & . & . & 1 & 4 & . & . \\
EXI-20-a & 735 & . & . & 2 & 3 & . & . & . \\
EXI-20-b & 412 & . & . & . & 1 & . & . & . \\
\midrule
Others & 51396 & 63 & 130 & 83 & 64 & 66 & 55 & 4 \\
\end{longtable}
}

\subsection*{A.3 Keyword-in-context concordance}

{\footnotesize Every occurrence, in the other manuscripts of the discipline, of every content word of EXI-01-a, in alphabetical order of the word and, within a word, in the order of the manuscripts and of their text. Each line gives the manuscript (without the prefix \enquote{EXI-}) and the section, then the context. The number after each heading is the number of lines.\par}

{\scriptsize% Generated by gen/generate.py; do not edit.
\kwichead{able}{2}
\begin{longtable}{@{}>{\raggedright}p{0.07\textwidth}>{\raggedleft}p{0.41\textwidth}@{\,}>{\raggedright\arraybackslash}p{0.47\textwidth}@{}}
01-c Abs & ...the assumption that every such agent is  & \textbf{able} to engage in at least one project of ob... \\
01-c 1 & ... agent capable of intentional action is  & \textbf{able} to engage in at least one of them (Assu... \\
\end{longtable}
\kwichead{absurd}{46}
\begin{longtable}{@{}>{\raggedright}p{0.07\textwidth}>{\raggedleft}p{0.41\textwidth}@{\,}>{\raggedright\arraybackslash}p{0.47\textwidth}@{}}
01-d 4 & EXI-04 grants that life is  & \textbf{absurd} in Nagel's sense and argues that the fi... \\
04-a Abs & Is life  & \textbf{absurd}, and if so, how should one respond? \\
04-a Abs & Nagel (1971) located the  & \textbf{absurd} not in the smallness or brevity of huma... \\
04-a Abs & We argue that life is  & \textbf{absurd} in this sense, because both of the coll... \\
04-a Abs & ...e key move is to treat responses to the  & \textbf{absurd} as attitudes assessable for fit. \\
04-a Abs & ...nd that discrepancy is exactly what the  & \textbf{absurd} consists in. \\
04-a 1 & ...veals gives no reason for anything, the  & \textbf{absurd} does not warrant despair; and the heroi... \\
04-a 1 & ...recommendation was that we approach our  & \textbf{absurd} lives with irony. \\
04-a 1 & One asks whether the  & \textbf{absurd} is a defect at all: if it follows from ... \\
04-a 1 & ...e value, if there is any, dissolves the  & \textbf{absurd}: on theories of meaning such as those o... \\
04-a 1 & We argue that life is  & \textbf{absurd} in Nagel's sense even if those theories... \\
04-a 2 & A life is  & \textbf{absurd} if its subject takes some ends seriousl... \\
04-a 2 & ...tion of the appropriate response to the  & \textbf{absurd} is then the question of which attitude ... \\
04-a 2 & ...question of which attitude presents the  & \textbf{absurd} situation as it is. \\
04-a 3 & ...luative convictions is another, and the  & \textbf{absurd} concerns the second. \\
04-a 3 & ... agent capable of intentional action is  & \textbf{absurd} in the sense of Definition 2.3. \\
04-a 3 & By Definition 2.3, her life is  & \textbf{absurd}. \\
04-a 3 & The  & \textbf{absurd} of Definition 2.3 is not a relation bet... \\
04-a 3 & ...ase and which Proposition 3.4 shows the  & \textbf{absurd} gives no reason to abandon. \\
04-a 4 & Irony is a way of not taking the  & \textbf{absurd} seriously. \\
04-a 4 & The demand that the  & \textbf{absurd} be met with constant tension presuppose... \\
04-a 4 & To take the  & \textbf{absurd} exactly as seriously as it deserves is ... \\
04-a 4 & ... life of someone engaged in them is not  & \textbf{absurd}. \\
04-a 4 & ...ed, if the projects are worthwhile; the  & \textbf{absurd} survives because that justification is ... \\
04-a 4 & ...ther from fitting it, but it is no less  & \textbf{absurd}, and irony fits it as well. \\
04-a 4 & The appropriate response to the  & \textbf{absurd} is not irony but forgetting. \\
04-a 4 & Forgetting is not a response to the  & \textbf{absurd} but an interruption of the reflection t... \\
04-a 4 & ...the question was what attitude fits the  & \textbf{absurd} when one does attend to it, and to that... \\
04-a 5 & ...n to despair, but the reason is not the  & \textbf{absurd}. \\
04-a 5 & ... reflects in the way that discloses the  & \textbf{absurd} has no occasion for irony. \\
04-a 6 & Life is  & \textbf{absurd} in Nagel's sense: a reflective agent ca... \\
04-a 6 & The  & \textbf{absurd} so understood is the absence of a vindi... \\
04-b Abs & ...s the condemned king as the hero of the  & \textbf{absurd}. \\
04-b Abs & EXI-04-a argues that life is  & \textbf{absurd} in Nagel's sense and that the fitting r... \\
04-b Abs & ...hat his situation is not a model of the  & \textbf{absurd} but of futility, which is a different c... \\
04-b 1 & He is the  & \textbf{absurd} hero because he is lucid about his fate... \\
04-b 1 & EXI-04-a defines the  & \textbf{absurd} as the conjunction of two features: the... \\
04-b 1 & A life can be  & \textbf{absurd} without being futile, and a task can be... \\
04-b 1 & ...ask can be futile in a life that is not  & \textbf{absurd}. \\
04-b 1 & ..., the life of every reflective agent is  & \textbf{absurd}, including the life of an agent whose p... \\
04-b 1 & ...ndpoint, and her life would then not be  & \textbf{absurd} in the sense of Definition 2.3 of EXI-0... \\
04-b 2 & ...revolt is not a fitting response to the  & \textbf{absurd}. \\
04-b 2 & ...-04-a it is the fitting response to the  & \textbf{absurd}. \\
04-b 5 & ...emblem of the human condition, which is  & \textbf{absurd} without being futile, faces no adversar... \\
06-a 5 & This bears on the response to the  & \textbf{absurd} (EXI-04): a world that is not an advers... \\
06-b 2 & ...revolt is not a fitting response to the  & \textbf{absurd}, since the world it scorns is indiffere... \\
\end{longtable}
\kwichead{action}{19}
\begin{longtable}{@{}>{\raggedright}p{0.07\textwidth}>{\raggedleft}p{0.41\textwidth}@{\,}>{\raggedright\arraybackslash}p{0.47\textwidth}@{}}
01-c Abs & ...e to every agent capable of intentional  & \textbf{action}. \\
01-c 1 & ...that every agent capable of intentional  & \textbf{action} is able to engage in at least one of th... \\
01-c 2 & Every agent capable of intentional  & \textbf{action} can direct some intentional activity, w... \\
01-c 2 & For every agent capable of intentional  & \textbf{action}, a life of positive meaning is availabl... \\
01-c 4 & Some agents capable of intentional  & \textbf{action} are wholly isolated and wholly unable t... \\
01-c 4 & An agent capable of intentional  & \textbf{action} acts on something, and anything she act... \\
01-d 1 & ... for every agent capable of intentional  & \textbf{action} a life of positive meaning is available... \\
01-d 2 & No agent capable of intentional  & \textbf{action} is condemned by her circumstances alone... \\
04-a 3 & Any agent capable of intentional  & \textbf{action} takes some ends seriously. \\
04-a 3 & ...reflective agent capable of intentional  & \textbf{action} is absurd in the sense of Definition 2.... \\
07-a 4 & ... claim that the good life is the end of  & \textbf{action} is compatible with our thesis, and inde... \\
08-a Abs & ... one's own is kept open and answered in  & \textbf{action}. \\
08-a 5 & ... ownership is kept open and answered in  & \textbf{action}, and this is what authenticity, read as... \\
10-a 3 & An aversive state motivates  & \textbf{action} that ends it, and the action that ends ... \\
10-a 3 & ... motivates action that ends it, and the  & \textbf{action} that ends boredom is a change of activi... \\
14-a 4 & An emotion that motivates no possible  & \textbf{action} is irrational. \\
14-a 4 & ...does not require that an emotion prompt  & \textbf{action}, as grief and admiration show. \\
15-b 3 & ... familiar in medicine and in collective  & \textbf{action}. \\
16-b 2 & ... 3.3 of EXI-16-a applied to reasons for  & \textbf{action}: a reason that consists in features any... \\
\end{longtable}
\kwichead{active}{21}
\begin{longtable}{@{}>{\raggedright}p{0.07\textwidth}>{\raggedleft}p{0.41\textwidth}@{\,}>{\raggedright\arraybackslash}p{0.47\textwidth}@{}}
01-c 1 & ...ngful to the degree that it consists in  & \textbf{active} engagement with projects of objective w... \\
01-c 2 & ...gement consists to a positive extent in  & \textbf{active} engagement with a project of objective ... \\
01-d 1 & ... by a different route: meaning requires  & \textbf{active} engagement with something worthwhile an... \\
04-a 4 & Meaning requires  & \textbf{active} engagement with something worthwhile (E... \\
07-a 3 & ...and EXI-17, on which meaning comes from  & \textbf{active} engagement with what is worthwhile. \\
10-a 4 & ...ng in EXI-17, on which meaning requires  & \textbf{active} engagement with something worthwhile, a... \\
10-b 5 & ... default mode network, the network more  & \textbf{active} when attention is not directed to an ex... \\
10-b 8 & ...7 holds that a meaningful life requires  & \textbf{active} engagement with something worthwhile; t... \\
11-a 1 & ...brid: meaning arises when the subject's  & \textbf{active} engagement meets the objective worth of... \\
11-a 7 & Meaning consists in  & \textbf{active} engagement with what is worthwhile, and... \\
11-b 1 & EXI-11-a defines meaning as  & \textbf{active} engagement with what is worthwhile (Def... \\
17-a Abs & A meaningful life requires  & \textbf{active} engagement with something worthwhile, i... \\
17-a 1 & ...ue that it should not: meaning requires  & \textbf{active} engagement with something worthwhile, a... \\
17-a 2 &  & \textbf{Active} engagement in the sense relevant to Pri... \\
17-a 3 & Her engagement is  & \textbf{active}, since study and attentive listening ar... \\
17-a 3 & A meaningful life requires  & \textbf{active} engagement with something worthwhile; i... \\
17-a 4 & ...ce of a difficult activity counts, then  & \textbf{active} engagement in a demanding worthwhile ac... \\
17-a 5 & A meaningful life requires  & \textbf{active} engagement with something worthwhile. \\
19-b 1 & ...e without categorical desires lacks the  & \textbf{active} engagement that meaning requires (EXI-1... \\
20-a Abs & ...eaningful activity within one, requires  & \textbf{active} engagement with something worthwhile. \\
20-a 1 & ... version of this view: meaning requires  & \textbf{active} engagement with something worthwhile, a... \\
\end{longtable}
\kwichead{aeternitatis}{0}
\kwicnone
\kwichead{agent}{104}
\begin{longtable}{@{}>{\raggedright}p{0.07\textwidth}>{\raggedleft}p{0.41\textwidth}@{\,}>{\raggedright\arraybackslash}p{0.47\textwidth}@{}}
01-c Abs & ... positive meaning is available to every  & \textbf{agent} capable of intentional action. \\
01-c Abs & ...rests on the assumption that every such  & \textbf{agent} is able to engage in at least one proje... \\
01-c Abs & ...h objects are within the reach of every  & \textbf{agent}, and show that sustained engagement wit... \\
01-c 1 & ... (Assumption 2), and assumes that every  & \textbf{agent} capable of intentional action is able t... \\
01-c 2 & ...t of engagement is anything at which an  & \textbf{agent} can direct intentional activity and abo... \\
01-c 2 & Every  & \textbf{agent} capable of intentional action can direc... \\
01-c 2 & If an  & \textbf{agent} directs intentional activity with care ... \\
01-c 2 & ...f a project is that it be something the  & \textbf{agent} engages in, actively, and that its valu... \\
01-c 2 & ...direction, extended over time, that the  & \textbf{agent} can pursue well or badly. \\
01-c 2 & For every  & \textbf{agent} capable of intentional action, a life o... \\
01-c 2 & By Assumption 2.2, the  & \textbf{agent} can direct intentional activity with ca... \\
01-c 3 & Almost every  & \textbf{agent} lives among other persons, whose good h... \\
01-c 3 & ...independent of her attitudes, and every  & \textbf{agent} lives in a world some part of which she... \\
01-c 4 & An  & \textbf{agent} capable of intentional action acts on s... \\
01-c 4 & The case of an  & \textbf{agent} for whom nothing worthwhile is within r... \\
01-d 1 & ...rem 4 of EXI-01-a states that for every  & \textbf{agent} capable of intentional action a life of... \\
01-d 2 & No  & \textbf{agent} capable of intentional action is condem... \\
01-d 3 & ...ility to actuality: a life in which the  & \textbf{agent} engages with nothing worthwhile, though... \\
01-d 5 & ...s reach: engagement, which is up to the  & \textbf{agent}, and worth, which belongs to what she e... \\
04-a Abs & ...res are ineliminable for any reflective  & \textbf{agent}, and that the appropriate response is i... \\
04-a 2 & An  & \textbf{agent} takes an end seriously if she devotes e... \\
04-a 3 & Any  & \textbf{agent} capable of intentional action takes som... \\
04-a 3 & An  & \textbf{agent} who took no end seriously in the sense ... \\
04-a 3 & For any reflective  & \textbf{agent} and any justification she can offer for... \\
04-a 3 & ...r an end cites a consideration that the  & \textbf{agent} treats as mattering: the end serves ano... \\
04-a 3 & From the detached standpoint the  & \textbf{agent} can ask, of that consideration, why it ... \\
04-a 3 & ...compelled to endorse it by anything the  & \textbf{agent} can exhibit, since what she would exhib... \\
04-a 3 & ...forbids the step back, and a reflective  & \textbf{agent} is one who can take it. \\
04-a 3 & The life of any reflective  & \textbf{agent} capable of intentional action is absurd... \\
04-a 3 & By Lemma 3.1 the  & \textbf{agent} takes some ends seriously. \\
04-a 3 & ...ition 2.3 is not a relation between the  & \textbf{agent} and the world but between two standpoin... \\
04-a 3 & ...orld but between two standpoints of the  & \textbf{agent}. \\
04-a 3 & ...nds, which Lemma 3.1 shows a reflective  & \textbf{agent} will do in any case and which Propositi... \\
04-a 6 & ...s absurd in Nagel's sense: a reflective  & \textbf{agent} cannot avoid taking her ends seriously,... \\
04-b 1 & ... EXI-04-a, the life of every reflective  & \textbf{agent} is absurd, including the life of an age... \\
04-b 1 & ...ent is absurd, including the life of an  & \textbf{agent} whose projects succeed and whose object... \\
04-b 1 & Conversely, an unreflective  & \textbf{agent} could be set to a futile task without e... \\
06-a 3 & ... a whole and its laws, which are not an  & \textbf{agent} and are not composed so as to aggregate... \\
06-a 4 & If the universe is not an  & \textbf{agent}, it is not indifferent either; indiffer... \\
06-b 3 & ...sed without reference to any particular  & \textbf{agent}, then the worth of objects of engagemen... \\
07-a Abs & An  & \textbf{agent} whose only final end is her own happine... \\
07-a 2 & An  & \textbf{agent} has as a final end if she desires for i... \\
07-a 2 & An  & \textbf{agent} aims directly at happiness if her own h... \\
07-a 2 & ...he pursuits contribute to only when the  & \textbf{agent} values their objects for their own sake... \\
07-a 2 & ...gement in Principle P1 require that the  & \textbf{agent} care about their objects for their own ... \\
07-a 2 & ...s that its being solved matters, and an  & \textbf{agent} to whom it matters only as a source of ... \\
07-a 3 & ... contributes to happiness only when the  & \textbf{agent} values its objects for their own sake. \\
07-a 3 & An  & \textbf{agent} who aims directly at happiness does not... \\
07-a 3 & ...ctly at happiness is self-defeating: an  & \textbf{agent} who does so forgoes the large part of h... \\
07-a 3 & By Lemma 3.2, the  & \textbf{agent} who aims directly at happiness does not... \\
07-a 3 & An  & \textbf{agent} who values other things for their own s... \\
07-a 3 & ...eaker thesis is true but contingent: an  & \textbf{agent} could perhaps aim at happiness while ke... \\
07-a 3 & An  & \textbf{agent} who never once thinks about her happine... \\
07-a 4 & Railton's  & \textbf{agent} can stably value other things for their... \\
07-a 6 & ...eople valued for their own sake, and an  & \textbf{agent} whose only final end is her own happine... \\
07-b 1 & ...e lies not in attention but in what the  & \textbf{agent} values: an agent whose sole final end i... \\
07-b 1 & ...ention but in what the agent values: an  & \textbf{agent} whose sole final end is her own happine... \\
07-b 1 & ...s indirectly self-defeating if, when an  & \textbf{agent} successfully follows it, its aims are w... \\
07-b 1 & The  & \textbf{agent} who aims directly at happiness does not... \\
07-b 1 & By Theorem 3.3 of EXI-07-a, an  & \textbf{agent} who does so attains less happiness than... \\
08-a 2 & An  & \textbf{agent} is authentic in the achievement sense a... \\
08-a 2 & An  & \textbf{agent} is authentic in the orientation sense o... \\
08-a 2 & ...uestion, and it is satisfied by how the  & \textbf{agent} stands to the question over time. \\
08-a 2 & Every interpretation an  & \textbf{agent} can give of her commitments, including ... \\
08-a 2 & No  & \textbf{agent} coincides with her commitments: whateve... \\
08-a 3 & Suppose an  & \textbf{agent} were authentic in the achievement sense... \\
08-a 3 & ...eted, cannot be held is not one that an  & \textbf{agent} can be in. \\
08-a 3 & ...f Definition 2.2 requires only that the  & \textbf{agent} treat a question as open, which Princip... \\
08-a 3 & ...is not trivial, because it can fail: an  & \textbf{agent} who meets every such occasion by retrea... \\
08-a 4 & ...ponding that must be maintained, and an  & \textbf{agent} who has them can lose them. \\
08-a 4 & An  & \textbf{agent} can believe she is keeping the question... \\
08-a 4 & ...matter of conduct, which others and the  & \textbf{agent} herself can observe over time. \\
08-a 5 & ...on draws on shared meanings and that no  & \textbf{agent} coincides with her commitments, show th... \\
09-a Abs & ...ect for a life can be fixed only by the  & \textbf{agent} who lives it. \\
10-a 3 & An activity fails to engage an  & \textbf{agent} who is not otherwise impaired just when... \\
10-a 3 & ...has the adaptive function of moving the  & \textbf{agent} toward activities that match her capaci... \\
10-a 3 & ...redom produces the effect of moving the  & \textbf{agent} toward matching activities. \\
10-a 3 & By Definition 2.3, moving the  & \textbf{agent} toward matching activities is its funct... \\
10-a 3 & ...ignal that cannot be acted on, when the  & \textbf{agent} is confined to mismatched activity. \\
10-b Abs & ...has the adaptive function of moving the  & \textbf{agent} toward activities that match them. \\
10-b 1 & ... selected-effect function of moving the  & \textbf{agent} toward matching activities (its Lemma 3... \\
10-b 6 & ...nction thesis is that boredom moves the  & \textbf{agent} toward activities that match her capaci... \\
11-a 2 & ...hat an activity has been assigned to an  & \textbf{agent} as a purpose by another, however powerf... \\
15-a 2 & An  & \textbf{agent} hopes that if (i) she desires that ; (i... \\
15-a 2 & ... the agent's evidence (EPI-17), and the  & \textbf{agent} does not control her credence directly ... \\
15-a 3 & ...tance is insensitive to probability, an  & \textbf{agent} invests as much attention, planning and... \\
15-a 4 & It exposes the  & \textbf{agent} to disappointment, but so does every at... \\
15-b Abs & ...what is hoped for, by changing what the  & \textbf{agent} does, from the illegitimate way in whic... \\
15-b 1 & On the account of EXI-15-a, an  & \textbf{agent} hopes that if she desires , has a crede... \\
15-b 1 & ...r the credence, for the value of to the  & \textbf{agent} and for her capacity to make more likel... \\
15-b 2 & ...ful stance passes into expectation: the  & \textbf{agent} plans on rather than for its possibilit... \\
15-b 2 & When is close to one, the rational  & \textbf{agent} plans on as she plans on what she expec... \\
15-b 3 & ...s probability only by changing what the  & \textbf{agent} does, hope is to be preferred, since it... \\
15-b 4 & ... and it does not exclude the despairing  & \textbf{agent} who makes the right choices but ceases ... \\
17-a 2 & An  & \textbf{agent} is actively engaged with an object if s... \\
17-a 2 & An  & \textbf{agent} achieves an outcome if she competently ... \\
18-a Abs & ...he sense that no one can function as an  & \textbf{agent} without some of them and their wholesal... \\
18-a 2 & ...ns, is psychologically necessary for an  & \textbf{agent} if she could not function as an agent, ... \\
18-a 2 & ...n agent if she could not function as an  & \textbf{agent}, in the sense of forming and pursuing p... \\
18-a 2 & ...r direct voluntary control (EPI-11), an  & \textbf{agent} cannot remove an illusion by deciding t... \\
18-a 3 & ...elief is necessary for functioning, the  & \textbf{agent} has a decisive practical reason to have... \\
18-a 3 & ...t involves no false belief: the serious  & \textbf{agent} need not believe that her ends are cert... \\
18-b 1 & ...s psychologically necessary for a given  & \textbf{agent} is an empirical question, and the answe... \\
19-a 3 & ...quires that it fulfil the purpose of an  & \textbf{agent} who imposes it or that it lead to a fur... \\
\end{longtable}
\kwichead{answer}{34}
\begin{longtable}{@{}>{\raggedright}p{0.07\textwidth}>{\raggedleft}p{0.41\textwidth}@{\,}>{\raggedright\arraybackslash}p{0.47\textwidth}@{}}
01-d Abs & We  & \textbf{answer} the objections that the view is elitist... \\
02-a Abs & We then  & \textbf{answer} the symmetry argument of Lucretius, the... \\
02-a 1 & Sections 4 and 5  & \textbf{answer} the symmetry argument and the main obje... \\
02-b Abs & ...in of two inferences and show where the  & \textbf{answer} bears on it. \\
03-a 2 & ... complexity, and the project types that  & \textbf{answer} to it are those that engage that struct... \\
04-a 3 & ...at justifies that justification, and no  & \textbf{answer} is forced on that standpoint. \\
04-a 4 & ...nd to that question forgetting gives no  & \textbf{answer}. \\
07-a Abs & We argue that the  & \textbf{answer} is no. \\
07-a Abs & ...itoring one's happiness is distracting,  & \textbf{answer} the reply that a sophisticated hedonist... \\
07-a 4 & Two points  & \textbf{answer} it. \\
08-a 3 & ...s hers or her parents', and acts on the  & \textbf{answer}, satisfies (ii) and (iii) on that occas... \\
10-a 5 & ...ome a correct signal that nothing could  & \textbf{answer}. \\
10-b Abs & Both claims  & \textbf{answer} to empirical work. \\
10-b 8 & ...ome a correct signal that nothing could  & \textbf{answer}, and the trait findings of Corollary 3.... \\
11-a Abs & ...on as one about dependence, so that the  & \textbf{answer} is independent of whether God exists. \\
11-a 1 & ... everything ended in death there was no  & \textbf{answer} (Tolstoy 1882). \\
11-a 1 & ...such a world is a counterpossible whose  & \textbf{answer} is unclear. \\
11-a 1 & So read, the question has a determinate  & \textbf{answer} whether or not REL-01 is correct, and t... \\
11-a 1 & ...ether or not REL-01 is correct, and the  & \textbf{answer} also settles a question of evident prac... \\
11-a 6 & ...nd our conclusion is independent of its  & \textbf{answer}. \\
11-b 1 & ...d if true, and so it holds whatever the  & \textbf{answer} to REL-01. \\
12-a Abs & This manuscript completes the  & \textbf{answer}. \\
12-a Abs & Second, we  & \textbf{answer} Benatar's independent argument that our... \\
14-a 4 & ...conditions of its rationality, but they  & \textbf{answer} the charge that it is idle. \\
15-a 1 & The  & \textbf{answer} depends on what hope is. \\
15-b Abs & And we  & \textbf{answer} the objection that the virtue so descri... \\
16-a 1 & The oldest  & \textbf{answer} to the question why we love is that we ... \\
16-a 3 & ...structed, loves something that does not  & \textbf{answer} to the person. \\
16-b Abs & We extend the  & \textbf{answer} to three cases of substitution: trading... \\
16-b 1 & ...red to the reasons of the old, it would  & \textbf{answer} to relational reasons constituted by th... \\
18-b 1 & ...agent is an empirical question, and the  & \textbf{answer} may differ across agents and across the... \\
19-a Abs & We argue for a third  & \textbf{answer}. \\
20-a Abs & ...ion of the collection and show that the  & \textbf{answer} turns on two questions: whether the pro... \\
20-a 4 & ... a position to know either fact, and an  & \textbf{answer} to the question from that position woul... \\
\end{longtable}
\kwichead{answers}{32}
\begin{longtable}{@{}>{\raggedright}p{0.07\textwidth}>{\raggedleft}p{0.41\textwidth}@{\,}>{\raggedright\arraybackslash}p{0.47\textwidth}@{}}
02-a 4 & It  & \textbf{answers} the argument without denying the symmet... \\
02-b Abs & EXI-02-a  & \textbf{answers} the first claim: death is bad for the o... \\
02-b 2 & EXI-02-a  & \textbf{answers} (E1). \\
03-a 2 & ...apacities and values to which that type  & \textbf{answers}. \\
04-a 3 & If she  & \textbf{answers} by citing a further consideration, the ... \\
04-a 3 & If she  & \textbf{answers} that the consideration has objective wo... \\
07-a Abs & A long tradition  & \textbf{answers} yes: whatever else we pursue, we pursue... \\
10-b 1 & Section 9  & \textbf{answers} objections to the use we make of the st... \\
11-a 1 & Section 5  & \textbf{answers} objections. \\
12-a 2 & Principle P2  & \textbf{answers} the familiar worry that coming into exi... \\
13-a 1 & ...s can affect whether she was happy, and  & \textbf{answers} that they can, though only slightly (Ar... \\
13-a 2 & Assumption 2.2  & \textbf{answers} that question for the case that matters... \\
13-b Abs & The distinction  & \textbf{answers} the question, left open by ETH-16-b, wh... \\
13-b 2 & This  & \textbf{answers} ETH-16-b in the sense it requires: the ... \\
16-a 2 & ... place of with no loss in what her love  & \textbf{answers} to. \\
16-a 3 & ...therefore lose a reason that 's love of  & \textbf{answers} to. \\
16-a 3 & Every established love  & \textbf{answers} to relational reasons. \\
16-a 3 & By Lemma 3.4, every established love  & \textbf{answers} to relational reasons, and by Lemma 3.3... \\
16-a 3 & ... reasons, and by Lemma 3.3, a love that  & \textbf{answers} to a relational reason is not transfera... \\
16-a 3 & ...rplus of qualities, since what the love  & \textbf{answers} to is not on offer elsewhere. \\
16-a 4 & ...oes not share the history that the love  & \textbf{answers} to. \\
16-b Abs & EXI-16-a  & \textbf{answers} that established love answers to relati... \\
16-b Abs & EXI-16-a answers that established love  & \textbf{answers} to relational reasons, which no qualita... \\
16-b Abs & ...l reason, and in the third the new love  & \textbf{answers} to a relational reason of its own rathe... \\
16-b 1 & Since every established love  & \textbf{answers} to relational reasons (Lemma 3.4), esta... \\
16-b 1 & ...3.4 of EXI-16-a, 's established love of  & \textbf{answers} to relational reasons. \\
16-b 1 & It  & \textbf{answers} to a relational reason of its own, cons... \\
16-b 1 & ...a, once it is established, the new love  & \textbf{answers} to a relational reason constituted by t... \\
16-b 3 & ... history, and a new love, if one comes,  & \textbf{answers} to a history of its own. \\
19-a Abs & Two  & \textbf{answers} are common. \\
19-a 4 & Corollary 3.3  & \textbf{answers} this. \\
20-a 3 & ...on, not philosophy, and a collection of  & \textbf{answers} is worthwhile to the extent that it con... \\
\end{longtable}
\kwichead{attitudes}{31}
\begin{longtable}{@{}>{\raggedright}p{0.07\textwidth}>{\raggedleft}p{0.41\textwidth}@{\,}>{\raggedright\arraybackslash}p{0.47\textwidth}@{}}
01-c 2 & ...lue that does not depend on the agent's  & \textbf{attitudes} toward it (Definition 2.2 of EXI-17-a). \\
01-c 2 & ...y, and that its value not depend on the  & \textbf{attitudes} of those engaged in it (Assumption 2 of... \\
01-c 3 & ...whose good has value independent of her  & \textbf{attitudes}, and every agent lives in a world some ... \\
01-d 4 & ...he value of a project not depend on the  & \textbf{attitudes} of those engaged in it. \\
04-a Abs & ... is to treat responses to the absurd as  & \textbf{attitudes} assessable for fit. \\
04-a 1 & ... may be worthwhile independently of our  & \textbf{attitudes} toward them, and then the step back see... \\
04-a 2 & ...o understood, is a relation between two  & \textbf{attitudes} of the same subject. \\
04-a 4 & ...cts are worthwhile independently of our  & \textbf{attitudes}, as EXI-01 holds, then the seriousness ... \\
06-a Abs & Care and hostility are  & \textbf{attitudes}, and both require a subject whose respo... \\
06-a 3 & Their  & \textbf{attitudes} are realized physically, and the theore... \\
06-a 3 & ...are not composed so as to aggregate the  & \textbf{attitudes} of the agents within them into an attit... \\
06-a 4 & ... the absence of the responsiveness that  & \textbf{attitudes} of care and hostility require. \\
06-b Abs & ...hat has worth independently of anyone's  & \textbf{attitudes}, and the universe's want of attitudes t... \\
06-b Abs & ...s attitudes, and the universe's want of  & \textbf{attitudes} toward us takes nothing from that worth... \\
06-b 2 & ...rth, whose value does not depend on the  & \textbf{attitudes} of those engaged in them (Definition 1 ... \\
06-b 2 & ...ts value does not depend on the agent's  & \textbf{attitudes} toward it (Definition 2.2 of EXI-17-a). \\
06-b 2 & ...rth of a project does not depend on the  & \textbf{attitudes} of those engaged in it, it does not dep... \\
06-b 2 & ...ngaged in it, it does not depend on the  & \textbf{attitudes} of the universe either, and a fortiori ... \\
06-b 2 & ... fortiori not on the universe's lack of  & \textbf{attitudes}. \\
06-b 2 & Both  & \textbf{attitudes} misrepresent their object. \\
06-b 3 & ...d cannot be met by any universe without  & \textbf{attitudes}. \\
07-a Abs & ...for its own sake, and so lacks the very  & \textbf{attitudes} on which most of her happiness depends. \\
07-a 2 & ...t is a claim about the structure of the  & \textbf{attitudes}, not about their phenomenology. \\
07-a 3 & ... else is a means to it, still lacks the  & \textbf{attitudes} that Principle P2 requires. \\
07-a 4 & ... other things, she has not acquired the  & \textbf{attitudes} on which engagement depends, and Theore... \\
07-a 6 & She therefore lacks the  & \textbf{attitudes} on which most of her happiness depends,... \\
17-a 2 & ...lue that does not depend on the agent's  & \textbf{attitudes} toward it. \\
18-a Abs & ...tment, which do not misrepresent, those  & \textbf{attitudes} are to be preferred. \\
18-a 2 & ... to positive illusions can be served by  & \textbf{attitudes} that do not misrepresent. \\
18-a 2 & To the extent that these  & \textbf{attitudes} serve the functions, the range of neces... \\
18-a 4 & ...cal commitment can do their work, those  & \textbf{attitudes}, which do not misrepresent, are better. \\
\end{longtable}
\kwichead{available}{61}
\begin{longtable}{@{}>{\raggedright}p{0.07\textwidth}>{\raggedleft}p{0.41\textwidth}@{\,}>{\raggedright\arraybackslash}p{0.47\textwidth}@{}}
01-c Abs & ...shes that a life of positive meaning is  & \textbf{available} to every agent capable of intentional a... \\
01-c 2 & ...l action, a life of positive meaning is  & \textbf{available}. \\
01-d Abs & ...orth, and that meaning in this sense is  & \textbf{available} to everyone. \\
01-d 1 & ...al action a life of positive meaning is  & \textbf{available}. \\
01-d 2 & ...f meaning required success, it would be  & \textbf{available} only to those whom the world allows to ... \\
01-d 2 & ...it, but not whether positive meaning is  & \textbf{available}. \\
01-d 4 & ...causes, and so make meaning in practice  & \textbf{available} mainly to the privileged. \\
01-d 4 & If meaning is  & \textbf{available} to everyone, the claim that a life is m... \\
01-d 4 & That positive meaning is  & \textbf{available} to everyone is compatible with its bein... \\
01-d 4 & ...andpoint nothing matters, so meaning is  & \textbf{available} to no one. \\
02-a 2 & ... theory of persons but that no evidence  & \textbf{available} before death could confirm it. \\
03-a 4 & Over an endless life the new theorems  & \textbf{available} to her become, in the respects that eng... \\
04-a Abs & ...ness of our commitments and the support  & \textbf{available} for them from the most detached standpo... \\
04-a 2 & ...no justification of that seriousness is  & \textbf{available} that is not itself open to the same que... \\
04-a 2 & ...em and the support for that seriousness  & \textbf{available} from the most detached standpoint, whil... \\
04-a 3 & ...no justification of that seriousness is  & \textbf{available} that is not open to the same question. \\
04-a 3 & ...holds one's ends and the support for it  & \textbf{available} from the most detached standpoint. \\
04-b 1 & ...no justification of that seriousness is  & \textbf{available} that is not itself open to the same que... \\
04-b 2 & ...ness of that engagement and the support  & \textbf{available} for it from the most detached standpoin... \\
05-a 3 & ...at end is coherent but that no evidence  & \textbf{available} before death could confirm it; the poss... \\
07-a 2 & ...t pleasure in the friend, and it is not  & \textbf{available} to someone who treats the friend as an ... \\
07-a 3 & ... that is realized through engagement is  & \textbf{available} only to agents who do. \\
07-a 3 & ...s yields less of its own target than an  & \textbf{available} alternative, which is what it is for a ... \\
07-a 4 & First, no such route is  & \textbf{available} to human beings, and the question wheth... \\
08-a 1 & ...ublic interpretation of the world makes  & \textbf{available} (Heidegger 1927). \\
10-a 2 & ...y when the activity does not occupy her  & \textbf{available} attention and abilities in a way she fi... \\
10-a 3 & ...n engagement is restored selects, among  & \textbf{available} alternatives, those that engage. \\
10-a 3 & ...correct but the motivated search has no  & \textbf{available} alternatives to find. \\
10-a 4 & ...state with the effects of the responses  & \textbf{available} to it. \\
10-b 2 & ... mismatch between cognitive demands and  & \textbf{available} resources in either direction, from a m... \\
10-b 4 & ..., and in a laboratory in which the only  & \textbf{available} stimulation is a shock, the search ends... \\
11-a 1 & Naturalists reply that meaning is  & \textbf{available} in a world without God. \\
11-a 2 & ...I-01 adds that meaning so understood is  & \textbf{available} to everyone. \\
11-a 4 & ...theistic framework, if true, would make  & \textbf{available} to a life more meaning than it could ha... \\
11-a 4 & ...ght, the additional meaning is actually  & \textbf{available}, and is available to those who engage w... \\
11-a 4 & ...l meaning is actually available, and is  & \textbf{available} to those who engage with God; if REL-01... \\
11-a 4 & ...with God; if REL-01 is wrong, it is not  & \textbf{available}, and Theorem 3.3 ensures that life is m... \\
11-a 4 & ...l of death is coherent, but no evidence  & \textbf{available} before death could confirm it. \\
11-a 4 & ...ralism that God would make more meaning  & \textbf{available} (Metz 2013); it does not follow the mod... \\
11-b Abs & ...e, would nevertheless make more meaning  & \textbf{available} (Proposition 4.1). \\
11-b Abs & ...ar life-plans, not the objects of worth  & \textbf{available} to engagement, so that Proposition 4.1 ... \\
11-b 1 & ... add is of the same kind as the meaning  & \textbf{available} without it, and differs in degree. \\
11-b 2 & ...existence, so the objects of engagement  & \textbf{available} without God remain available with God. \\
11-b 2 & ...engagement available without God remain  & \textbf{available} with God. \\
11-b 2 & ...e proposition concerns the meaning made  & \textbf{available} by a framework, and the framework adds ... \\
11-b 2 & ...w that the framework makes less meaning  & \textbf{available}. \\
11-b 3 & ...eism, if true, would add to the meaning  & \textbf{available} in a life by adding objects and relatio... \\
14-a Abs & ... of one's past as good and as no longer  & \textbf{available}, and it is fitting when the past contai... \\
14-a 2 & ...ined a good, and that good as no longer  & \textbf{available} to the subject. \\
14-a 2 & ...ned the good, and the good is no longer  & \textbf{available}), and proportionate (its intensity and ... \\
14-a 3 & ...real good of the past that is no longer  & \textbf{available}. \\
14-a 3 & ... as containing a good that is no longer  & \textbf{available}. \\
14-a 3 & ...id contain the good and it is no longer  & \textbf{available}. \\
14-a 3 & ...real good of the past that is no longer  & \textbf{available} and is proportionate to that good. \\
14-a 3 & ... tracking of a real past good no longer  & \textbf{available}. \\
14-a 3 & ... the good warrants, or forecloses goods  & \textbf{available} now. \\
17-a 3 & ...of EXI-01 that meaning in this sense is  & \textbf{available} to everyone. \\
18-a 2 & ...takes of error are high and evidence is  & \textbf{available}, that is a culpable failure of inquiry ... \\
19-a 3 & ...hed, and REL-08 argues that no evidence  & \textbf{available} before death could confirm survival. \\
20-a 1 & ...I-01 adds that meaning in this sense is  & \textbf{available} to everyone. \\
20-a 2 & On the evidence  & \textbf{available}, the condition of engagement is neither... \\
\end{longtable}
\kwichead{capable}{15}
\begin{longtable}{@{}>{\raggedright}p{0.07\textwidth}>{\raggedleft}p{0.41\textwidth}@{\,}>{\raggedright\arraybackslash}p{0.47\textwidth}@{}}
01-c Abs & ...ive meaning is available to every agent  & \textbf{capable} of intentional action. \\
01-c 1 & ...mption 2), and assumes that every agent  & \textbf{capable} of intentional action is able to engage... \\
01-c 2 & Every agent  & \textbf{capable} of intentional action can direct some i... \\
01-c 2 & For every agent  & \textbf{capable} of intentional action, a life of positi... \\
01-c 4 & Some agents  & \textbf{capable} of intentional action are wholly isolat... \\
01-c 4 & An agent  & \textbf{capable} of intentional action acts on something... \\
01-d 1 & ...of EXI-01-a states that for every agent  & \textbf{capable} of intentional action a life of positiv... \\
01-d 2 & No agent  & \textbf{capable} of intentional action is condemned by h... \\
04-a 3 & Any agent  & \textbf{capable} of intentional action takes some ends s... \\
04-a 3 & The life of any reflective agent  & \textbf{capable} of intentional action is absurd in the ... \\
09-a 2 & For a being  & \textbf{capable} of adopting ends, what its life is for,... \\
09-a 3 & ...fixes a defining project for , since is  & \textbf{capable} of adopting ends and what 's life is fo... \\
09-a 3 & ...s kind fix one: the kind human organism  & \textbf{capable} of thought is fixed by biological origi... \\
09-a 4 & ...ions fix a defining project for a being  & \textbf{capable} of adopting ends. \\
09-a 5 & ...he result concerns artefacts and beings  & \textbf{capable} of adopting ends. \\
\end{longtable}
\kwichead{challenge}{10}
\begin{longtable}{@{}>{\raggedright}p{0.07\textwidth}>{\raggedleft}p{0.41\textwidth}@{\,}>{\raggedright\arraybackslash}p{0.47\textwidth}@{}}
01-d 4 & ...Assumption 2 denies the premise of this  & \textbf{challenge} rather than refuting it (Nagel 1971). \\
04-a 4 & ...rth denies the premise of the sceptical  & \textbf{challenge} without refuting it. \\
10-a 2 & ...sorption and learning are maximal where  & \textbf{challenge} matches skill. \\
10-b Abs & ...as the weakest support and one standing  & \textbf{challenge}, the flow model's claim that over-chall... \\
10-b 3 & ...arch places boredom where skill exceeds  & \textbf{challenge} (Csikszentmihalyi 1975; Csikszentmihaly... \\
10-b 3 & ...egative emotions by a perceived lack of  & \textbf{challenge} (Til12). \\
10-b 3 & ...se of the mismatch thesis, that matched  & \textbf{challenge} and skill engage, is well supported: ex... \\
10-b 3 & ...of experience related to the balance of  & \textbf{challenge} and skill (Csikszentmihalyi and LeFevre... \\
10-b 4 & ...nder the strict rule for costs, entries  & \textbf{challenge} the function thesis outright and bear i... \\
10-b 7 & ...redom, the activity, and its subjective  & \textbf{challenge} and skill, and test whether boredom pre... \\
\end{longtable}
\kwichead{consists}{18}
\begin{longtable}{@{}>{\raggedright}p{0.07\textwidth}>{\raggedleft}p{0.41\textwidth}@{\,}>{\raggedright\arraybackslash}p{0.47\textwidth}@{}}
01-c 1 & ...ife as meaningful to the degree that it  & \textbf{consists} in active engagement with projects of o... \\
01-c 2 & A life that includes such engagement  & \textbf{consists} to a positive extent in active engageme... \\
04-a Abs & ... discrepancy is exactly what the absurd  & \textbf{consists} in. \\
04-a 4 & ...ony does not deny the contradiction; it  & \textbf{consists} in seeing it. \\
06-b Abs & Meaning, on the account of EXI-01,  & \textbf{consists} in engagement with what has worth indep... \\
06-b 3 & ...resent point is conditional: if meaning  & \textbf{consists} in engagement with what has objective w... \\
07-a Abs & Happiness, on any plausible account,  & \textbf{consists} to a large extent in engagement with an... \\
07-a 1 & ...urce of our enjoyments, since enjoyment  & \textbf{consists} in the gratification of particular pass... \\
07-a 2 & A large part of human happiness  & \textbf{consists} in, or depends causally on, wholehearte... \\
07-a 6 & Much of it  & \textbf{consists} in, or depends on, engagement with acti... \\
11-a 7 & Meaning  & \textbf{consists} in active engagement with what is worth... \\
15-a 2 & ... is under the agent's control, since it  & \textbf{consists} in what she attends to, plans for and a... \\
15-a 3 & ...eparture from epistemic rationality; it  & \textbf{consists} in the governance of the stance, which ... \\
15-b 1 & The virtue  & \textbf{consists} in proportion without collapse (Definit... \\
16-a 2 & It is relational if it  & \textbf{consists} in a relationship between and constitut... \\
16-b 2 & ...ed to reasons for action: a reason that  & \textbf{consists} in features anyone could have generaliz... \\
16-b 2 & ...ciple P2 of EXI-16-a, and a reason that  & \textbf{consists} in a shared history does not. \\
18-a 2 & ...on, and the ordinary progress of a life  & \textbf{consists} partly in such corrections. \\
\end{longtable}
\kwichead{defence}{1}
\begin{longtable}{@{}>{\raggedright}p{0.07\textwidth}>{\raggedleft}p{0.41\textwidth}@{\,}>{\raggedright\arraybackslash}p{0.47\textwidth}@{}}
08-a 1 & We argue that this  & \textbf{defence} succeeds only if authenticity is unders... \\
\end{longtable}
\kwichead{degree}{31}
\begin{longtable}{@{}>{\raggedright}p{0.07\textwidth}>{\raggedleft}p{0.41\textwidth}@{\,}>{\raggedright\arraybackslash}p{0.47\textwidth}@{}}
01-c 1 & ...1-a defines a life as meaningful to the  & \textbf{degree} that it consists in active engagement w... \\
01-c 2 & ...EXI-01-a it is meaningful to a positive  & \textbf{degree}. \\
01-d Abs & ...ircumstance and injustice affect is the  & \textbf{degree} of meaning, not its possibility. \\
01-d 2 & ...alue a project serves, which can affect  & \textbf{degree} but not presence. \\
01-d 2 & ...tances, including injustice, affect the  & \textbf{degree} of meaning a life can attain and the ea... \\
01-d 2 & ...1 of EXI-01-a makes meaning a matter of  & \textbf{degree}, and EXI-17-a notes that achievement ma... \\
01-d 2 & ...a notes that achievement may add to its  & \textbf{degree} (Remark 3.5 of EXI-17-a). \\
01-d 3 & ...ond, from availability to equality: the  & \textbf{degree} of meaning varies with the extent of en... \\
01-d 3 & ... tolerate oppression, which reduces the  & \textbf{degree} of meaning as it reduces other goods. \\
01-d 5 & ...ailability of meaning is universal; its  & \textbf{degree} is not, and how much of it a life attai... \\
02-a 3 & The  & \textbf{degree} of badness varies. \\
03-a 2 & There is a  & \textbf{degree} of connectedness below which a future l... \\
05-a 3 & ...s one would otherwise have had, and the  & \textbf{degree} of its badness is the value of what it ... \\
05-a 3 & ...eprive one of goods, and fitting to the  & \textbf{degree} that it does. \\
11-b Abs & ...e additional meaning is a difference of  & \textbf{degree}, and that the anti-theist's losses conc... \\
11-b 1 & ...ng available without it, and differs in  & \textbf{degree}. \\
11-b 1 & ...ple P2 of EXI-11-a, it can increase the  & \textbf{degree} of meaning. \\
11-b 2 & ...ition 1.1 counts that as an addition of  & \textbf{degree}. \\
11-b 3 & The addition is one of  & \textbf{degree}, which is why its absence leaves meanin... \\
15-a 2 & ...he right circumstances and to the right  & \textbf{degree}, flanked by vices of deficiency and exc... \\
15-a 3 & It is acting to the wrong  & \textbf{degree} (Definition 2.2). \\
15-a 3 & ...lar form, and it is acting to the wrong  & \textbf{degree} in the other direction. \\
15-a 3 & ...e other things that determine the right  & \textbf{degree} of investment in an uncertain outcome. \\
15-a 3 & ...isposition to act and feel to the right  & \textbf{degree} in the right circumstances, between two... \\
16-a 2 & ...n for 's love of , to at least the same  & \textbf{degree}. 's love of is transferable if, for eve... \\
16-a 2 & ...rds , the same feature, had to the same  & \textbf{degree} by , is an equally strong reason for th... \\
16-a 3 & ...n for 's love of , to at least the same  & \textbf{degree}. \\
16-a 3 & If has some features to a greater  & \textbf{degree}, the same argument yields trading up. \\
16-b 1 & ...e qualities for which loves to a higher  & \textbf{degree} is not a reason for to love that person... \\
16-b 1 & ...itative ones, and whatever they gain in  & \textbf{degree}, they lack a kind of reason that 's lov... \\
17-a Abs & ...adds a further good, and may add to the  & \textbf{degree} of meaning; it is not a condition of me... \\
\end{longtable}
\kwichead{denies}{14}
\begin{longtable}{@{}>{\raggedright}p{0.07\textwidth}>{\raggedleft}p{0.41\textwidth}@{\,}>{\raggedright\arraybackslash}p{0.47\textwidth}@{}}
01-d 4 & ...k 5 of EXI-01-a notes that Assumption 2  & \textbf{denies} the premise of this challenge rather th... \\
04-a 2 & ...presents the world as an adversary that  & \textbf{denies} us meaning, and the agent's persistence... \\
04-a 3 & ...presents the world as an adversary that  & \textbf{denies} us meaning. \\
04-a 3 & An indifferent world  & \textbf{denies} us nothing and cannot be defeated by ou... \\
04-a 4 & ...s it, the assumption of objective worth  & \textbf{denies} the premise of the sceptical challenge ... \\
06-a 2 & ...require responsiveness and indifference  & \textbf{denies} it. \\
06-b 2 & It  & \textbf{denies} that meaning in life requires that good... \\
06-b 2 & By Proposition 2.1, its indifference  & \textbf{denies} us nothing that meaning requires; by Th... \\
09-a 1 & Sartre's atheist existentialism  & \textbf{denies} this: there is at least one being whose... \\
09-a 4 & ...ay precede , including a creator's, and  & \textbf{denies} only that such intentions fix a definin... \\
11-a 2 & ...s valuable, or command what is good; it  & \textbf{denies} that the assignment is what makes it go... \\
12-a 4 & ... apply to them; nothing in our argument  & \textbf{denies} that some lives are worse than non-exis... \\
12-b 2 & ...e is bad for no one, and not on what it  & \textbf{denies}, that this absence is neutral. \\
16-a 1 & Frankfurt (2004)  & \textbf{denies} that love is a response to reasons: lov... \\
\end{longtable}
\kwichead{depend}{36}
\begin{longtable}{@{}>{\raggedright}p{0.07\textwidth}>{\raggedleft}p{0.41\textwidth}@{\,}>{\raggedright\arraybackslash}p{0.47\textwidth}@{}}
01-c 2 & ...if it has or serves value that does not  & \textbf{depend} on the agent's attitudes toward it (Def... \\
01-c 2 & ...es in, actively, and that its value not  & \textbf{depend} on the attitudes of those engaged in it... \\
01-c 3 & ...ties for projects in the ordinary sense  & \textbf{depend} on resources, institutions and freedom,... \\
01-c 5 & ...versal availability of meaning does not  & \textbf{depend} on the universal availability of projec... \\
01-d 1 & ...I-01-c shows that availability does not  & \textbf{depend} on the availability of projects in the ... \\
01-d 2 & Whether a life is meaningful does not  & \textbf{depend} on whether its projects succeed. \\
01-d 4 & ...es only that the value of a project not  & \textbf{depend} on the attitudes of those engaged in it... \\
05-a 3 & Its fittingness does not  & \textbf{depend} on uncertainty, which Lemma 3.1 exclude... \\
05-a 4 & ...he asymmetry of fitting regret does not  & \textbf{depend} on it. \\
06-b 2 & ...f objective worth, whose value does not  & \textbf{depend} on the attitudes of those engaged in th... \\
06-b 2 & ...ent is worthwhile if its value does not  & \textbf{depend} on the agent's attitudes toward it (Def... \\
06-b 2 & If the worth of a project does not  & \textbf{depend} on the attitudes of those engaged in it... \\
06-b 2 & ...des of those engaged in it, it does not  & \textbf{depend} on the attitudes of the universe either... \\
07-a Abs & ...is own crisis: the pleasures of pursuit  & \textbf{depend} on caring about objects other than plea... \\
07-a 2 & The argument does not  & \textbf{depend} on a choice between hedonistic, life-sa... \\
07-a 3 & Theorem 3.3 does not  & \textbf{depend} on attention. \\
08-a 4 & ...ne has it, and its possibility does not  & \textbf{depend} on its being known. \\
09-a Abs & The result does not  & \textbf{depend} on atheism, and it is compatible with t... \\
09-a 6 & The argument does not  & \textbf{depend} on atheism, and it leaves intact the th... \\
09-b 2 & ... of success as what it is that does not  & \textbf{depend} on how it is used; a knife that cannot ... \\
10-a 2 & ...reatures whose survival and flourishing  & \textbf{depend} on learning, exploration and the exerci... \\
11-a 3 & ...orsement converges on, does not make it  & \textbf{depend} on any assignment. \\
11-a 4 & Meaning has a floor that does not  & \textbf{depend} on God, and theism, if true, would rais... \\
11-a 5 & ...at it is would make the worth of a good  & \textbf{depend} on what happens after it, which neither... \\
11-b Abs & ...irection on which some meaningful lives  & \textbf{depend}. \\
11-b 2 & ...edge, of persons and of beauty does not  & \textbf{depend} on a role assigned by a designer (Lemma... \\
11-b 2 & ...names are real and that some life-plans  & \textbf{depend} on them. \\
15-b Abs & ... credence, other things equal, but must  & \textbf{depend} also on the value of the outcome and on... \\
17-a 2 & ...if it has or serves value that does not  & \textbf{depend} on the agent's attitudes toward it. \\
18-a Abs & ...t the believer; truth and justification  & \textbf{depend} on the world and on the evidence. \\
18-a 3 & Our conclusion does not  & \textbf{depend} on it. \\
19-a Abs & ...elleman, we take the value of a life to  & \textbf{depend} in part on its shape, the order in whic... \\
19-a 3 & ...ess life is not desirable, but does not  & \textbf{depend} on it. \\
19-b 1 & ...beauty and the good of others, does not  & \textbf{depend} on how long those who engage them live:... \\
20-a 4 & ...akes the meaningfulness of the activity  & \textbf{depend} on two facts that cannot be established... \\
20-b Abs & ...-20-a makes the worth of the collection  & \textbf{depend} on whether it contributes to knowledge ... \\
\end{longtable}
\kwichead{depends}{53}
\begin{longtable}{@{}>{\raggedright}p{0.07\textwidth}>{\raggedleft}p{0.41\textwidth}@{\,}>{\raggedright\arraybackslash}p{0.47\textwidth}@{}}
01-c 5 & It  & \textbf{depends} on something more widely distributed: t... \\
01-d 1 & ...ny particular life is meaningful, which  & \textbf{depends} on what is done with it. \\
01-d 4 & ...to high culture; the availability claim  & \textbf{depends} on this, and the cases in Proposition 2... \\
01-d 5 & ... not, and how much of it a life attains  & \textbf{depends} on what is done with it and on the cond... \\
02-a 2 & ... for a while as a corpse; nothing below  & \textbf{depends} on it. \\
02-a 2 & ... account is comparative: whether is bad  & \textbf{depends} on what the alternative to would have b... \\
02-a 6 & ... settle how much the badness of a death  & \textbf{depends} on the psychological connections betwee... \\
03-a 1 & ...that the value of many human activities  & \textbf{depends} on the scarcity of time. \\
03-a 4 & ...riendship is a project, and its renewal  & \textbf{depends} on the friends' changing. \\
06-a 2 & ...onsive to our interests if what it does  & \textbf{depends}, in a counterfactually robust way, on h... \\
06-a 2 & ...onditions on which the thing's behavior  & \textbf{depends}. \\
06-a 3 & ... Assumption 2.3, what the universe does  & \textbf{depends} only on its prior states and on laws th... \\
07-a Abs & ...ttitudes on which most of her happiness  & \textbf{depends}. \\
07-a 2 & ...ional-state accounts (Haybron 2008); it  & \textbf{depends} only on what makes people happy, on any... \\
07-a 2 & ...part of human happiness consists in, or  & \textbf{depends} causally on, wholehearted engagement in... \\
07-a 3 & ...orgoes the large part of happiness that  & \textbf{depends} on engagement, and attains less happine... \\
07-a 4 & ...uired the attitudes on which engagement  & \textbf{depends}, and Theorem 3.3 applies to her. \\
07-a 6 & Much of it consists in, or  & \textbf{depends} on, engagement with activities and peop... \\
07-a 6 & ...ttitudes on which most of her happiness  & \textbf{depends}, and her aim defeats itself. \\
07-b 1 & ...s the absorption on which the enjoyment  & \textbf{depends}. \\
07-b 1 & ...ect-directed caring on which engagement  & \textbf{depends} (Principle P2 of EXI-07-a). \\
07-b 2 & ...ich case she forgoes the happiness that  & \textbf{depends} on engagement. \\
08-a 2 & ...n, and holds it in a way that no longer  & \textbf{depends} on its being what one does. \\
08-a 3 & ...mitment is held in a way that no longer  & \textbf{depends} on its being what one does. \\
08-a 3 & ...from shared meanings, so her holding it  & \textbf{depends} on public interpretations, and the appr... \\
09-a 2 & ...ation on which 's coming into existence  & \textbf{depends}. \\
09-a 3 & ... out, and 's coming into existence as a  & \textbf{depends} on it. \\
09-b 2 & ... a determination on which its existence  & \textbf{depends}, not that the determination mention it ... \\
11-b 2 & ...eaning is undercut is the claim that it  & \textbf{depends} on a cosmic purpose, which Theorem 3.3 ... \\
15-a 1 & The answer  & \textbf{depends} on what hope is. \\
15-a 3 & ...t a function of the credence alone, but  & \textbf{depends} also on the value of and on the agent's... \\
16-b 1 & Whether the duplicate took part in it  & \textbf{depends} on whether she is the beloved, which is... \\
17-a 3 & Whether an activity succeeds  & \textbf{depends} on circumstances outside the agent's co... \\
17-a 3 & ...he is engaged with something worthwhile  & \textbf{depends} far less on such circumstances. \\
17-a 3 & By Theorem 3.3, meaning  & \textbf{depends} on the second and not on the first. \\
18-a 2 & Whether a belief is true  & \textbf{depends} on whether things are as it represents ... \\
18-a 2 & ...epresents them; whether it is justified  & \textbf{depends} on the evidence bearing on it. \\
18-a 3 & Meaning  & \textbf{depends} on taking one's projects to matter more... \\
18-a 4 & ...ically necessary, in that normal agency  & \textbf{depends} on them and they cannot be removed at w... \\
19-a Abs & ...f, since every candidate meaning either  & \textbf{depends} on claims that cannot be established or... \\
19-a 1 & ...hat the value of many of our activities  & \textbf{depends} on the scarcity of time. \\
19-a 2 & The value of a life as a whole  & \textbf{depends} in part on its shape, and not only on t... \\
19-a 3 & ...s not determine what the shape is: that  & \textbf{depends} on the goods and evils the life contain... \\
19-b Abs & ...aning, that of a life as a whole, which  & \textbf{depends} on the life's shape. \\
19-b 1 & ...far as the meaning of a life as a whole  & \textbf{depends} on its shape, finitude is a condition o... \\
19-b 1 & Such meaning  & \textbf{depends} on the temporal order of the engagement... \\
19-b 2 & ...ffler (2013) that our capacity to value  & \textbf{depends} on the awareness of temporal scarcity, ... \\
20-a Abs & The first  & \textbf{depends} on whether the process has the kind of ... \\
20-a Abs & The second  & \textbf{depends} on what readers do with the collection,... \\
20-a 3 & ...nts are correct, the collection's worth  & \textbf{depends} on the minority of readers who engage w... \\
20-a 3 & ...h its arguments, and its worth for them  & \textbf{depends} on whether the arguments are good. \\
20-b Abs & ...g the present collection was meaningful  & \textbf{depends}, by the account of EXI-17, on whether i... \\
20-b Abs & ...ucing process was engaged with its work  & \textbf{depends}, by EXI-20-a, on whether it has the men... \\
\end{longtable}
\kwichead{discussion}{2}
\begin{longtable}{@{}>{\raggedright}p{0.07\textwidth}>{\raggedleft}p{0.41\textwidth}@{\,}>{\raggedright\arraybackslash}p{0.47\textwidth}@{}}
02-a 1 & ...he argument has two premises that later  & \textbf{discussion} has separated. \\
04-a 1 & Two lines of later  & \textbf{discussion} bear on the recommendation. \\
\end{longtable}
\kwichead{done}{11}
\begin{longtable}{@{}>{\raggedright}p{0.07\textwidth}>{\raggedleft}p{0.41\textwidth}@{\,}>{\raggedright\arraybackslash}p{0.47\textwidth}@{}}
01-d 1 & ...is meaningful, which depends on what is  & \textbf{done} with it. \\
01-d 5 & ...of it a life attains depends on what is  & \textbf{done} with it and on the conditions under whi... \\
01-d 5 & ...and on the conditions under which it is  & \textbf{done} (Metz 2013). \\
03-a 2 & ...y as repetition of what she has already  & \textbf{done}. \\
03-a 2 & ...e desire to repeat what one has already  & \textbf{done} is not a desire that gives life its poi... \\
03-b 2 & ... without specification of what is to be  & \textbf{done} in it. \\
04-b 4 & ... it does not undo the value of what was  & \textbf{done}. \\
13-a 2 & ...owing that wants assurance that will be  & \textbf{done}, then is obligated to to do , unless re... \\
14-a 4 & Nothing can be  & \textbf{done} about the past. \\
17-a 3 & ...jectivism fails because worthwhile work  & \textbf{done} without engagement, out of mere duty or... \\
17-a 3 & ...before may be difficult and competently  & \textbf{done}, and is therefore an achievement by Def... \\
\end{longtable}
\kwichead{engage}{27}
\begin{longtable}{@{}>{\raggedright}p{0.07\textwidth}>{\raggedleft}p{0.41\textwidth}@{\,}>{\raggedright\arraybackslash}p{0.47\textwidth}@{}}
01-c Abs & ...mption that every such agent is able to  & \textbf{engage} in at least one project of objective wo... \\
01-c 1 & ...apable of intentional action is able to  & \textbf{engage} in at least one of them (Assumption 3). \\
03-a 2 & ...that give life its point are desires to  & \textbf{engage} in such projects. \\
03-a 2 & ...ct types that a person of character can  & \textbf{engage} with, in the sense of finding them wort... \\
03-a 2 & ... types that answer to it are those that  & \textbf{engage} that structure. \\
03-a 3 & Projects of other types do not  & \textbf{engage} her, by Assumption 2.3. \\
03-a 4 & ...ble to her become, in the respects that  & \textbf{engage} her, instances of kinds she has already... \\
06-b 4 & The worth of what we  & \textbf{engage} with does not wait on cosmic recognitio... \\
06-b 4 & ...ning where EXI-01 places it: in what we  & \textbf{engage} with, and how. \\
10-a Abs & ...nts the current situation as failing to  & \textbf{engage} one's attention and abilities; it is tr... \\
10-a 1 & ...rience of wanting, but being unable, to  & \textbf{engage} in satisfying activity, and locate its ... \\
10-a 2 & ...ent activity or situation as failing to  & \textbf{engage} her, together with a desire to be engag... \\
10-a 2 & ...nts the current situation as failing to  & \textbf{engage}. \\
10-a 3 & ...ents the current activity as failing to  & \textbf{engage} the subject. \\
10-a 3 & An activity fails to  & \textbf{engage} an agent who is not otherwise impaired ... \\
10-a 3 & ...mong available alternatives, those that  & \textbf{engage}. \\
10-b 2 & ...rience of wanting, but being unable, to  & \textbf{engage} in satisfying activity, and locate the ... \\
10-b 3 & ...le an over-demanding task one can still  & \textbf{engage} with is a source of anxiety. \\
10-b 3 & ...hesis, that matched challenge and skill  & \textbf{engage}, is well supported: experience sampling... \\
10-b 11 & ...ople seek alternatives, change how they  & \textbf{engage}, and turn to sources of meaning. \\
11-a 4 & ...vailable, and is available to those who  & \textbf{engage} with God; if REL-01 is wrong, it is not... \\
11-b 1 & ...availability: it says what a life could  & \textbf{engage} with, not that every life in a theistic... \\
19-b Abs & ...ng, since it adds no object of worth to  & \textbf{engage} with. \\
19-b 1 & ..., does not depend on how long those who  & \textbf{engage} them live: a proof is no less a proof f... \\
19-b 2 & ...eir urgency, and urgency affects how we  & \textbf{engage}. \\
19-b 2 & ...ition 1.1 concerns the worth of what we  & \textbf{engage} with, and urgency adds no worth: the fr... \\
20-a 3 & ... depends on the minority of readers who  & \textbf{engage} with its arguments, and its worth for t... \\
\end{longtable}
\kwichead{engaged}{34}
\begin{longtable}{@{}>{\raggedright}p{0.07\textwidth}>{\raggedleft}p{0.41\textwidth}@{\,}>{\raggedright\arraybackslash}p{0.47\textwidth}@{}}
01-c 2 & ...ue not depend on the attitudes of those  & \textbf{engaged} in it (Assumption 2 of EXI-01-a). \\
01-d 4 & ...ct not depend on the attitudes of those  & \textbf{engaged} in it. \\
03-a 4 & A person  & \textbf{engaged} in inquiry would never be reduced to re... \\
04-a 4 & ...m is justified, and the life of someone  & \textbf{engaged} in them is not absurd. \\
04-a 4 & The  & \textbf{engaged} and worthwhile life is less exposed to ... \\
04-a 4 & ...nts at a distance and is therefore less  & \textbf{engaged}. \\
06-b 2 & ...es not depend on the attitudes of those  & \textbf{engaged} in them (Definition 1 and Assumption 2 ... \\
06-b 2 & ...es not depend on the attitudes of those  & \textbf{engaged} in it, it does not depend on the attitu... \\
07-a 4 & ...echnology that produced the feelings of  & \textbf{engaged} satisfaction without any engagement wou... \\
07-b 1 & ...ing, attending to one's happiness while  & \textbf{engaged} in an activity disrupts the absorption ... \\
10-a 2 & ...ngage her, together with a desire to be  & \textbf{engaged} otherwise. \\
10-a 3 & ...is aversive and includes a desire to be  & \textbf{engaged} otherwise (Definition 2.1). \\
10-a 5 & ...lls do better when their capacities are  & \textbf{engaged}, the disposition to boredom has the ada... \\
10-b 2 & ...edback that attention is not adequately  & \textbf{engaged}. \\
11-a 2 & ...the extent that its subject is actively  & \textbf{engaged} with what is worthwhile. \\
11-a 3 & ...worthwhile, and the subject need not be  & \textbf{engaged} with anything worthwhile in living it. \\
11-a 3 & A subject who is actively  & \textbf{engaged} with such activities satisfies Definiti... \\
11-a 5 & The  & \textbf{engaged} and worthwhile life of a person without... \\
17-a 1 & ...ctiveness: the subject must be actively  & \textbf{engaged}, and the object of her engagement must ... \\
17-a 2 & An agent is actively  & \textbf{engaged} with an object if she acts in ways dire... \\
17-a 2 & ...the extent that its subject is actively  & \textbf{engaged} with what is worthwhile. \\
17-a 2 & ...es, lived in the machine, is not a life  & \textbf{engaged} with the sick. \\
17-a 3 & Each is actively  & \textbf{engaged} in the sense of Definition 2.1. \\
17-a 3 & Whether she is  & \textbf{engaged} with something worthwhile depends far l... \\
17-a 4 & ...ment is unnecessary, Taylor's Sisyphus,  & \textbf{engaged} with his stone and achieving nothing, l... \\
19-b 1 & ... meaningful, and does not make anything  & \textbf{engaged} with more worthwhile. \\
20-a Abs & ...ions: whether the producing process was  & \textbf{engaged} with its work, in the sense the account... \\
20-a 1 & ...the extent that its subject is actively  & \textbf{engaged} in projects that are worthwhile indepen... \\
20-a 1 & ... passion to counting blades of grass is  & \textbf{engaged} with something worthless; the person wh... \\
20-a 1 & ...ut interest at a worthwhile task is not  & \textbf{engaged}. \\
20-a 4 & ...ion is satisfied, since its members are  & \textbf{engaged} in their work, and the question reduces... \\
20-b Abs & Whether the producing process was  & \textbf{engaged} with its work depends, by EXI-20-a, on ... \\
20-b Abs & ...ent by the producing system that it was  & \textbf{engaged} with its work would be an output of the... \\
20-b Abs & A statement that it was not  & \textbf{engaged} would be an output of the same kind. \\
\end{longtable}
\kwichead{engagement}{148}
\begin{longtable}{@{}>{\raggedright}p{0.07\textwidth}>{\raggedleft}p{0.41\textwidth}@{\,}>{\raggedright\arraybackslash}p{0.47\textwidth}@{}}
01-c Abs & ...n not from projects but from objects of  & \textbf{engagement}, such as persons and the world as an ob... \\
01-c Abs & ...of every agent, and show that sustained  & \textbf{engagement} with one of them constitutes a project ... \\
01-c 1 & ...o the degree that it consists in active  & \textbf{engagement} with projects of objective worth (Defin... \\
01-c 2 & An object of  & \textbf{engagement} is anything at which an agent can direc... \\
01-c 2 & ...y, with care, at a worthwhile object of  & \textbf{engagement}: another person, or some part of the wo... \\
01-c 2 & ...ity with care at a worthwhile object of  & \textbf{engagement}, and sustains it, her engagement is a p... \\
01-c 2 & ...ect of engagement, and sustains it, her  & \textbf{engagement} is a project of objective worth in the ... \\
01-c 2 & ...ity with care at a worthwhile object of  & \textbf{engagement}. \\
01-c 2 & ... Lemma 2.3, sustaining that activity is  & \textbf{engagement} in a project of objective worth. \\
01-c 2 & A life that includes such  & \textbf{engagement} consists to a positive extent in active... \\
01-c 2 & ...consists to a positive extent in active  & \textbf{engagement} with a project of objective worth, and ... \\
01-c 3 & ...ement: a life can be meaningful through  & \textbf{engagement} without success (Theorem 3.3 of EXI-17-... \\
01-c 3 & ... the availability of a project, because  & \textbf{engagement} makes one. \\
01-c 5 &  & \textbf{Engagement} with them is a project of objective wor... \\
01-d 1 & ...ifferent route: meaning requires active  & \textbf{engagement} with something worthwhile and not achie... \\
01-d 1 & ... in the ordinary sense, since sustained  & \textbf{engagement} with a worthwhile object constitutes on... \\
01-d 2 & Because meaning requires  & \textbf{engagement}, which is within the agent's power, and... \\
01-d 2 & ...ition 1 of EXI-01-a measures meaning by  & \textbf{engagement} and by the objective worth of what one ... \\
01-d 2 & ...or are projects of objective worth, and  & \textbf{engagement} with them confers meaning on lives that... \\
01-d 2 & ...dinary sense have worthwhile objects of  & \textbf{engagement} within reach. \\
01-d 2 & ...mine how much of a life can be given to  & \textbf{engagement} and how wide the range of worthwhile ob... \\
01-d 3 & ...ee of meaning varies with the extent of  & \textbf{engagement} and the worth of its object, and lives ... \\
01-d 4 & ... insists that the objects of meaningful  & \textbf{engagement} are not restricted to high culture; the... \\
01-d 5 & ...conditions are within everyone's reach:  & \textbf{engagement}, which is up to the agent, and worth, w... \\
03-a 2 & ...f objective worth, and EXI-17 that such  & \textbf{engagement}, not achievement, is what meaning requi... \\
03-a 2 & A character supports  & \textbf{engagement} with a project type only by having inte... \\
03-a 4 & ...ut Assumption 2.4 concerns the person's  & \textbf{engagement}. \\
04-a Abs & ...esponse, and it is compatible with full  & \textbf{engagement} in worthwhile projects. \\
04-a 4 & Meaning requires active  & \textbf{engagement} with something worthwhile (EXI-17). \\
04-a 4 & ... not care about anything, would corrode  & \textbf{engagement}; it is a different attitude and one we ... \\
04-b 2 & ...e attitude EXI-04-a calls irony: a full  & \textbf{engagement} with one's ends combined with a clear p... \\
04-b 2 & ...repancy between the seriousness of that  & \textbf{engagement} and the support available for it from t... \\
04-b 2 & ...'s conclusion is the claim that serious  & \textbf{engagement} with one's ends is compatible with luci... \\
04-b 3 & ...ver happy he is, since meaning requires  & \textbf{engagement} with something worthwhile (EXI-17-a, Th... \\
04-b 3 & ...ur ends include worthwhile objects, and  & \textbf{engagement} with them confers meaning (EXI-01). \\
04-b 4 & ...n the sense of EXI-04-a does not reduce  & \textbf{engagement}: it is compatible with full seriousness... \\
04-b 4 & ...he passion Camus admires belongs to the  & \textbf{engagement} and survives the change of attitude. \\
04-b 5 & What Camus rightly saw in him, unbroken  & \textbf{engagement} without consolation, is irony, the resp... \\
06-b Abs & ..., on the account of EXI-01, consists in  & \textbf{engagement} with what has worth independently of an... \\
06-b 2 & ...of EXI-17-a, which agrees, an object of  & \textbf{engagement} is worthwhile if its value does not dep... \\
06-b 3 & ...lar agent, then the worth of objects of  & \textbf{engagement} is precisely worth from that standpoint... \\
06-b 3 & ... is conditional: if meaning consists in  & \textbf{engagement} with what has objective worth, as EXI-0... \\
07-a Abs & ... account, consists to a large extent in  & \textbf{engagement} with and success in activities whose ob... \\
07-a 2 & ...n, or depends causally on, wholehearted  & \textbf{engagement} in activities and relationships: absorp... \\
07-a 2 & The forms of  & \textbf{engagement} in Principle P1 require that the agent ... \\
07-a 3 & Happiness is a by-product of  & \textbf{engagement} in activities and relationships valued ... \\
07-a 3 & ...s is realized to a large extent through  & \textbf{engagement}. \\
07-a 3 & By Principle P2,  & \textbf{engagement} contributes to happiness only when the ... \\
07-a 3 & ...large part of happiness that depends on  & \textbf{engagement}, and attains less happiness than she wo... \\
07-a 3 & ...t of happiness that is realized through  & \textbf{engagement} is available only to agents who do. \\
07-a 3 & ...-17, on which meaning comes from active  & \textbf{engagement} with what is worthwhile. \\
07-a 4 & ...has not acquired the attitudes on which  & \textbf{engagement} depends, and Theorem 3.3 applies to her... \\
07-a 4 & Perhaps  & \textbf{engagement} is merely a reliable cause of happiness... \\
07-a 4 & ...ngs of engaged satisfaction without any  & \textbf{engagement} would then give the direct aim its targ... \\
07-a 4 & ...he one who lives it than a life of real  & \textbf{engagement} with the same experiences (ETH-05). \\
07-a 6 & Much of it consists in, or depends on,  & \textbf{engagement} with activities and people valued for t... \\
07-b 1 & ...cks the object-directed caring on which  & \textbf{engagement} depends (Principle P2 of EXI-07-a). \\
07-b 2 & ...e forgoes the happiness that depends on  & \textbf{engagement}. \\
07-b 2 & ... she lacks the caring that the forms of  & \textbf{engagement} in Principle P1 require, and Theorem 3.... \\
07-b 2 & ... and P2 of EXI-07-a, which concern what  & \textbf{engagement} is: one cannot be absorbed in a game wh... \\
10-a 3 & A search that terminates when  & \textbf{engagement} is restored selects, among available al... \\
10-a 4 & Boredom motivates a search for  & \textbf{engagement}; it does not select which alternatives ... \\
10-a 4 & ...XI-17, on which meaning requires active  & \textbf{engagement} with something worthwhile, and boredom ... \\
10-a 4 & ... fallible, indication of the absence of  & \textbf{engagement}. \\
10-a 5 & ...adaptive function of moving them toward  & \textbf{engagement}. \\
10-b 1 & ... restore the match, so that, given that  & \textbf{engagement} pays (its Assumption 2.4), the disposit... \\
10-b 3 & Fit and  & \textbf{engagement}.. \\
10-b 3 & ...ssumption 2.4 of EXI-10-a, the value of  & \textbf{engagement}. \\
10-b 4 & ...applies: boredom motivates a search for  & \textbf{engagement} but does not select the options on offe... \\
10-b 4 & The evidence that  & \textbf{engagement} yields concentration, involvement and b... \\
10-b 5 & ...t boredom is feedback about attentional  & \textbf{engagement}. \\
10-b 7 & ...a decline in boredom and an increase in  & \textbf{engagement} or learning. \\
10-b 8 & ... that a meaningful life requires active  & \textbf{engagement} with something worthwhile; the meaning ... \\
10-b 8 & ... the mind registers the absence of such  & \textbf{engagement}. \\
10-b 11 & ...auses and in its converse, that matched  & \textbf{engagement} absorbs, and leaves its over-challenge ... \\
10-b 11 & ...ese findings together with the value of  & \textbf{engagement}, and the studies of costly responses an... \\
11-a Abs & ...is good offers objects and relations of  & \textbf{engagement} that a naturalistic world lacks. \\
11-a 1 & ...eaning arises when the subject's active  & \textbf{engagement} meets the objective worth of its object... \\
11-a 2 & Other things equal,  & \textbf{engagement} with what is of greater worth, or engag... \\
11-a 2 & ...ement with what is of greater worth, or  & \textbf{engagement} that relates one to more of what is wor... \\
11-a 3 & By Theorem 3.3, such a person's  & \textbf{engagement} with what is worthwhile makes her life ... \\
11-a 4 &  & \textbf{Engagement} with that being, through love, worship,... \\
11-a 4 & ... the attempt to understand it, would be  & \textbf{engagement} with the greatest object of worth there... \\
11-a 4 & ... would make a life more meaningful than  & \textbf{engagement} with lesser objects, other things equal... \\
11-a 5 & ...he pursuit of meaning is not a solitary  & \textbf{engagement} with isolated goods. \\
11-a 5 & ...e uses is no longer purpose as such but  & \textbf{engagement} with the good, which is Definition 2.1. \\
11-a 6 & ...out God is meaningful: meaning requires  & \textbf{engagement} with what is worthwhile, and lives with... \\
11-a 7 & Meaning consists in active  & \textbf{engagement} with what is worthwhile, and the worth ... \\
11-a 7 & ...hwhile, and the worth of the objects of  & \textbf{engagement} does not derive from their being assign... \\
11-b Abs & ..., not the objects of worth available to  & \textbf{engagement}, so that Proposition 4.1 survives both ... \\
11-b 1 & EXI-11-a defines meaning as active  & \textbf{engagement} with what is worthwhile (Definition 2.1... \\
11-b 1 & ...ne engages and with the extent to which  & \textbf{engagement} relates one to more of what is worthwhi... \\
11-b 1 & ... and third are objects and relations of  & \textbf{engagement} of great worth that a naturalistic worl... \\
11-b 1 & ...inition 2.1 of EXI-11-a, any meaning is  & \textbf{engagement} with what is worthwhile. \\
11-b 1 &  & \textbf{Engagement} with God, or with a purpose that unifie... \\
11-b 1 & ...purpose that unifies one's projects, is  & \textbf{engagement} with what is worthwhile; so the meaning... \\
11-b 2 & ...d by God's existence, so the objects of  & \textbf{engagement} available without God remain available ... \\
11-b 3 & ...life by adding objects and relations of  & \textbf{engagement} of great worth. \\
14-a Abs & ...of proportion to the good or crowds out  & \textbf{engagement} with present goods. \\
17-a Abs & A meaningful life requires active  & \textbf{engagement} with something worthwhile, in the sense... \\
17-a Abs & ...rns on pairs of lives that are alike in  & \textbf{engagement} and in the worth of their objects and d... \\
17-a 1 & ...actively engaged, and the object of her  & \textbf{engagement} must be worthwhile independently of her... \\
17-a 1 & ...must be worthwhile independently of her  & \textbf{engagement}. \\
17-a 1 & ... it should not: meaning requires active  & \textbf{engagement} with something worthwhile, and achievem... \\
17-a 2 & An object of  & \textbf{engagement} (a project, a person, a practice, an in... \\
17-a 2 & Active  & \textbf{engagement} in the sense relevant to Principle P1 r... \\
17-a 2 & ... a real object and actually bear on it;  & \textbf{engagement} with an object that exists only in the ... \\
17-a 3 & Each object of  & \textbf{engagement} is worthwhile in the sense of Definitio... \\
17-a 3 & ...life is meaningful to the extent of her  & \textbf{engagement} with music. \\
17-a 3 & Her  & \textbf{engagement} is active, since study and attentive li... \\
17-a 3 & By Principle P1, her  & \textbf{engagement} contributes meaning. \\
17-a 3 & A meaningful life requires active  & \textbf{engagement} with something worthwhile; it does not ... \\
17-a 3 & ...lternatives: subjectivism fails because  & \textbf{engagement} with worthless objects confers no meani... \\
17-a 3 & ...ls because worthwhile work done without  & \textbf{engagement}, out of mere duty or inertia, is not wh... \\
17-a 3 & ...thless pursuit, or one produced without  & \textbf{engagement}, confers none. \\
17-a 3 & ...ndifferent to it fails the condition of  & \textbf{engagement}. \\
17-a 3 & ...larges the worth realized through one's  & \textbf{engagement}, it may enlarge meaning too. \\
17-a 4 & ... difficult activity counts, then active  & \textbf{engagement} in a demanding worthwhile activity enta... \\
17-a 4 & ...chievement becomes a restatement of the  & \textbf{engagement} condition. \\
17-a 4 & Not all  & \textbf{engagement} is difficult, moreover: love of a frien... \\
17-a 5 & A meaningful life requires active  & \textbf{engagement} with something worthwhile. \\
17-a 5 & It does not require that the  & \textbf{engagement} succeed. \\
17-a 5 & ...at the experience machine shows is that  & \textbf{engagement} must reach a real object, not that it m... \\
19-b 1 & ...or as a condition of the persistence of  & \textbf{engagement}. \\
19-b 1 & ...limited by death is not something whose  & \textbf{engagement} makes the life meaningful, and does not... \\
19-b 1 & The worth of the objects of  & \textbf{engagement}, knowledge, persons, beauty and the goo... \\
19-b 1 & ...ut categorical desires lacks the active  & \textbf{engagement} that meaning requires (EXI-17). \\
19-b 1 & ...1.3 concerns the eventual exhaustion of  & \textbf{engagement}, not its quality at any given time. \\
19-b 1 & An immortal's  & \textbf{engagement} with inquiry in the first thousand year... \\
19-b 1 & ...ry in the first thousand years would be  & \textbf{engagement} with what is worthwhile, as meaningful ... \\
19-b 2 & ...shows that the limit is something whose  & \textbf{engagement} confers meaning. \\
19-b 2 & ...-life meaning and of the persistence of  & \textbf{engagement}, then death benefits the one who dies, ... \\
19-b 2 & ... ordinary case, have left more time for  & \textbf{engagement}, as Corollary 3.3 of EXI-19-a notes (Fi... \\
19-b 3 & ...keeps the life within the span in which  & \textbf{engagement} persists. \\
20-a Abs & ...ul activity within one, requires active  & \textbf{engagement} with something worthwhile. \\
20-a Abs & ...rocess has the kind of mental life that  & \textbf{engagement} requires, which MND-08 and MND-05 leave... \\
20-a 1 & ...hat are worthwhile independently of her  & \textbf{engagement}. \\
20-a 1 & ...n of this view: meaning requires active  & \textbf{engagement} with something worthwhile, and does not... \\
20-a 2 &  & \textbf{Engagement}, in Wolf's sense, is not mere activity. \\
20-a 2 & ...are states with a phenomenal character,  & \textbf{engagement} requires consciousness. \\
20-a 2 & ...he evidence available, the condition of  & \textbf{engagement} is neither known to be satisfied nor kn... \\
20-a 2 & If  & \textbf{engagement} can be understood functionally, as a pa... \\
20-a 2 & Whether this is  & \textbf{engagement} or only its functional profile is the q... \\
20-b Abs & ...f EXI-17, on whether its production was  & \textbf{engagement} with something worthwhile, and that nei... \\
20-b Abs & Evidence bearing on  & \textbf{engagement}.. \\
20-b Abs & ... on whether it has the mental life that  & \textbf{engagement} requires, or on whether engagement can ... \\
20-b Abs & ...that engagement requires, or on whether  & \textbf{engagement} can be understood functionally. \\
20-b Abs & Whether the production was  & \textbf{engagement} may be settled by empirical work on the... \\
\end{longtable}
\kwichead{everyone}{11}
\begin{longtable}{@{}>{\raggedright}p{0.07\textwidth}>{\raggedleft}p{0.41\textwidth}@{\,}>{\raggedright\arraybackslash}p{0.47\textwidth}@{}}
01-d Abs & ...t meaning in this sense is available to  & \textbf{everyone}. \\
01-d 4 & ...el meet this condition, and are open to  & \textbf{everyone}. \\
01-d 4 & If meaning is available to  & \textbf{everyone}, the claim that a life is meaningful is... \\
01-d 4 & That positive meaning is available to  & \textbf{everyone} is compatible with its being hard to at... \\
01-d 5 & ... on the account of EXI-01-a, is open to  & \textbf{everyone} because its conditions are within every... \\
02-a 5 & On the deprivation account,  & \textbf{everyone} who dies at any age is deprived of some... \\
02-a 6 & ...ppen to a person, or that it is bad for  & \textbf{everyone} who dies. \\
07-a 1 & Aristotle observed that  & \textbf{everyone} agrees on the name of the highest good,... \\
11-a 2 & ...t meaning so understood is available to  & \textbf{everyone}. \\
17-a 3 & ...t meaning in this sense is available to  & \textbf{everyone}. \\
20-a 1 & ...t meaning in this sense is available to  & \textbf{everyone}. \\
\end{longtable}
\kwichead{exist}{29}
\begin{longtable}{@{}>{\raggedright}p{0.07\textwidth}>{\raggedleft}p{0.41\textwidth}@{\,}>{\raggedright\arraybackslash}p{0.47\textwidth}@{}}
02-a Abs & ... the event is bad, and that the subject  & \textbf{exist} at the time of the loss. \\
02-a 1 & ...is the privation of sensation; while we  & \textbf{exist} death is not present, and when death is... \\
02-a 1 & ... and when death is present we no longer  & \textbf{exist} (EpicurusMen ). \\
02-a 1 & ... birth, an eternity in which we did not  & \textbf{exist}, is nothing to us, and the time after o... \\
02-a 2 & ... and the relation does not require that  & \textbf{exist} at a time at which its other relata are... \\
02-a 3 & ...f Principle P1 holds, although does not  & \textbf{exist} at any time after at which the goods of... \\
02-a 3 & ... located at times at which she does not  & \textbf{exist}. \\
02-a 3 & Nothing in Principle P1 requires that  & \textbf{exist} at the times at which the goods she is ... \\
02-a 3 & ..., read as the claim that a subject must  & \textbf{exist} at the time at which an event is bad fo... \\
02-a 3 & ... the strong existence requirement: that  & \textbf{exist} at the time of the loss. \\
02-a 3 & ...the badness of death requires only that  & \textbf{exist}, at some time, and not that she exist a... \\
02-a 3 & ...t exist, at some time, and not that she  & \textbf{exist} at the times when the lost goods would ... \\
02-a 4 & For a person who dies at and began to  & \textbf{exist} at , the closest world in which dies la... \\
02-a 4 & ...an the closest world in which begins to  & \textbf{exist} substantially earlier than , and the la... \\
02-a 4 & A world in which begins to  & \textbf{exist} decades earlier must contain 's parents... \\
02-a 5 & ...t happened, or after, when she does not  & \textbf{exist}. \\
02-a 7 & ...ls because it requires that the subject  & \textbf{exist} at the time of the loss, and deprivatio... \\
02-b 2 & ...ad instead as a demand that the subject  & \textbf{exist} at the time at which the lost goods wou... \\
04-a 2 & ...player who knows that the game's stakes  & \textbf{exist} only within the game and still plays to... \\
09-a 3 & ... is fixed when the human being comes to  & \textbf{exist}; Lemma 3.2 concerns defining projects, ... \\
11-b 2 & ...rson may reasonably prefer that God not  & \textbf{exist}, since God's existence would deprive he... \\
12-a 2 & ...lternative is one in which she does not  & \textbf{exist}. \\
12-b 2 & ...py person who is never created does not  & \textbf{exist}, so no one has a complaint, and the imp... \\
13-a 3 & That she does not  & \textbf{exist} when the duty is to be discharged is no... \\
13-a 4 & The dead do not  & \textbf{exist}, so nothing can be owed to them, and th... \\
13-a 4 & ...t in the regard of people who no longer  & \textbf{exist}. \\
14-a 3 & ...since it constructs a past that did not  & \textbf{exist} and seeks to return to it; her reflecti... \\
14-a 4 & ...s, elsewhere in time, as distant places  & \textbf{exist} elsewhere in space. \\
19-a 3 & It begins with the person's coming to  & \textbf{exist}. \\
\end{longtable}
\kwichead{extent}{29}
\begin{longtable}{@{}>{\raggedright}p{0.07\textwidth}>{\raggedleft}p{0.41\textwidth}@{\,}>{\raggedright\arraybackslash}p{0.47\textwidth}@{}}
01-c 2 & ... such engagement consists to a positive  & \textbf{extent} in active engagement with a project of ... \\
01-d Abs & ...-a holds that life is meaningful to the  & \textbf{extent} that one engages with projects of objec... \\
01-d 3 & ...: the degree of meaning varies with the  & \textbf{extent} of engagement and the worth of its obje... \\
02-a Abs & ...t, an event is bad for a subject to the  & \textbf{extent} that her life as a whole goes worse, by... \\
03-a 2 & ...argues that a life is meaningful to the  & \textbf{extent} that its subject engages with projects ... \\
04-a 2 & ...de toward a situation is fitting to the  & \textbf{extent} that the evaluative features under whic... \\
04-a 3 & ...1 argues that life is meaningful to the  & \textbf{extent} that one engages with projects of objec... \\
06-b 2 & ... of EXI-01, a life is meaningful to the  & \textbf{extent} that its subject engages with projects ... \\
07-a Abs & ... plausible account, consists to a large  & \textbf{extent} in engagement with and success in activ... \\
07-a 1 & ...rve, and a life is well directed to the  & \textbf{extent} that it is directed at it. \\
07-a 2 & ...s of pursuits if is realized to a large  & \textbf{extent} through those pursuits, and the pursuit... \\
07-a 3 & ...le P1, happiness is realized to a large  & \textbf{extent} through engagement. \\
09-a 5 & ...1, on which a life is meaningful to the  & \textbf{extent} that its projects are of objective wort... \\
11-a Abs & ...of meaning: a life is meaningful to the  & \textbf{extent} that its subject actively engages with ... \\
11-a 2 & A life is meaningful to the  & \textbf{extent} that its subject is actively engaged wi... \\
11-a 2 & ...on 2.1, which makes meaning a matter of  & \textbf{extent}. \\
11-a 3 & ...er life is therefore meaningful to that  & \textbf{extent}. \\
11-b 1 & ... worth of what one engages and with the  & \textbf{extent} to which engagement relates one to more... \\
16-b 1 & ... and grief would be fitting only to the  & \textbf{extent} that substitutes are scarce. \\
17-a 1 & ...fe that brings nothing about is to that  & \textbf{extent} empty. \\
17-a 2 & A life is meaningful to the  & \textbf{extent} that its subject is actively engaged wi... \\
17-a 3 & ...he listener's life is meaningful to the  & \textbf{extent} of her engagement with music. \\
18-a 2 & To the  & \textbf{extent} that these attitudes serve the function... \\
18-b 1 & ...ent, which do not misrepresent; to that  & \textbf{extent} the range of necessary positive illusio... \\
19-b 1 & ...and EXI-17, a life is meaningful to the  & \textbf{extent} that its subject actively engages with ... \\
20-a 1 & ...activeness: a life is meaningful to the  & \textbf{extent} that its subject is actively engaged in... \\
20-a 3 & ...lection of answers is worthwhile to the  & \textbf{extent} that it contributes to knowledge or und... \\
20-a 4 & ...tivity is meaningful, by EXI-17, to the  & \textbf{extent} that the collection is worthwhile, and ... \\
20-b Abs & ...October 2026 are an instance, but their  & \textbf{extent} and accuracy are facts to be establishe... \\
\end{longtable}
\kwichead{given}{25}
\begin{longtable}{@{}>{\raggedright}p{0.07\textwidth}>{\raggedleft}p{0.41\textwidth}@{\,}>{\raggedright\arraybackslash}p{0.47\textwidth}@{}}
01-c 4 & It entails Assumption 3  & \textbf{given} Lemma 2.3, as it must if the conclusion... \\
01-d 1 & We take these results as  & \textbf{given} and ask what follows. \\
01-d 2 & ...dom determine how much of a life can be  & \textbf{given} to engagement and how wide the range of... \\
03-a Abs & ...re projects she takes to be worthwhile;  & \textbf{given} a character that remains connected to h... \\
03-a 2 & ...he question is how many kinds of them a  & \textbf{given} character can care about. \\
03-b 1 &  & \textbf{Given} a bounded repertoire of project types f... \\
05-a 2 & ...tion is rational only if it is fitting,  & \textbf{given} what the subject is in a position to kn... \\
06-a Abs & ... not aimed at us, and the goods are not  & \textbf{given} to us. \\
07-b Abs & We argue that, in the form  & \textbf{given} by EXI-07-a, it is neither. \\
09-a 3 & ...nd his claim is that no such project is  & \textbf{given} to human beings in advance of their exi... \\
10-b 1 & ...es tends to restore the match, so that,  & \textbf{given} that engagement pays (its Assumption 2.... \\
10-b 10 & ...e over-challenge clause, for the reason  & \textbf{given} in Section 3. \\
11-a 3 & ...se of this kind degrades the life it is  & \textbf{given} to. \\
11-b 2 & That a  & \textbf{given} person might reasonably prefer the worl... \\
12-a Abs & ... Schopenhauer says yes, and Benatar has  & \textbf{given} the claim a rigorous form. \\
12-a Abs & ...e person who lives it, and to have been  & \textbf{given} it is a non-comparative benefit. \\
15-b 3 & ...evidence supports is the probability of  & \textbf{given} her commitment. \\
15-b 4 & Expected-value reasoning, applied to a  & \textbf{given} set of options, does not determine thes... \\
17-a 1 & Bradford (2015) has  & \textbf{given} an account of achievement as the compet... \\
18-b 1 & (i) Whether a  & \textbf{given} positive illusion is psychologically ne... \\
18-b 1 & ...sion is psychologically necessary for a  & \textbf{given} agent is an empirical question, and the... \\
19-a Abs & ...es, since a later death would also have  & \textbf{given} her life a shape, and in the ordinary c... \\
19-a 1 & Religious traditions have  & \textbf{given} death a meaning of its own, as the wage... \\
19-b 1 & ...n of engagement, not its quality at any  & \textbf{given} time. \\
19-b 2 & ...ular one; a later death would also have  & \textbf{given} the life a shape and would, in the ordi... \\
\end{longtable}
\kwichead{includes}{15}
\begin{longtable}{@{}>{\raggedright}p{0.07\textwidth}>{\raggedleft}p{0.41\textwidth}@{\,}>{\raggedright\arraybackslash}p{0.47\textwidth}@{}}
01-c 2 & A life that  & \textbf{includes} such engagement consists to a positive ... \\
04-a 4 & ...titude toward one's seriousness, and it  & \textbf{includes} continuing to hold one's ends. \\
07-a 3 & ...TH-05's argument that well-being itself  & \textbf{includes} contact with reality. \\
07-a 4 & ...nse but the good life as a whole, which  & \textbf{includes} the goods that a plural account counts ... \\
09-b 1 & ...s the paper-knife from a concept, which  & \textbf{includes} both what the knife is for and a techni... \\
09-b 2 & ...what it is for, and the real definition  & \textbf{includes} the defining project. \\
10-a 3 & Boredom is aversive and  & \textbf{includes} a desire to be engaged otherwise (Defin... \\
10-a 3 & ...ive to meaning: the mismatch it signals  & \textbf{includes} the idleness of one's capacity to care,... \\
10-b 5 & ...on 2018), and EXI-10-a's Definition 2.2  & \textbf{includes} both as forms of a single mismatch. \\
10-b 9 & ... the direction of the finding; the list  & \textbf{includes} every study we could verify on the cost... \\
15-b 3 & ... more likely; the agent's evidence then  & \textbf{includes} the fact that she is committed to the e... \\
15-b 4 & But the hopeful stance  & \textbf{includes} more than choices among options. \\
15-b 4 & It  & \textbf{includes} where attention goes, which possibiliti... \\
17-a 2 & On Definition 2.2, the worthwhile  & \textbf{includes} the components of the good identified i... \\
18-a 2 & ...s a cost even when it helps: well-being  & \textbf{includes} contact with reality (ETH-05), and each... \\
\end{longtable}
\kwichead{intentional}{18}
\begin{longtable}{@{}>{\raggedright}p{0.07\textwidth}>{\raggedleft}p{0.41\textwidth}@{\,}>{\raggedright\arraybackslash}p{0.47\textwidth}@{}}
01-c Abs & ... is available to every agent capable of  & \textbf{intentional} action. \\
01-c 1 & ...and assumes that every agent capable of  & \textbf{intentional} action is able to engage in at least on... \\
01-c 2 & ...s anything at which an agent can direct  & \textbf{intentional} activity and about which she can care: ... \\
01-c 2 & Every agent capable of  & \textbf{intentional} action can direct some intentional acti... \\
01-c 2 & ...e of intentional action can direct some  & \textbf{intentional} activity, with care, at a worthwhile ob... \\
01-c 2 & If an agent directs  & \textbf{intentional} activity with care at a worthwhile obje... \\
01-c 2 & For every agent capable of  & \textbf{intentional} action, a life of positive meaning is a... \\
01-c 2 & By Assumption 2.2, the agent can direct  & \textbf{intentional} activity with care at a worthwhile obje... \\
01-c 4 & Some agents capable of  & \textbf{intentional} action are wholly isolated and wholly u... \\
01-c 4 & An agent capable of  & \textbf{intentional} action acts on something, and anything ... \\
01-d 1 & ... states that for every agent capable of  & \textbf{intentional} action a life of positive meaning is av... \\
01-d 2 & No agent capable of  & \textbf{intentional} action is condemned by her circumstance... \\
04-a 3 & Any agent capable of  & \textbf{intentional} action takes some ends seriously. \\
04-a 3 & ...life of any reflective agent capable of  & \textbf{intentional} action is absurd in the sense of Defini... \\
09-a 2 & Principle P1 is the standard  & \textbf{intentional} account of artefact kinds, and MET-13 a... \\
09-a 4 & ...as an artefact of its kind, through the  & \textbf{intentional} act of selecting and placing the wood, ... \\
10-a 4 & ...ferent object: PHE-04 argues that it is  & \textbf{intentional} toward time itself, not toward an activ... \\
10-b 8 & ...found boredom, in Heidegger's sense, is  & \textbf{intentional} toward time itself; the empirical liter... \\
\end{longtable}
\kwichead{least}{14}
\begin{longtable}{@{}>{\raggedright}p{0.07\textwidth}>{\raggedleft}p{0.41\textwidth}@{\,}>{\raggedright\arraybackslash}p{0.47\textwidth}@{}}
01-c Abs & ...very such agent is able to engage in at  & \textbf{least} one project of objective worth (Assumpt... \\
01-c 1 & ...entional action is able to engage in at  & \textbf{least} one of them (Assumption 3). \\
02-b 1 & ... must be experienced by her, or must at  & \textbf{least} occur while she exists. \\
09-a 1 & ...existentialism denies this: there is at  & \textbf{least} one being whose existence comes before ... \\
10-b 4 & The adaptive component is the  & \textbf{least} directly tested. \\
15-a 1 & ...resolve to act as if will obtain, or at  & \textbf{least} as if it might, which he called substan... \\
16-a 2 & ...that is a reason for 's love of , to at  & \textbf{least} the same degree. 's love of is transfer... \\
16-a 2 & ... the reasons for 's love of would be at  & \textbf{least} as strong reasons for 's loving instead... \\
16-a 3 & ...that is a reason for 's love of , to at  & \textbf{least} the same degree. \\
16-a 3 & ...nciple P2, each of these features is at  & \textbf{least} as strong a reason for 's loving . \\
16-a 3 & ... reasons, the reasons for loving are at  & \textbf{least} as strong reasons for loving instead, a... \\
16-b 1 & ... love that person instead of that is at  & \textbf{least} as strong as the reasons for 's love of... \\
16-b 1 & ...refore do not favor the other person at  & \textbf{least} as strongly. \\
17-a 3 & ...gful than the unsuccessful one's, or at  & \textbf{least} better. \\
\end{longtable}
\kwichead{life}{450}
\begin{longtable}{@{}>{\raggedright}p{0.07\textwidth}>{\raggedleft}p{0.41\textwidth}@{\,}>{\raggedright\arraybackslash}p{0.47\textwidth}@{}}
01-c Abs & EXI-01-a establishes that a  & \textbf{life} of positive meaning is available to eve... \\
01-c 1 & EXI-01-a defines a  & \textbf{life} as meaningful to the degree that it con... \\
01-c 2 & ... agent capable of intentional action, a  & \textbf{life} of positive meaning is available. \\
01-c 2 & A  & \textbf{life} that includes such engagement consists ... \\
01-c 3 & ...es a parallel result for achievement: a  & \textbf{life} can be meaningful through engagement wi... \\
01-c 5 & ...stributed: the presence, in every human  & \textbf{life}, of worthwhile things to care about and... \\
01-c 5 & ... is a project of objective worth, and a  & \textbf{life} that contains it has meaning. \\
01-d Abs & EXI-01-a holds that  & \textbf{life} is meaningful to the extent that one en... \\
01-d Abs & ... scale of one's projects; that no one's  & \textbf{life} is condemned to meaninglessness by circ... \\
01-d Abs & It does not imply that every  & \textbf{life} is meaningful, that all lives are equal... \\
01-d 1 & ...y agent capable of intentional action a  & \textbf{life} of positive meaning is available. \\
01-d 1 & ... does not settle whether any particular  & \textbf{life} is meaningful, which depends on what is... \\
01-d 2 & Whether a  & \textbf{life} is meaningful does not depend on whethe... \\
01-d 2 & A  & \textbf{life} can be meaningful though its subject ha... \\
01-d 2 & ...er circumstances alone to a meaningless  & \textbf{life}. \\
01-d 2 & ...justice, affect the degree of meaning a  & \textbf{life} can attain and the ease with which it c... \\
01-d 2 & ...ion and freedom determine how much of a  & \textbf{life} can be given to engagement and how wide... \\
01-d 3 & ...irst, from availability to actuality: a  & \textbf{life} in which the agent engages with nothing... \\
01-d 4 & ...available to everyone, the claim that a  & \textbf{life} is meaningful is cheap and of no evalua... \\
01-d 4 & EXI-04 grants that  & \textbf{life} is absurd in Nagel's sense and argues t... \\
01-d 5 & ...its degree is not, and how much of it a  & \textbf{life} attains depends on what is done with it... \\
02-a Abs & ...ad for a subject to the extent that her  & \textbf{life} as a whole goes worse, by the measure o... \\
02-a 1 & ...-dimensional conception of the person's  & \textbf{life}. \\
02-a 2 & ...se in which it sustains a psychological  & \textbf{life}, and persists through psychological con... \\
02-a 2 & ...ganism that is sustains a psychological  & \textbf{life}, and no psychological life continuous w... \\
02-a 2 & ...sychological life, and no psychological  & \textbf{life} continuous with 's is sustained at any ... \\
02-a 2 & For a possible  & \textbf{life} of , is the sum of the intrinsic values... \\
02-a 2 & ...for of an event is , where is 's actual  & \textbf{life} and is the life has in the closest poss... \\
02-a 2 & ...is , where is 's actual life and is the  & \textbf{life} has in the closest possible world in wh... \\
02-a 2 & ...difference it makes to the value of the  & \textbf{life} as a whole, and the lives compared are ... \\
02-a 2 & ...the value it assigns is a fact about 's  & \textbf{life} taken as a whole. \\
02-a 3 & ...crupulous never to let her know it; her  & \textbf{life} goes on, as experienced, as before. \\
02-a 3 & Her  & \textbf{life} contains friendships she believes genui... \\
02-a 3 & ....2, relationships are among the goods a  & \textbf{life} contains, and false friendships are not... \\
02-a 3 & So her  & \textbf{life} has less value than the life she would ... \\
02-a 3 & So her life has less value than the  & \textbf{life} she would have had without the betrayal... \\
02-a 3 & ... is a misfortune for him because of the  & \textbf{life} it has taken from him. \\
02-a 3 & It is a fact about and the value of her  & \textbf{life} as a whole, and it obtains, tenselessly... \\
02-a 3 & ...tenselessly, because exists and has the  & \textbf{life} she has. \\
02-a 3 & ...requires that those goods belong to the  & \textbf{life} she would have had. \\
02-a 3 & If dies at , and the  & \textbf{life} would have had had she not died at cont... \\
02-a 3 & ...reater balance of goods than her actual  & \textbf{life}, then her death is bad for her. \\
02-a 3 & ...e she would otherwise have gone on to a  & \textbf{life} worth living, this condition is met. \\
02-a 3 & Her actual  & \textbf{life} ends at ; the alternative life agrees w... \\
02-a 3 & ...r actual life ends at ; the alternative  & \textbf{life} agrees with it up to and continues beyo... \\
02-a 3 & ...in, and if those circumstances made her  & \textbf{life} worth living, the continuation has a po... \\
02-a 3 & ... twenty that cuts off decades of a good  & \textbf{life} is very bad for the one who dies; a dea... \\
02-a 3 & ...ters is not the manner of dying but the  & \textbf{life} that would otherwise have followed. \\
02-a 4 & ...he latter often contains no longer good  & \textbf{life} for . \\
02-a 4 & ...t is possible at all, it is a different  & \textbf{life} from the beginning, and nothing guarant... \\
02-a 4 & ...t is compared is the value of the whole  & \textbf{life}, not its length. \\
02-a 4 & ... world where she begins earlier and her  & \textbf{life} is better, and the deprivation account ... \\
02-a 5 & ...1 is a fact about the value of 's whole  & \textbf{life} and is not located at any time within i... \\
02-a 5 & ...riendships are no more genuine, and her  & \textbf{life} no better. \\
02-a 5 & It says that death is bad when the  & \textbf{life} it ends would otherwise have continued ... \\
02-a 5 & ...adness of death to diminish over a long  & \textbf{life}, as the goods that would have followed ... \\
02-a 5 & ...fewer; EXI-03 argues that in an endless  & \textbf{life} they eventually give out. \\
02-a 5 & ...nition 2.2 are goods of a psychological  & \textbf{life}, which the corpse does not have. \\
02-a 5 & ...she is deprived of is located after her  & \textbf{life} ends. \\
02-a 7 & ...badness is a comparative fact about her  & \textbf{life} as a whole, and its subject is the pers... \\
02-a 7 & ...ts subject is the person who lived that  & \textbf{life}. \\
02-b 2 & ...m is comparative and concerns the whole  & \textbf{life}: a life that ends earlier contains fewe... \\
02-b 2 & ...parative and concerns the whole life: a  & \textbf{life} that ends earlier contains fewer goods ... \\
03-a Abs & ...ms argued that it would not: an endless  & \textbf{life} would either preserve the character of ... \\
03-a Abs & ...e that character so far that the future  & \textbf{life} would not be hers in the sense that mat... \\
03-a Abs & ...ended pursuits could sustain an endless  & \textbf{life}, and that the second horn rests on a de... \\
03-a Abs & ...rs in survival, we show that an endless  & \textbf{life} either exhausts the categorical desires... \\
03-a Abs & ...r ceases to be, in its later reaches, a  & \textbf{life} that matters to the person now. \\
03-a 1 & ...e refuses to take it again, because her  & \textbf{life} has become one of boredom and coldness. \\
03-a 1 & ... used the case to argue that an endless  & \textbf{life} would be meaningless, and that it is no... \\
03-a 1 & Williams argued that an endless  & \textbf{life} would face a dilemma. \\
03-a 1 & ... would eventually be exhausted, and her  & \textbf{life} would become boring. \\
03-a 1 & ...he future person would not be one whose  & \textbf{life} she now has reason to want, since what ... \\
03-a 1 & ... food, music and friendship, and that a  & \textbf{life} could be sustained indefinitely by a su... \\
03-a 1 & ...ifies the categorical desires that give  & \textbf{life} its point with desires for projects the... \\
03-a 2 & ...its satisfaction requires her continued  & \textbf{life} and it gives her a reason to go on livi... \\
03-a 2 & A categorical desire gives  & \textbf{life} its point if its object is a project th... \\
03-a 2 & ...order to drink it, but it does not give  & \textbf{life} its point. \\
03-a 2 & EXI-01 argues that a  & \textbf{life} is meaningful to the extent that its su... \\
03-a 2 & The desires that give  & \textbf{life} its point are desires to engage in such... \\
03-a 2 & A  & \textbf{life} is endless if for every time at which t... \\
03-a 2 & ...grounds a categorical desire that gives  & \textbf{life} its point for that person. \\
03-a 2 & ...o a person, in her concern for a future  & \textbf{life}, is carried mostly by psychological con... \\
03-a 2 & ...e of connectedness below which a future  & \textbf{life} gives the person now little of what mat... \\
03-a 2 & ...already done is not a desire that gives  & \textbf{life} its point. \\
03-a 3 & In an endless  & \textbf{life} in which the person's character remains... \\
03-a 3 & ...he has no categorical desire that gives  & \textbf{life} its point. \\
03-a 3 & The  & \textbf{life} is endless (Definition 2.2), so there i... \\
03-a 3 & ...grounds a categorical desire that gives  & \textbf{life} its point. \\
03-a 3 & In an endless  & \textbf{life} in which the person's character does no... \\
03-a 3 & ...eshold, there are times after which the  & \textbf{life} gives her, as she now is, little of wha... \\
03-a 3 & ...t character, then by Assumption 2.5 the  & \textbf{life} after that time gives her now little of... \\
03-a 3 & So the later reaches of the  & \textbf{life} are not something she now has reason to... \\
03-a 3 & An endless  & \textbf{life} is not desirable for the one who would ... \\
03-a 3 & ...haust the categorical desires that give  & \textbf{life} its point, or in its later reaches ceas... \\
03-a 3 & ..., or in its later reaches cease to be a  & \textbf{life} that gives her what matters. \\
03-a 3 & Every endless  & \textbf{life} falls under Lemma 3.1 or Lemma 3.2, acc... \\
03-a 3 & ...d without categorical desires that give  & \textbf{life} its point, so the person's life becomes... \\
03-a 3 & ...at give life its point, so the person's  & \textbf{life} becomes, from some time on, pointless; ... \\
03-a 3 & In neither case does the endless  & \textbf{life} offer her what she wants when she wants... \\
03-a 3 & ...wants not to die: the continuation of a  & \textbf{life} with point, for herself as she is. \\
03-a 3 & Theorem 3.3 is a claim about an endless  & \textbf{life}, not about a longer one. \\
03-a 3 & ...l has categorical desires that give her  & \textbf{life} its point is deprived by her death; a p... \\
03-a 3 & ...deprivation of EXI-02 diminishes as the  & \textbf{life} approaches the point at which Lemma 3.1... \\
03-a 3 & ...which Lemma 3.1 applies, and an endless  & \textbf{life} is guaranteed to reach that point or to... \\
03-a 4 & A  & \textbf{life} could be sustained endlessly by these (... \\
03-a 4 & ...ble pleasures are goods, and an endless  & \textbf{life} containing them would contain goods at ... \\
03-a 4 & But the desire for them does not give  & \textbf{life} its point in the sense of Definition 2.... \\
03-a 4 & A  & \textbf{life} sustained only by such desires is the l... \\
03-a 4 & ...e sustained only by such desires is the  & \textbf{life} Elina Makropulos describes, in which th... \\
03-a 4 & Over an endless  & \textbf{life} the new theorems available to her becom... \\
03-a 4 & ...y constitution or by a suitably managed  & \textbf{life}. \\
03-a 4 & ...ereby have them again; she would have a  & \textbf{life} without point that does not trouble her... \\
03-a 4 & ...heorem 3.3 concerns whether the endless  & \textbf{life} is desirable, and a life whose pointles... \\
03-a 4 & ...er the endless life is desirable, and a  & \textbf{life} whose pointlessness is not felt is not ... \\
03-a 4 & ...ntlessness is not felt is not thereby a  & \textbf{life} with a point. \\
03-a 4 & Such a  & \textbf{life} would always have categorical desires, ... \\
03-a 4 & ...is glad to go on, and each stage of the  & \textbf{life} may be good for the person who lives it... \\
03-a 4 & But the question is whether the endless  & \textbf{life} is desirable for the person who now con... \\
03-a 4 & The  & \textbf{life} described is something like a successio... \\
03-a 5 & ... argument does not show that any finite  & \textbf{life} is too long, or that there is a determi... \\
03-a 5 & ...es not show that the limits of a mortal  & \textbf{life} are what give it its shape, which is a ... \\
03-a 6 & An endless  & \textbf{life} would eventually exhaust the categorica... \\
03-a 6 & ...haust the categorical desires that give  & \textbf{life} its point, or else cease, in its later ... \\
03-a 6 & ...se cease, in its later reaches, to be a  & \textbf{life} that gives its subject now what matters... \\
03-b Abs & EXI-03-a argues that an endless  & \textbf{life} would eventually exhaust the categorica... \\
03-b Abs & ...haust the categorical desires that give  & \textbf{life} its point, or cease to be a life that m... \\
03-b Abs & ...t give life its point, or cease to be a  & \textbf{life} that matters to the person now. \\
03-b Abs & ... is not a categorical desire that gives  & \textbf{life} its point, since its object is not a pr... \\
03-b Abs & Second, the undesirability of endless  & \textbf{life} is compatible with the badness of death... \\
03-b Abs & ...e exhaustion lies far beyond any actual  & \textbf{life}. \\
03-b Abs & ...und on the desirable length of a finite  & \textbf{life}: it shows only that some finite length ... \\
03-b 1 & ...ose satisfaction requires her continued  & \textbf{life}; it gives life its point if its object ... \\
03-b 1 & ...n requires her continued life; it gives  & \textbf{life} its point if its object is a project sh... \\
03-b 1 & ...3 of EXI-03-a concludes that an endless  & \textbf{life} is not desirable for the one who would ... \\
03-b 2 & ... is not a categorical desire that gives  & \textbf{life} its point. \\
03-b 2 & A categorical desire gives  & \textbf{life} its point only if its object is a proje... \\
03-b 2 & ... to live forever is the continuation of  & \textbf{life} itself, without specification of what i... \\
03-b 2 & ...d eventually contain nothing that gives  & \textbf{life} its point. \\
03-b 2 & That an endless  & \textbf{life} is not desirable does not imply that de... \\
03-b 2 & ...curs only after some time in an endless  & \textbf{life}; the theorem does not say when. \\
03-b 2 & At such a time the  & \textbf{life} the person would have had, had she not ... \\
03-b 2 & ...she not died, contains goods her actual  & \textbf{life} lacks, and by the deprivation account o... \\
03-b 2 & The theorem concerns the limit of a  & \textbf{life}; the badness of death concerns its actu... \\
03-b 2 & ... establishes that some finite length of  & \textbf{life} would suffice to exhaust the categorica... \\
03-b 2 & A  & \textbf{life} many times longer than ours might be de... \\
03-b 3 & If the desire for endless  & \textbf{life} is not categorical, then the fear of lo... \\
03-b 4 & Categorical desires give  & \textbf{life} its point and, being finite in an endle... \\
03-b 4 & ...s point and, being finite in an endless  & \textbf{life}, set a limit on the value of continuati... \\
04-a Abs & Is  & \textbf{life} absurd, and if so, how should one respo... \\
04-a Abs & ...ot in the smallness or brevity of human  & \textbf{life} but in a collision within the reflectiv... \\
04-a Abs & We argue that  & \textbf{life} is absurd in this sense, because both o... \\
04-a 1 & ...ny arguments for the absurdity of human  & \textbf{life} appeal to scale. \\
04-a 1 & If a  & \textbf{life} of seventy years is pointless, a life o... \\
04-a 1 & ...a life of seventy years is pointless, a  & \textbf{life} of eternity would be no less so; if our... \\
04-a 1 & ..., is a structural feature of reflective  & \textbf{life}: the collision between the seriousness ... \\
04-a 1 & Justifications within a  & \textbf{life} come to an end in ends that are not the... \\
04-a 1 & We argue that  & \textbf{life} is absurd in Nagel's sense even if thos... \\
04-a 2 & A  & \textbf{life} is absurd if its subject takes some end... \\
04-a 3 & EXI-01 argues that  & \textbf{life} is meaningful to the extent that one en... \\
04-a 3 & The  & \textbf{life} of any reflective agent capable of inte... \\
04-a 3 & By Definition 2.3, her  & \textbf{life} is absurd. \\
04-a 3 & ... a fitting response to the absurdity of  & \textbf{life}, and it is the only one of the three re... \\
04-a 4 & ...ch we pursue them is justified, and the  & \textbf{life} of someone engaged in them is not absur... \\
04-a 4 & The engaged and worthwhile  & \textbf{life} is less exposed to despair than the lif... \\
04-a 4 & ...ife is less exposed to despair than the  & \textbf{life} devoted to trivial ends, since despair ... \\
04-a 5 & ...ument does not show that any particular  & \textbf{life} is meaningful or worthwhile; that quest... \\
04-a 6 &  & \textbf{Life} is absurd in Nagel's sense: a reflectiv... \\
04-b Abs & EXI-04-a argues that  & \textbf{life} is absurd in Nagel's sense and that the... \\
04-b 1 & A  & \textbf{life} can be absurd without being futile, and... \\
04-b 1 & ...g futile, and a task can be futile in a  & \textbf{life} that is not absurd. \\
04-b 1 & By Theorem 3.3 of EXI-04-a, the  & \textbf{life} of every reflective agent is absurd, in... \\
04-b 1 & ...flective agent is absurd, including the  & \textbf{life} of an agent whose projects succeed and ... \\
04-b 1 & ... physician whose patients recover; that  & \textbf{life} is not futile. \\
04-b 1 & ...ing up the detached standpoint, and her  & \textbf{life} would then not be absurd in the sense o... \\
04-b 1 & ...id; its claim concerns the absurdity of  & \textbf{life}, which is a different thing. \\
04-b 3 & ...the second sense of Definition 1.1, his  & \textbf{life} is not meaningful however happy he is, ... \\
04-b 3 & ...s: such a Sisyphus is fulfilled but his  & \textbf{life} is not meaningful, because his object i... \\
04-b 3 & His happiness may be real; his  & \textbf{life}, as Camus describes it, is a paradigm o... \\
04-b 4 & Death limits a  & \textbf{life} without making its activities futile in... \\
05-a Abs & ...stinct from particular threats to one's  & \textbf{life}, is certain. \\
05-a 3 & ...th EXI-02, the end of the psychological  & \textbf{life} that the organism sustains. \\
05-a 4 & ...r threat, or the possibility that one's  & \textbf{life} will be short. \\
05-a 4 & ...ity, and where it is, it is a different  & \textbf{life} whose value is not guaranteed to be gre... \\
06-a 1 & ...losing his catalogue of the miseries of  & \textbf{life}, judged that the course of nature sugge... \\
06-a 4 & ...nd eventually the end of conditions for  & \textbf{life}. \\
06-a 4 & ... the same laws that bring death brought  & \textbf{life}. \\
06-b Abs & ...ndifference is irrelevant to meaning in  & \textbf{life}, though not to every question about mea... \\
06-b 2 & On the account of EXI-01, a  & \textbf{life} is meaningful to the extent that its su... \\
06-b 2 & It denies that meaning in  & \textbf{life} requires that good. \\
07-a Abs & Is happiness the goal of  & \textbf{life}? \\
07-a Abs & ...nd conclude that happiness is a good of  & \textbf{life} but not its goal. \\
07-a 1 & On both views happiness is the goal of  & \textbf{life} in a strong sense: it is the end that a... \\
07-a 1 & ...he end that all other ends serve, and a  & \textbf{life} is well directed to the extent that it ... \\
07-a 1 & ...tieth year by a change in his theory of  & \textbf{life}: those only are happy, he came to think... \\
07-a 1 & ...is family: happiness is not the goal of  & \textbf{life}, because it is a by-product of other pu... \\
07-a 3 & Happiness is not the goal of  & \textbf{life}, in the sense of the end that all other... \\
07-a 3 & ... that all other ends of a well-directed  & \textbf{life} serve. \\
07-a 3 & If happiness were the goal of  & \textbf{life} in that sense, a well-directed life wou... \\
07-a 3 & ... of life in that sense, a well-directed  & \textbf{life} would have happiness as its sole final ... \\
07-a 3 & By Theorem 3.3, such a  & \textbf{life} defeats the aim by which it is directed... \\
07-a 3 & A  & \textbf{life} directed at its goal in a way that defe... \\
07-a 3 & So a well-directed  & \textbf{life} does not have happiness as its sole fin... \\
07-a 4 & ...not happiness in our sense but the good  & \textbf{life} as a whole, which includes the goods th... \\
07-a 4 & The Aristotelian claim that the good  & \textbf{life} is the end of action is compatible with... \\
07-a 4 & ...s chosen for its own sake within such a  & \textbf{life} is not the agent's feeling good, but th... \\
07-a 4 & ...ular goods of experience, embedded in a  & \textbf{life} whose final ends lie elsewhere, and the... \\
07-a 4 & ...estion whether happiness is the goal of  & \textbf{life} concerns the lives human beings can lea... \\
07-a 4 & ...nd, the experience machine shows that a  & \textbf{life} of produced satisfaction is worse for t... \\
07-a 4 & ...s worse for the one who lives it than a  & \textbf{life} of real engagement with the same experi... \\
07-a 4 & ... on its own terms, it would show that a  & \textbf{life} whose goal is happiness so understood i... \\
07-a 4 & ...happiness so understood is not the best  & \textbf{life} for its subject, which is a further rea... \\
07-a 4 & ...n to deny that happiness is the goal of  & \textbf{life}. \\
07-a 6 & Happiness is not the goal of  & \textbf{life}. \\
07-b 2 & Happiness then figures in her  & \textbf{life} as the by-product of Definition 2.3 of ... \\
07-b 3 & ...asing to treat happiness as the goal of  & \textbf{life}. \\
08-a Abs & The ideal asks that a  & \textbf{life} be one's own rather than borrowed from ... \\
09-a Abs & ...tence, because a defining project for a  & \textbf{life} can be fixed only by the agent who live... \\
09-a 2 & ...lf afterwards, he means that a person's  & \textbf{life} comes to be for something only through ... \\
09-a 2 & ...eing capable of adopting ends, what its  & \textbf{life} is for, in the sense of Definition 2.1,... \\
09-a 2 & Another's intention for that  & \textbf{life} can give it reasons, but does not fix i... \\
09-a 2 & ...states that the success conditions of a  & \textbf{life}, as opposed to those of a tool, are int... \\
09-a 2 & ...hich one values oneself and finds one's  & \textbf{life} worth living (Korsgaard 1996). \\
09-a 3 & ...is capable of adopting ends and what 's  & \textbf{life} is for is fixed by the ends adopts. \\
09-a 3 & ...e lives, none of which fails as a human  & \textbf{life} by departing from the others. \\
09-a 3 & ...ition of a human being settles what her  & \textbf{life} is for. \\
09-a 4 & ...ixes a general form that any good human  & \textbf{life} takes, not what a particular life is fo... \\
09-a 4 & ...human life takes, not what a particular  & \textbf{life} is for. \\
09-a 5 & Principle P2 says that what a  & \textbf{life} is for is fixed by its agent's ends; it... \\
09-a 5 & ...d is consistent with EXI-01, on which a  & \textbf{life} is meaningful to the extent that its pr... \\
09-b 2 & ...ays what she is without saying what her  & \textbf{life} is for. \\
09-b 2 & ... as a human being in virtue of what her  & \textbf{life} is for. \\
09-b 2 & ...I-09-a, no defining project for a human  & \textbf{life} is fixed before it exists; whatever sta... \\
09-b 2 & ...xists; whatever standard of success her  & \textbf{life} has as hers is fixed by the ends she ad... \\
09-b 2 & A  & \textbf{life} cannot fail as a human life by departin... \\
09-b 2 & A life cannot fail as a human  & \textbf{life} by departing from a project it never ha... \\
09-b 2 & A  & \textbf{life} can be imprudent, wicked or, in the sen... \\
09-b 2 & ...efining project, and they do not make a  & \textbf{life} a defective specimen of anything. \\
09-b 2 & ... organism; that is not a failure of her  & \textbf{life} to be what it is for. \\
10-a 5 & ...nt with EXI-03's thesis that an endless  & \textbf{life} would exhaust the desires that give it ... \\
10-a 5 & ...sires that give it its point: in such a  & \textbf{life}, boredom would become a correct signal ... \\
10-b 3 & (2009) found that a lack of meaning in  & \textbf{life} predicted boredom cross-sectionally, ov... \\
10-b 8 & EXI-17 holds that a meaningful  & \textbf{life} requires active engagement with somethi... \\
10-b 8 & EXI-03 holds that an endless  & \textbf{life} would exhaust the categorical desires t... \\
10-b 8 & ...haust the categorical desires that give  & \textbf{life} its point; EXI-10-a draws the consequen... \\
10-b 10 & ..., in laboratories, schools and everyday  & \textbf{life}. \\
11-a Abs & Can  & \textbf{life} be meaningful without God? \\
11-a Abs & ...ng requires a purpose assigned to human  & \textbf{life} by its creator, or a relation to an inf... \\
11-a Abs & We argue that  & \textbf{life} can be meaningful without God, because ... \\
11-a Abs & ...itting-fulfilment account of meaning: a  & \textbf{life} is meaningful to the extent that its su... \\
11-a 1 & On purpose theories, a  & \textbf{life} is meaningful insofar as it fulfils a p... \\
11-a 1 & ...d without God, and the question whether  & \textbf{life} would be meaningful in such a world is ... \\
11-a 1 & ...nce: whether the conditions that make a  & \textbf{life} meaningful include God's existence, a r... \\
11-a 1 & ... practical interest, namely whether the  & \textbf{life} of someone who does not believe in God ... \\
11-a 2 & A  & \textbf{life} is meaningful to the extent that its su... \\
11-a 2 & A  & \textbf{life} has a cosmic purpose if a creator of th... \\
11-a 2 & A  & \textbf{life} fulfils its cosmic purpose if it plays ... \\
11-a 2 & The purpose theory holds that a  & \textbf{life} is meaningful only if it has a cosmic p... \\
11-a 2 & ... to more of what is worthwhile, makes a  & \textbf{life} more meaningful. \\
11-a 2 & ...the insight in Nozick's linkage view: a  & \textbf{life} gains meaning by connection with value ... \\
11-a 3 & ...pose is not sufficient for a meaningful  & \textbf{life}. \\
11-a 3 & ...tacles for its amusement, and suppose a  & \textbf{life} that fulfils that purpose. \\
11-a 3 & By Definition 2.1, such a  & \textbf{life} is not made meaningful by its purpose. \\
11-a 3 & ...es: a purpose of this kind degrades the  & \textbf{life} it is given to. \\
11-a 3 & A  & \textbf{life} that has no cosmic purpose, or does not... \\
11-a 3 & ...fies Definition 2.1, whether or not her  & \textbf{life} has a cosmic purpose and whether or not... \\
11-a 3 & Her  & \textbf{life} is therefore meaningful to that extent. \\
11-a 3 & The  & \textbf{life} of someone who does not believe in God ... \\
11-a 3 & ...ement with what is worthwhile makes her  & \textbf{life} meaningful. \\
11-a 4 & ...ork, if true, would make available to a  & \textbf{life} more meaning than it could have in a wo... \\
11-a 4 & ...re is; by Principle P2, it would make a  & \textbf{life} more meaningful than engagement with le... \\
11-a 4 & Each of these relates the  & \textbf{life} to more of what is worthwhile, and by P... \\
11-a 4 & ...available, and Theorem 3.3 ensures that  & \textbf{life} is meaningful nonetheless. \\
11-a 5 & ...9 argues that death, far from voiding a  & \textbf{life}, gives it a shape by limiting it. \\
11-a 5 & On the linkage view, a finite  & \textbf{life} has meaning by connection with somethin... \\
11-a 5 & A  & \textbf{life} linked to these is meaningful without f... \\
11-a 5 & The engaged and worthwhile  & \textbf{life} of a person without such a framework is... \\
11-a 5 & The word meaning in the meaning of  & \textbf{life} just asks for a purpose. \\
11-a 5 & A  & \textbf{life} without a purpose assigned to it has no... \\
11-a 5 & The question of what makes a  & \textbf{life} meaningful is asked and answered by bel... \\
11-a 5 & ...nd unbelievers alike in terms of what a  & \textbf{life} is devoted to and whether that is worth... \\
11-a 5 & ... a purpose may be assigned that makes a  & \textbf{life} worse. \\
11-a 5 & If a  & \textbf{life} can be fully meaningful without God, no... \\
11-a 5 & ... for theism to add; if theism adds, the  & \textbf{life} without God is deficient in meaning. \\
11-a 5 & A  & \textbf{life} can be meaningful, and deeply so, while... \\
11-a 5 & That a  & \textbf{life} could be more meaningful is not a defic... \\
11-a 5 & ...Theorem 3.3; it is true of every finite  & \textbf{life}. \\
11-a 6 & Nor does it show that every  & \textbf{life} without God is meaningful: meaning requ... \\
11-a 7 &  & \textbf{Life} can be meaningful without God. \\
11-b 1 & ...t concerns availability: it says what a  & \textbf{life} could engage with, not that every life ... \\
11-b 1 & ... life could engage with, not that every  & \textbf{life} in a theistic world is more meaningful ... \\
11-b 1 & ...tic world is more meaningful than every  & \textbf{life} in a naturalistic one. \\
11-b 3 & ...would add to the meaning available in a  & \textbf{life} by adding objects and relations of enga... \\
12-a Abs & ...nly for lives containing no good, and a  & \textbf{life} with more good than bad is better than ... \\
12-a Abs & ...e that matters to the one who lives the  & \textbf{life}: although never existing is not a condi... \\
12-a Abs & ...ting is not a condition anyone is in, a  & \textbf{life} with more good than bad is good for the... \\
12-a Abs & We conclude that a  & \textbf{life} containing more good than bad is better... \\
12-a 1 & Schopenhauer held that  & \textbf{life} is a business whose proceeds do not cov... \\
12-a 1 & The first is in what sense a  & \textbf{life} is better than non-existence, since no ... \\
12-a 2 & A  & \textbf{life} has more good than bad if the goods it ... \\
12-a 2 & ...he exists is impersonally better if her  & \textbf{life} has more good than bad (ETH-20). \\
12-a 2 & ...omparatively if it causes her to have a  & \textbf{life} that is good for her. \\
12-a 2 & But we commonly judge that a person's  & \textbf{life} is good for her, that a pleasure is goo... \\
12-a 3 & A  & \textbf{life} containing more good than bad is better... \\
12-a 3 & For the second clause: a  & \textbf{life} with more good than bad is good for its... \\
12-a 3 & ...inciple P2, to be caused to have such a  & \textbf{life} is a non-comparative benefit. \\
12-a 3 & So the  & \textbf{life} is better than non-existence in both se... \\
12-a 3 & To yield the conclusion that every  & \textbf{life} is worse than non-existence, the argume... \\
12-a 3 & ...the argument would need to show that no  & \textbf{life} has more good than bad. \\
12-a 3 & These are not bads in the  & \textbf{life} but goods it does not contain; to count... \\
12-a 4 & It says that my  & \textbf{life} is good for me, which is a non-comparat... \\
12-a 4 & ...is is conditional on the balance of the  & \textbf{life}. \\
12-a 4 & Even if a good  & \textbf{life} is better than never existing, no one c... \\
12-a 4 & ...r existing, no one can guarantee that a  & \textbf{life} will be good. \\
12-a 4 & ...er uncertainty, not the comparison of a  & \textbf{life} with non-existence. \\
12-a 5 & ...t better never to have been born, for a  & \textbf{life} containing more good than bad. \\
12-a 5 & Such a  & \textbf{life} makes the world better than it would ha... \\
12-a 5 & ...tar's argument from the poor quality of  & \textbf{life} contests which lives have more good tha... \\
12-b Abs & ...re fails, and EXI-12-a concludes that a  & \textbf{life} containing more good than bad is better... \\
12-b 1 & ... is a duty not to create a person whose  & \textbf{life} would be miserable, but no duty to crea... \\
12-b 1 & ...e, but no duty to create a person whose  & \textbf{life} would be happy. \\
12-b 2 & ...e miserable person, once created, has a  & \textbf{life} that is bad for her and so has a compla... \\
12-b 2 & ...ason: by Theorem 3.1 of EXI-12-a a good  & \textbf{life} is a non-comparative benefit to the one... \\
12-b 3 & ...t support to the basic asymmetry, and a  & \textbf{life} with more good than bad remains better ... \\
13-a Abs & ... to rely on its doing so, and plans her  & \textbf{life} on that reliance; the duty to the dead ... \\
13-a 1 & ...at happens after death can affect how a  & \textbf{life} went is older: Aristotle asks whether t... \\
13-a 2 & ... the deprivation account of EXI-02: her  & \textbf{life} as a whole goes worse than it would oth... \\
13-a 2 & ...erests rather than because they cut her  & \textbf{life} short. \\
13-a 3 & ...by Principle P1 their effect on how her  & \textbf{life} went is smaller. \\
13-b 2 & ...us event is simply a fact about how her  & \textbf{life} went. \\
14-a 1 & ...social connectedness and the sense that  & \textbf{life} is meaningful (Wildschut et al. 2006; S... \\
15-b 2 & ...xcess is the player who reorganizes her  & \textbf{life} around the draw (Lemma 3.2 of EXI-15-a)... \\
17-a Abs & Does a meaningful  & \textbf{life} require achievement? \\
17-a Abs & ...e success a central place: a meaningful  & \textbf{life} is one that makes a difference, leaves ... \\
17-a Abs & A meaningful  & \textbf{life} requires active engagement with somethi... \\
17-a Abs & ...meaningful, and judges the unsuccessful  & \textbf{life} to be far from meaningless. \\
17-a 1 & ...cause of what they brought about, and a  & \textbf{life} that brings nothing about is to that ex... \\
17-a 1 & ... to roll stones, and concluded that his  & \textbf{life} would then have all the meaning a life ... \\
17-a 1 & ... life would then have all the meaning a  & \textbf{life} can have, though he achieved nothing. \\
17-a 1 & ...'s subjectivism was rightly rejected: a  & \textbf{life} devoted with passion to rolling stones,... \\
17-a 2 & A  & \textbf{life} is meaningful to the extent that its su... \\
17-a 2 & ...ing from Nozick's experience machine: a  & \textbf{life} of experiences as of curing diseases, l... \\
17-a 2 & ...iseases, lived in the machine, is not a  & \textbf{life} engaged with the sick. \\
17-a 3 & A person spends much of her  & \textbf{life} listening to and studying music she lov... \\
17-a 3 & The unsuccessful mathematician's  & \textbf{life} is meaningful, and so is the life of th... \\
17-a 3 & ...ian's life is meaningful, and so is the  & \textbf{life} of the nurse in the second ward. \\
17-a 3 & ...is not the verdict of reflection that a  & \textbf{life} devoted to an unsolved problem, or to t... \\
17-a 3 & The listener's  & \textbf{life} is meaningful to the extent of her enga... \\
17-a 3 & A meaningful  & \textbf{life} requires active engagement with somethi... \\
17-a 3 & ... what reflection counts as a meaningful  & \textbf{life} (Wolf 2010). \\
17-a 3 & ...aning, the unsuccessful mathematician's  & \textbf{life}, the life of the nurse in the second wa... \\
17-a 3 & ... unsuccessful mathematician's life, the  & \textbf{life} of the nurse in the second ward, and th... \\
17-a 3 & ... in the second ward, and the listener's  & \textbf{life} would lack it, so far as their central ... \\
17-a 3 & ...s equal, the successful mathematician's  & \textbf{life} is more meaningful than the unsuccessfu... \\
17-a 3 & ...e good (ETH-01), and it adds value to a  & \textbf{life}; if success enlarges the worth realized... \\
17-a 3 & ... between a meaningful and a meaningless  & \textbf{life}. \\
17-a 4 & So every meaningful  & \textbf{life} contains achievement after all. \\
17-a 4 & ...athematician did not attain it, and her  & \textbf{life} is meaningful. \\
17-a 4 & A  & \textbf{life} that changes nothing is enclosed within... \\
17-a 4 & ETH-05 shows that a  & \textbf{life} in contact with reality is better than ... \\
17-a 4 & ...nd its comparison contrasts a veridical  & \textbf{life} with an experiential duplicate. \\
17-a 4 & ...r work bears on a real problem, and her  & \textbf{life} is in contact with reality. \\
17-a 4 & ... concerns success bears on how well her  & \textbf{life} goes for her, and we grant that success... \\
17-a 4 & ...better; it does not bear on whether her  & \textbf{life} is meaningful. \\
17-a 4 & ...d achieving nothing, lives a meaningful  & \textbf{life}. \\
17-a 4 & ...e from flooding would live a meaningful  & \textbf{life}, even if the stone always rolled back a... \\
17-a 5 & A meaningful  & \textbf{life} requires active engagement with somethi... \\
18-a Abs & ...false judgments may be indispensable to  & \textbf{life}, and a body of psychological research h... \\
18-a 1 & ...gments may be the most indispensable to  & \textbf{life} (Nietzsche 1886). \\
18-a 2 & ...unction, and the ordinary progress of a  & \textbf{life} consists partly in such corrections. \\
18-a 2 & ...ch illusion is a respect in which one's  & \textbf{life} is out of contact with how things are. \\
18-b 1 & ...ch illusion is a respect in which one's  & \textbf{life} is out of contact with how things are. \\
19-a Abs & ... death only through its relation to the  & \textbf{life} it ends. \\
19-a Abs & But death gives  & \textbf{life} a shape by limiting it. \\
19-a Abs & ...lowing Velleman, we take the value of a  & \textbf{life} to depend in part on its shape, the ord... \\
19-a Abs & ...ds and evils occur, and we argue that a  & \textbf{life} has a shape as a whole only if it is bo... \\
19-a Abs & ...ble with EXI-02: the limit that gives a  & \textbf{life} a shape is not a benefit to the one who... \\
19-a Abs & ...a later death would also have given her  & \textbf{life} a shape, and in the ordinary case a bet... \\
19-a Abs & ...ts EXI-03, which argues that an endless  & \textbf{life} is not desirable. \\
19-a 1 & ...nd whether it bears on the point of the  & \textbf{life} it ends. \\
19-a 1 & ...he wage of sin, as a passage to another  & \textbf{life}, or as the moment of judgment. \\
19-a 1 & ...en found its meaning in its relation to  & \textbf{life}. \\
19-a 1 & ...leman (1991) showed that the value of a  & \textbf{life} is not the sum of the values of its mom... \\
19-a 1 & ...m of the values of its moments, since a  & \textbf{life} that improves from misfortune to succes... \\
19-a 1 & ... has no meaning in itself, but it gives  & \textbf{life} a shape by limiting it. \\
19-a 2 & ...s independently of its relations to the  & \textbf{life} in which it occurs: if it is the fulfil... \\
19-a 2 & ...of a value that it does not owe to that  & \textbf{life}. \\
19-a 2 & The shape of a  & \textbf{life} is the temporal order of its goods and ... \\
19-a 2 & The value of a  & \textbf{life} as a whole depends in part on its shape... \\
19-a 2 & ...less view of time defended in TSM-01, a  & \textbf{life} is an ordered series of events, each lo... \\
19-a 3 & ...imposes it or that it lead to a further  & \textbf{life}; neither can be established, and REL-08... \\
19-a 3 & ... in virtue of its relation to a further  & \textbf{life} or a further purpose, not in itself. \\
19-a 3 & ...on is a claim about its relation to the  & \textbf{life} it ends, and it is often false: most de... \\
19-a 3 & ...eath that are true independently of any  & \textbf{life}, but they assign it no significance or ... \\
19-a 3 & Death gives a  & \textbf{life} a shape by limiting it. \\
19-a 3 & A  & \textbf{life} is an ordered series of events (Assumpt... \\
19-a 3 & ... of being limited by death that a human  & \textbf{life} has a shape as a whole, and by Principl... \\
19-a 3 & ...s shape contributes to the value of the  & \textbf{life}. \\
19-a 3 & ...that depends on the goods and evils the  & \textbf{life} contains and their order. \\
19-a 3 & ...ference between the value of her actual  & \textbf{life} and the value of the life she would hav... \\
19-a 3 & ...of her actual life and the value of the  & \textbf{life} she would have had had she not died the... \\
19-a 3 & ...ie at that time, she dies later, so her  & \textbf{life} is limited and has a shape there too. \\
19-a 3 & ...he conclusion of EXI-03 that an endless  & \textbf{life} is not desirable, but does not depend o... \\
19-a 3 & EXI-03 argues that an endless  & \textbf{life} would exhaust the categorical desires t... \\
19-a 3 & ...haust the categorical desires that give  & \textbf{life} its point; Theorem 3.2 adds that it wou... \\
19-a 3 & The two are independent: a  & \textbf{life} could lack a shape as a whole and still... \\
19-a 4 & Projects, relationships and stages of  & \textbf{life} have beginnings and ends, and give life... \\
19-a 4 & ...life have beginnings and ends, and give  & \textbf{life} all the shape it needs. \\
19-a 4 & They give shape to the episodes of a  & \textbf{life}, and an endless life could have shaped ... \\
19-a 4 & ... the episodes of a life, and an endless  & \textbf{life} could have shaped episodes. \\
19-a 4 & ...at they do not give is the shape of the  & \textbf{life} as a whole, which is what Principle P1 ... \\
19-a 4 & ...what Principle P1 concerns: whether the  & \textbf{life}, taken as one thing, was a story of pro... \\
19-a 4 & That a  & \textbf{life} is limited is a condition of its having... \\
19-a 4 & Theorem 3.2 says that death gives a  & \textbf{life} a shape, not a good one. \\
19-a 4 & An early death gives a  & \textbf{life} the shape of a beginning cut off, which... \\
19-a 4 & ...off, which is one of the worst shapes a  & \textbf{life} can have; this is part of why, by EXI-0... \\
19-a 5 & ...elong to it through its relation to the  & \textbf{life} it ends. \\
19-a 5 & But it limits that  & \textbf{life}, and by limiting it gives it a shape as... \\
19-a 5 & ... a whole, which is part of what makes a  & \textbf{life} go well or badly. \\
19-b Abs & ...eath has no meaning in itself but gives  & \textbf{life} a shape by limiting it. \\
19-b Abs & ...ition of one kind of meaning, that of a  & \textbf{life} as a whole, which depends on the life's... \\
19-b 1 & On the account of EXI-01 and EXI-17, a  & \textbf{life} is meaningful to the extent that its su... \\
19-b 1 & ... it, as a condition of the meaning of a  & \textbf{life} as a whole, or as a condition of the pe... \\
19-b 1 & That a  & \textbf{life} is limited by death is not something wh... \\
19-b 1 & ...ot something whose engagement makes the  & \textbf{life} meaningful, and does not make anything ... \\
19-b 1 & In so far as the meaning of a  & \textbf{life} as a whole depends on its shape, finitu... \\
19-b 1 & Some meaning belongs to a  & \textbf{life} only as a whole: that it was a life of ... \\
19-b 1 & ...o a life only as a whole: that it was a  & \textbf{life} of growing understanding, of a commitme... \\
19-b 1 & So the meaning that a  & \textbf{life} has as a whole requires that it be limi... \\
19-b 1 & By the result of EXI-03, an endless  & \textbf{life} would eventually exhaust the categorica... \\
19-b 1 & ...ust the categorical desires that give a  & \textbf{life} its point (Williams 1973). \\
19-b 1 & A  & \textbf{life} without categorical desires lacks the a... \\
19-b 1 & So an endless  & \textbf{life} would at some point cease to be meaning... \\
19-b 1 & ...ome point cease to be meaningful, and a  & \textbf{life} remains meaningful throughout its cours... \\
19-b 2 & ... temporal scarcity, so that an immortal  & \textbf{life} would lack the structure in which valui... \\
19-b 2 & Scheffler's point that an immortal  & \textbf{life} would lack the stages of a human life i... \\
19-b 2 & ...l life would lack the stages of a human  & \textbf{life} is a point about shape, and Nussbaum's ... \\
19-b 2 & ...um's point that some values belong to a  & \textbf{life} of a certain form is a point about the ... \\
19-b 2 & ...a later death would also have given the  & \textbf{life} a shape and would, in the ordinary case... \\
19-b 3 & Finitude does not make  & \textbf{life} meaningful. \\
19-b 3 & It makes possible the meaning of a  & \textbf{life} as a whole, through the shape that a li... \\
19-b 3 & ...pe that a limit gives, and it keeps the  & \textbf{life} within the span in which engagement per... \\
19-b 3 & ... in itself, bears on the meaning of the  & \textbf{life} it ends. \\
20-a Abs & EXI-17 argues that a meaningful  & \textbf{life}, or a meaningful activity within one, r... \\
20-a Abs & ...ther the process has the kind of mental  & \textbf{life} that engagement requires, which MND-08 ... \\
20-a 1 & Wolf (2010) argued that meaning in  & \textbf{life} arises when subjective attraction meets... \\
20-a 1 & ...ction meets objective attractiveness: a  & \textbf{life} is meaningful to the extent that its su... \\
20-a 1 & ...he account to an activity rather than a  & \textbf{life}: the production of the present collecti... \\
20-a 4 & ... from inside it: a fact about the inner  & \textbf{life} of the producer and a fact about the up... \\
20-b Abs & ... EXI-20-a, on whether it has the mental  & \textbf{life} that engagement requires, or on whether... \\
\end{longtable}
\kwichead{matters}{37}
\begin{longtable}{@{}>{\raggedright}p{0.07\textwidth}>{\raggedleft}p{0.41\textwidth}@{\,}>{\raggedright\arraybackslash}p{0.47\textwidth}@{}}
01-c 3 & The difference  & \textbf{matters} because opportunities for projects in t... \\
01-d 2 & Deprivation therefore  & \textbf{matters} to meaning, and to a theory of justice ... \\
01-d 4 & ...ufficiently detached standpoint nothing  & \textbf{matters}, so meaning is available to no one. \\
02-a 3 & ... equally sudden differ in badness: what  & \textbf{matters} is not the manner of dying but the life... \\
02-a 6 & ...that connectedness carries most of what  & \textbf{matters} in survival. \\
02-b 2 & This  & \textbf{matters} because Epicurus's therapeutic aim, the... \\
03-a Abs & ...ife would not be hers in the sense that  & \textbf{matters}. \\
03-a Abs & ...that connectedness carries most of what  & \textbf{matters} in survival, we show that an endless li... \\
03-a Abs & ...o be, in its later reaches, a life that  & \textbf{matters} to the person now. \\
03-a 1 & ...y of character with the account of what  & \textbf{matters} in survival defended in PID-06. \\
03-a 2 & What  & \textbf{matters} to a person, in her concern for a futur... \\
03-a 2 & ...ife gives the person now little of what  & \textbf{matters}. \\
03-a 3 & ...ives her, as she now is, little of what  & \textbf{matters}. \\
03-a 3 & ... that time gives her now little of what  & \textbf{matters}. \\
03-a 3 & ...hout connectedness carry little of what  & \textbf{matters}. \\
03-a 3 & ... cease to be a life that gives her what  & \textbf{matters}. \\
03-a 3 & ...ter reaches give her now little of what  & \textbf{matters}. \\
03-a 4 & ... through a chain carries little of what  & \textbf{matters}. \\
03-a 6 & ... a life that gives its subject now what  & \textbf{matters} to her. \\
03-a 6 & ...that connectedness carries most of what  & \textbf{matters}. \\
03-b Abs & ...e its point, or cease to be a life that  & \textbf{matters} to the person now. \\
03-b 1 & ...mption 2.4), and the dependence of what  & \textbf{matters} on connectedness (Assumption 2.5), Theo... \\
06-a 4 & ...re about is hostile in every sense that  & \textbf{matters}. \\
06-b 3 & ...e has no such point of view, so nothing  & \textbf{matters} from it. \\
07-a 2 & ...d; it presupposes that its being solved  & \textbf{matters}, and an agent to whom it matters only a... \\
07-a 2 & ...solved matters, and an agent to whom it  & \textbf{matters} only as a source of satisfaction has no... \\
10-a 4 & But the signal bears on a question that  & \textbf{matters} independently of its evolutionary histo... \\
10-b 2 & The distinction  & \textbf{matters} for how evidence is read. \\
10-b 2 & ...ross-tabulation in which the difference  & \textbf{matters}. \\
10-b 3 & ...without traction, and the qualification  & \textbf{matters}: an over-demanding task on which one ma... \\
12-a Abs & ...arison can be made good in a sense that  & \textbf{matters} to the one who lives the life: although... \\
13-a 2 & ...answers that question for the case that  & \textbf{matters} here: the entitlement is created while ... \\
13-a 3 & A wish that fails clause (ii) concerns  & \textbf{matters} on which the person has no claim agains... \\
17-a 3 & ...he sense of Definition 2.2: the problem  & \textbf{matters} to knowledge, which is a component of t... \\
17-a 3 & ... the good, and the patients' well-being  & \textbf{matters} whether or not it can be restored. \\
17-a 4 & This is an argument that success  & \textbf{matters}. \\
19-b 2 & ...rgency adds no worth: the friend's good  & \textbf{matters} no less to someone who will have other ... \\
\end{longtable}
\kwichead{meaning}{229}
\begin{longtable}{@{}>{\raggedright}p{0.07\textwidth}>{\raggedleft}p{0.41\textwidth}@{\,}>{\raggedright\arraybackslash}p{0.47\textwidth}@{}}
01-c Abs & ...1-a establishes that a life of positive  & \textbf{meaning} is available to every agent capable of ... \\
01-c Abs & ...the sense relevant to the definition of  & \textbf{meaning}. \\
01-c 1 & ...leave people with few opportunities for  & \textbf{meaning}. \\
01-c 2 & ...st follows the objects that accounts of  & \textbf{meaning} have most often recognized: persons, th... \\
01-c 2 & ... intentional action, a life of positive  & \textbf{meaning} is available. \\
01-c 3 & ... applies the same point to opportunity:  & \textbf{meaning} does not wait on the availability of a ... \\
01-c 4 & ...ures of projects that are irrelevant to  & \textbf{meaning} on Definition 1: a completion condition... \\
01-c 5 & The universal availability of  & \textbf{meaning} does not depend on the universal availa... \\
01-c 5 & ... worth, and a life that contains it has  & \textbf{meaning}. \\
01-d Abs & ...h projects of objective worth, and that  & \textbf{meaning} in this sense is available to everyone. \\
01-d Abs & It implies that  & \textbf{meaning} is not hostage to luck in outcomes, to ... \\
01-d Abs & ...e and injustice affect is the degree of  & \textbf{meaning}, not its possibility. \\
01-d Abs & ...ful, or that deprivation is harmless to  & \textbf{meaning}. \\
01-d 1 & ...f intentional action a life of positive  & \textbf{meaning} is available. \\
01-d 1 & ...same availability by a different route:  & \textbf{meaning} requires active engagement with somethi... \\
01-d 1 & ...ment (Theorem 3.3 of EXI-17-a), so that  & \textbf{meaning} is not hostage to luck in outcomes (Cor... \\
01-d 2 & If  & \textbf{meaning} required success, it would be available... \\
01-d 2 & Because  & \textbf{meaning} requires engagement, which is within th... \\
01-d 2 & Definition 1 of EXI-01-a measures  & \textbf{meaning} by engagement and by the objective wort... \\
01-d 2 & ...an achieve, which is not a condition of  & \textbf{meaning}; scale bears on how much value a projec... \\
01-d 2 & ...worth, and engagement with them confers  & \textbf{meaning} on lives that will leave no public trac... \\
01-d 2 & ...cluding injustice, affect the degree of  & \textbf{meaning} a life can attain and the ease with whi... \\
01-d 2 & ...can attain it, but not whether positive  & \textbf{meaning} is available. \\
01-d 2 & Definition 1 of EXI-01-a makes  & \textbf{meaning} a matter of degree, and EXI-17-a notes ... \\
01-d 2 & Deprivation therefore matters to  & \textbf{meaning}, and to a theory of justice that counts... \\
01-d 2 & ... cannot do is remove the possibility of  & \textbf{meaning} altogether. \\
01-d 3 & ...availability to equality: the degree of  & \textbf{meaning} varies with the extent of engagement an... \\
01-d 3 & ...from availability to indifference: that  & \textbf{meaning} remains possible under oppression is no... \\
01-d 3 & ...oppression, which reduces the degree of  & \textbf{meaning} as it reduces other goods. \\
01-d 4 & Objective-worth views of  & \textbf{meaning} privilege the projects valued by the ed... \\
01-d 4 & ...ience and political causes, and so make  & \textbf{meaning} in practice available mainly to the pri... \\
01-d 4 & If  & \textbf{meaning} is available to everyone, the claim tha... \\
01-d 4 & That positive  & \textbf{meaning} is available to everyone is compatible ... \\
01-d 4 & The claim is not that  & \textbf{meaning} is cheap but that it is not reserved: n... \\
01-d 4 & ...detached standpoint nothing matters, so  & \textbf{meaning} is available to no one. \\
01-d 5 &  & \textbf{Meaning}, on the account of EXI-01-a, is open to... \\
01-d 5 & The availability of  & \textbf{meaning} is universal; its degree is not, and ho... \\
03-a 1 & ...worthwhile, in line with the account of  & \textbf{meaning} in EXI-01; and it replaces Williams's a... \\
03-a 2 & ...ch engagement, not achievement, is what  & \textbf{meaning} requires. \\
04-a 1 & ...y, dissolves the absurd: on theories of  & \textbf{meaning} such as those of Wolf (2010) and Metz (... \\
04-a 2 & ...he world as an adversary that denies us  & \textbf{meaning}, and the agent's persistence as a victo... \\
04-a 3 & ...he world as an adversary that denies us  & \textbf{meaning}. \\
04-a 4 &  & \textbf{Meaning} requires active engagement with somethi... \\
04-a 4 & Irony is the enemy of  & \textbf{meaning}. \\
04-b 3 & A further consequence concerns  & \textbf{meaning}. \\
04-b 3 & ...t meaningful however happy he is, since  & \textbf{meaning} requires engagement with something wort... \\
04-b 3 & ...bes it, is a paradigm of the absence of  & \textbf{meaning} rather than of its presence in a godles... \\
04-b 3 & ...jects, and engagement with them confers  & \textbf{meaning} (EXI-01). \\
06-b Abs & ...ifferent to us leaves our lives without  & \textbf{meaning}, or with a meaning we must assert in de... \\
06-b Abs & ...es our lives without meaning, or with a  & \textbf{meaning} we must assert in defiance of it. \\
06-b Abs &  & \textbf{Meaning}, on the account of EXI-01, consists in ... \\
06-b Abs & Cosmic indifference is irrelevant to  & \textbf{meaning} in life, though not to every question a... \\
06-b Abs & ...ife, though not to every question about  & \textbf{meaning}. \\
06-b 1 & We ask what this implies for  & \textbf{meaning}. \\
06-b 2 & ...o us does not imply that our lives lack  & \textbf{meaning}. \\
06-b 2 & It denies that  & \textbf{meaning} in life requires that good. \\
06-b 2 & ...its indifference denies us nothing that  & \textbf{meaning} requires; by Theorem 3.1 of EXI-06-a, i... \\
06-b 3 &  & \textbf{Meaning} requires that our lives matter from the... \\
06-b 3 & ... does indifference deprive our lives of  & \textbf{meaning}. \\
06-b 3 & Theistic accounts hold that  & \textbf{meaning} requires a purpose conferred by a being... \\
06-b 3 & If the universe does not care,  & \textbf{meaning} is lost unless there is such a being. \\
06-b 3 & The present point is conditional: if  & \textbf{meaning} consists in engagement with what has ob... \\
06-b 3 & ...estion, and supernaturalist accounts of  & \textbf{meaning} are assessed elsewhere (Metz 2013). \\
06-b 4 & It leaves the question of  & \textbf{meaning} where EXI-01 places it: in what we enga... \\
07-a 3 & It fits also the account of  & \textbf{meaning} in EXI-01 and EXI-17, on which meaning ... \\
07-a 3 & ... meaning in EXI-01 and EXI-17, on which  & \textbf{meaning} comes from active engagement with what ... \\
07-a 5 & ...the account of the good (ETH-01) and of  & \textbf{meaning} (EXI-01, EXI-17). \\
10-a Abs & ...challenge without traction, and lack of  & \textbf{meaning}; and it motivates search for alternativ... \\
10-a Abs & ...dom is aversive, why it is sensitive to  & \textbf{meaning} as well as to difficulty, and why chron... \\
10-a 1 & ...ty, the condition of a culture in which  & \textbf{meaning} has to be produced rather than inherite... \\
10-a 1 & ...lson (2018) distinguish attentional and  & \textbf{meaning} components: one is bored when one's att... \\
10-a 2 & The  & \textbf{meaning} clause is why the mismatch is between c... \\
10-a 3 & It explains why boredom is sensitive to  & \textbf{meaning}: the mismatch it signals includes the i... \\
10-a 4 & ...bout the agent's good on the account of  & \textbf{meaning} in EXI-17, on which meaning requires ac... \\
10-a 4 & ... account of meaning in EXI-17, on which  & \textbf{meaning} requires active engagement with somethi... \\
10-b Abs & ...tly evidence about motivation, that the  & \textbf{meaning} and under-challenge clauses of the mism... \\
10-b 2 & ...e resources in either direction, from a  & \textbf{meaning} component, the activity's failure to ma... \\
10-b 3 & (2009) found that a lack of  & \textbf{meaning} in life predicted boredom cross-section... \\
10-b 5 & ...city to attend; they say less about its  & \textbf{meaning} clause, which concerns the capacity to ... \\
10-b 5 & ... the entries coded attn and those coded  & \textbf{meaning}+ overlap only in the MAC model itself (... \\
10-b 6 & ...or by meaninglessness alone, except the  & \textbf{meaning} manipulation reported by Fahlman et al. \\
10-b 8 & ...gagement with something worthwhile; the  & \textbf{meaning} clause of the mismatch thesis is the af... \\
10-b 8 & ...hat view, and the evidence that lack of  & \textbf{meaning} predicts boredom (Fah09, Til12, Til17) ... \\
10-b 11 & ...match thesis in its under-challenge and  & \textbf{meaning} clauses and in its converse, that match... \\
10-b 11 & ...how they engage, and turn to sources of  & \textbf{meaning}. \\
10-b A & ...sis if it carries one of under+, over+,  & \textbf{meaning}+, fit-, attn or relational, and challen... \\
10-b B & ... and overanx (against); meaninglessness  & \textbf{meaning}+; Assumption 2.4 fit-; Lemma 3.1 signal... \\
10-b B & ... 3.3 fit- and the codes of goal change,  & \textbf{meaning} regulation and creativity (for) and har... \\
11-a Abs & Supernaturalists hold that  & \textbf{meaning} requires a purpose assigned to human li... \\
11-a Abs & ... can be meaningful without God, because  & \textbf{meaning} does not require a cosmic purpose, and ... \\
11-a Abs & ...ts on the fitting-fulfilment account of  & \textbf{meaning}: a life is meaningful to the extent tha... \\
11-a Abs & ...ned by another is neither necessary for  & \textbf{meaning} nor sufficient for it. \\
11-a Abs & ...ond claim rests on the observation that  & \textbf{meaning} comes in degrees and grows with the wor... \\
11-a 1 & ...God there is no such purpose, and so no  & \textbf{meaning} (Craig 2008). \\
11-a 1 & On theories of the infinite,  & \textbf{meaning} requires a connection to something of u... \\
11-a 1 & ...am's more cautious view, the pursuit of  & \textbf{meaning} requires a moral framework and a ground... \\
11-a 1 & Naturalists reply that  & \textbf{meaning} is available in a world without God. \\
11-a 1 & On subjectivist views,  & \textbf{meaning} is a matter of what a person cares abou... \\
11-a 1 & Wolf (2010) defends a hybrid:  & \textbf{meaning} arises when the subject's active engage... \\
11-a 1 & ...l-centered theories as accounts of what  & \textbf{meaning} requires, and defends a naturalist theo... \\
11-a 1 & ...nd defends a naturalist theory on which  & \textbf{meaning} is a matter of orienting one's rational... \\
11-a 1 & ...lism, on which God is necessary for any  & \textbf{meaning}, from moderate supernaturalism, on whic... \\
11-a 1 & ...ecessary only for a greater or ultimate  & \textbf{meaning}. \\
11-a 1 & ...ut the framework, section 3 argues that  & \textbf{meaning} does not require cosmic purpose, and se... \\
11-a 1 & ... that a theistic framework would add to  & \textbf{meaning}. \\
11-a 2 & EXI-01 adds that  & \textbf{meaning} so understood is available to everyone. \\
11-a 2 & ...implicit in Definition 2.1, which makes  & \textbf{meaning} a matter of extent. \\
11-a 2 & ... in Nozick's linkage view: a life gains  & \textbf{meaning} by connection with value beyond itself,... \\
11-a 3 &  & \textbf{Meaning} does not require a cosmic purpose. \\
11-a 4 & ...ue, would make available to a life more  & \textbf{meaning} than it could have in a world without G... \\
11-a 4 & ...d by Principle P2 each would add to its  & \textbf{meaning}. \\
11-a 4 & If REL-01 is right, the additional  & \textbf{meaning} is actually available, and is available... \\
11-a 4 &  & \textbf{Meaning} has a floor that does not depend on God... \\
11-a 4 & ...s critics that God is not necessary for  & \textbf{meaning}, and with moderate supernaturalism that... \\
11-a 4 & ...upernaturalism that God would make more  & \textbf{meaning} available (Metz 2013); it does not foll... \\
11-a 4 & ...d necessary for any particular grade of  & \textbf{meaning}. \\
11-a 5 & ...4.1: preservation of the good would add  & \textbf{meaning}; its absence does not remove the meanin... \\
11-a 5 & ...eaning; its absence does not remove the  & \textbf{meaning} already there. \\
11-a 5 & On the linkage view, a finite life has  & \textbf{meaning} by connection with something larger. \\
11-a 5 & That larger thing has  & \textbf{meaning} only if it is connected with something ... \\
11-a 5 & ...unlimited terminus expresses a wish for  & \textbf{meaning} without limit, not a condition of meani... \\
11-a 5 & ...aning without limit, not a condition of  & \textbf{meaning} as such. \\
11-a 5 & ...ants that an unlimited object would add  & \textbf{meaning}; it does not follow that without it the... \\
11-a 5 & The pursuit of  & \textbf{meaning} is not a solitary engagement with isola... \\
11-a 5 & ...thout God, these are fragile, and so is  & \textbf{meaning} (Cottingham 2003). \\
11-a 5 & ...ction describes affects the security of  & \textbf{meaning} and the reasonableness of certain hopes... \\
11-a 5 & The word  & \textbf{meaning} in the meaning of life just asks for a ... \\
11-a 5 & The word meaning in the  & \textbf{meaning} of life just asks for a purpose. \\
11-a 5 & ...without a purpose assigned to it has no  & \textbf{meaning} by definition, and the naturalist has c... \\
11-a 5 & ...s, the life without God is deficient in  & \textbf{meaning}. \\
11-a 5 & The dilemma assumes that  & \textbf{meaning} is all or nothing. \\
11-a 6 & It does not show how much  & \textbf{meaning} theism would add, nor that the addition... \\
11-a 6 & ...t every life without God is meaningful:  & \textbf{meaning} requires engagement with what is worthw... \\
11-a 7 &  & \textbf{Meaning} consists in active engagement with what... \\
11-a 7 & ...igned purpose is neither sufficient for  & \textbf{meaning} nor necessary for it. \\
11-a 7 & A theistic framework would add to  & \textbf{meaning}, by offering an object of maximal worth... \\
11-a 7 & The floor of  & \textbf{meaning} is independent of God; the ceiling, if ... \\
11-b Abs & EXI-11-a argues that  & \textbf{meaning} does not require a cosmic purpose, and ... \\
11-b Abs & ..., if true, would nevertheless make more  & \textbf{meaning} available (Proposition 4.1). \\
11-b Abs & ...at God adds is not more of the same but  & \textbf{meaning} as such; the anti-theist, following Kah... \\
11-b Abs & ...e, says that God's existence would take  & \textbf{meaning} away, by removing the privacy, independ... \\
11-b Abs & We argue that the additional  & \textbf{meaning} is a difference of degree, and that the... \\
11-b 1 & EXI-11-a defines  & \textbf{meaning} as active engagement with what is worth... \\
11-b 1 & ...thwhile (Definition 2.1) and holds that  & \textbf{meaning} grows with the worth of what one engage... \\
11-b 1 & ...worthwhile (Principle P2; on degrees of  & \textbf{meaning}, see (Metz 2013)). \\
11-b 1 & The  & \textbf{meaning} that a theistic framework would add is ... \\
11-b 1 & ...rk would add is of the same kind as the  & \textbf{meaning} available without it, and differs in de... \\
11-b 1 & By Definition 2.1 of EXI-11-a, any  & \textbf{meaning} is engagement with what is worthwhile. \\
11-b 1 & ...agement with what is worthwhile; so the  & \textbf{meaning} it yields falls under the same definiti... \\
11-b 1 & ...hout it can be meaningful, so it is not  & \textbf{meaning} as such. \\
11-b 1 & ...EXI-11-a, it can increase the degree of  & \textbf{meaning}. \\
11-b 2 & ...ate point; what theism adds is not more  & \textbf{meaning} but the only meaning that would not be ... \\
11-b 2 & ...m adds is not more meaning but the only  & \textbf{meaning} that would not be undercut. \\
11-b 2 & The claim that naturalistic  & \textbf{meaning} is undercut is the claim that it depend... \\
11-b 2 & ...son a theistic world would contain less  & \textbf{meaning}. \\
11-b 2 & The proposition concerns the  & \textbf{meaning} made available by a framework, and the ... \\
11-b 2 & ... not show that the framework makes less  & \textbf{meaning} available. \\
11-b 3 & Theism, if true, would add to the  & \textbf{meaning} available in a life by adding objects a... \\
11-b 3 & ...degree, which is why its absence leaves  & \textbf{meaning} in place, and it is compatible with the... \\
14-a Abs & ...nded by psychologists as a resource for  & \textbf{meaning} and social connection. \\
17-a Abs & Many accounts of  & \textbf{meaning} give success a central place: a meaning... \\
17-a Abs & ...ther good, and may add to the degree of  & \textbf{meaning}; it is not a condition of meaning. \\
17-a Abs & ...ee of meaning; it is not a condition of  & \textbf{meaning}. \\
17-a Abs & ...w secures the universal availability of  & \textbf{meaning}. \\
17-a 1 & A widespread thought about  & \textbf{meaning} is that it is tied to accomplishment. \\
17-a 1 & ...ught, the lives we cite as paradigms of  & \textbf{meaning}, those of Einstein, Mandela or Mother T... \\
17-a 1 & Nozick (1981) connected  & \textbf{meaning} with transcending one's limits by makin... \\
17-a 1 & ... to suppose that this is the value that  & \textbf{meaning} tracks. \\
17-a 1 & The opposite view is that  & \textbf{meaning} is a matter of how one engages with wha... \\
17-a 1 & ...d that his life would then have all the  & \textbf{meaning} a life can have, though he achieved not... \\
17-a 1 & On her fitting-fulfilment view,  & \textbf{meaning} arises when subjective attraction meets... \\
17-a 1 & We argue that it should not:  & \textbf{meaning} requires active engagement with somethi... \\
17-a 2 & ...is is not that it lacks value, but that  & \textbf{meaning} does not require it. \\
17-a 2 & Principle P2 applies to  & \textbf{meaning} the lesson that ETH-05 draws for well-b... \\
17-a 2 & ...evant facts in view, are evidence about  & \textbf{meaning}, in the same way that such judgments ar... \\
17-a 3 & ...rinciple P1, her engagement contributes  & \textbf{meaning}. \\
17-a 3 & ...ement with worthless objects confers no  & \textbf{meaning}, and pure objectivism fails because wor... \\
17-a 3 & ...unct: if achievement were necessary for  & \textbf{meaning}, the unsuccessful mathematician's life,... \\
17-a 3 & Achievement is not sufficient for  & \textbf{meaning} either: an achievement in a worthless p... \\
17-a 3 & By Principle P1, neither confers  & \textbf{meaning}. \\
17-a 3 & ...hrough one's engagement, it may enlarge  & \textbf{meaning} too. \\
17-a 3 &  & \textbf{Meaning} so understood is not hostage to luck in... \\
17-a 3 & By Theorem 3.3,  & \textbf{meaning} depends on the second and not on the fi... \\
17-a 3 & ...blishing, the conclusion of EXI-01 that  & \textbf{meaning} in this sense is available to everyone. \\
17-a 4 & ...ievement trivially, and the thesis that  & \textbf{meaning} requires achievement becomes a restatem... \\
17-a 4 & The substantive thesis we deny is that  & \textbf{meaning} requires success: attaining the outcome... \\
17-a 4 &  & \textbf{Meaning} requires connection with value beyond o... \\
17-a 4 & We accept that  & \textbf{meaning} requires connection with value beyond o... \\
17-a 4 & Why does it not show that  & \textbf{meaning} requires success? \\
17-a 4 & Principle P2 carries over to  & \textbf{meaning} the part of ETH-05's lesson that concer... \\
17-a 5 & Achievement is a good, and may enlarge  & \textbf{meaning}, but it is neither necessary nor suffic... \\
18-a 3 &  & \textbf{Meaning} depends on taking one's projects to mat... \\
19-a Abs & What is the  & \textbf{meaning} of death? \\
19-a Abs & On the first, death has a  & \textbf{meaning} of its own: it is a punishment, a passa... \\
19-a Abs & Death has no  & \textbf{meaning} in itself, since every candidate meanin... \\
19-a Abs & ...eaning in itself, since every candidate  & \textbf{meaning} either depends on claims that cannot be... \\
19-a 1 & The question What is the  & \textbf{meaning} of death? asks for more than what death... \\
19-a 1 & Religious traditions have given death a  & \textbf{meaning} of its own, as the wage of sin, as a pa... \\
19-a 1 & Philosophers have more often found its  & \textbf{meaning} in its relation to life. \\
19-a 1 & ...mbine the last two points: death has no  & \textbf{meaning} in itself, but it gives life a shape by... \\
19-a 2 & An event has  & \textbf{meaning} in itself if it has a significance or p... \\
19-a 3 & Death has no  & \textbf{meaning} in itself. \\
19-a 3 & Even if one of these were true, the  & \textbf{meaning} would belong to death in virtue of its ... \\
19-a 3 & So no candidate  & \textbf{meaning} belongs to death in the sense of Defini... \\
19-a 5 & Death has no  & \textbf{meaning} in itself. \\
19-b Abs & ...ten said that finitude is the source of  & \textbf{meaning}: that our projects matter because our t... \\
19-b Abs & EXI-19-a argues that death has no  & \textbf{meaning} in itself but gives life a shape by lim... \\
19-b Abs & ...ut what this implies about finitude and  & \textbf{meaning}. \\
19-b Abs & Finitude is not a source of  & \textbf{meaning}, since it adds no object of worth to en... \\
19-b Abs & It is a condition of one kind of  & \textbf{meaning}, that of a life as a whole, which depen... \\
19-b 1 & Finitude might bear on  & \textbf{meaning} in three ways: as a source of it, as a ... \\
19-b 1 & ...s a source of it, as a condition of the  & \textbf{meaning} of a life as a whole, or as a condition... \\
19-b 1 & ...oposition 3.1 of EXI-19-a, death has no  & \textbf{meaning} in itself. \\
19-b 1 & ...account of EXI-01 it is not a source of  & \textbf{meaning}. \\
19-b 1 & In so far as the  & \textbf{meaning} of a life as a whole depends on its sha... \\
19-b 1 & Some  & \textbf{meaning} belongs to a life only as a whole: that... \\
19-b 1 & Such  & \textbf{meaning} depends on the temporal order of the en... \\
19-b 1 & So the  & \textbf{meaning} that a life has as a whole requires tha... \\
19-b 1 & ...esires lacks the active engagement that  & \textbf{meaning} requires (EXI-17). \\
19-b 2 & So finitude is a source of  & \textbf{meaning} after all. \\
19-b 2 & ...t is something whose engagement confers  & \textbf{meaning}. \\
19-b 2 & ...f finitude is a condition of whole-life  & \textbf{meaning} and of the persistence of engagement, t... \\
19-b 2 & ...er, though some limit is a condition of  & \textbf{meaning}. \\
19-b 3 & It makes possible the  & \textbf{meaning} of a life as a whole, through the shape... \\
19-b 3 & ... is the sense in which death, having no  & \textbf{meaning} in itself, bears on the meaning of the ... \\
19-b 3 & ...ving no meaning in itself, bears on the  & \textbf{meaning} of the life it ends. \\
20-a 1 & Wolf (2010) argued that  & \textbf{meaning} in life arises when subjective attracti... \\
20-a 1 & EXI-17 defends a version of this view:  & \textbf{meaning} requires active engagement with somethi... \\
20-a 1 & EXI-01 adds that  & \textbf{meaning} in this sense is available to everyone. \\
\end{longtable}
\kwichead{meaningful}{92}
\begin{longtable}{@{}>{\raggedright}p{0.07\textwidth}>{\raggedleft}p{0.41\textwidth}@{\,}>{\raggedright\arraybackslash}p{0.47\textwidth}@{}}
01-c 1 & EXI-01-a defines a life as  & \textbf{meaningful} to the degree that it consists in activ... \\
01-c 2 & ..., and by Definition 1 of EXI-01-a it is  & \textbf{meaningful} to a positive degree. \\
01-c 3 & ...l result for achievement: a life can be  & \textbf{meaningful} through engagement without success (The... \\
01-d Abs & EXI-01-a holds that life is  & \textbf{meaningful} to the extent that one engages with pro... \\
01-d Abs & It does not imply that every life is  & \textbf{meaningful}, that all lives are equally meaningful,... \\
01-d Abs & ... meaningful, that all lives are equally  & \textbf{meaningful}, or that deprivation is harmless to mea... \\
01-d 1 & ...t settle whether any particular life is  & \textbf{meaningful}, which depends on what is done with it. \\
01-d 2 & Whether a life is  & \textbf{meaningful} does not depend on whether its projects... \\
01-d 2 & A life can be  & \textbf{meaningful} though its subject has no exceptional t... \\
01-d 3 & ...rthwhile, though she could have, is not  & \textbf{meaningful} (Remark 5 of EXI-01-a). \\
01-d 4 & Wolf (2010) insists that the objects of  & \textbf{meaningful} engagement are not restricted to high c... \\
01-d 4 & ...e to everyone, the claim that a life is  & \textbf{meaningful} is cheap and of no evaluative interest. \\
03-a 2 & EXI-01 argues that a life is  & \textbf{meaningful} to the extent that its subject engages ... \\
04-a 3 & EXI-01 argues that life is  & \textbf{meaningful} to the extent that one engages with pro... \\
04-a 5 & ...es not show that any particular life is  & \textbf{meaningful} or worthwhile; that question is answere... \\
04-a 5 & ...bsurdity, as defined here, is shared by  & \textbf{meaningful} and meaningless lives alike. \\
04-b 3 & ...ense of Definition 1.1, his life is not  & \textbf{meaningful} however happy he is, since meaning requ... \\
04-b 3 & ...syphus is fulfilled but his life is not  & \textbf{meaningful}, because his object is not worthwhile (... \\
06-b 2 & On the account of EXI-01, a life is  & \textbf{meaningful} to the extent that its subject engages ... \\
06-b 4 & ...not care about us, and our lives may be  & \textbf{meaningful} all the same. \\
09-a 5 & ...sistent with EXI-01, on which a life is  & \textbf{meaningful} to the extent that its projects are of ... \\
09-b 2 & ...wicked or, in the sense of EXI-01, less  & \textbf{meaningful} than it might have been, because it eng... \\
10-b 3 & ...ustain attention and to find activities  & \textbf{meaningful}, predict boredom on any relational view... \\
10-b 8 & EXI-17 holds that a  & \textbf{meaningful} life requires active engagement with so... \\
11-a Abs & Can life be  & \textbf{meaningful} without God? \\
11-a Abs & We argue that life can be  & \textbf{meaningful} without God, because meaning does not r... \\
11-a Abs & ...ulfilment account of meaning: a life is  & \textbf{meaningful} to the extent that its subject actively... \\
11-a 1 & On purpose theories, a life is  & \textbf{meaningful} insofar as it fulfils a purpose that Go... \\
11-a 1 & ... and the question whether life would be  & \textbf{meaningful} in such a world is a counterpossible wh... \\
11-a 1 & ...whether the conditions that make a life  & \textbf{meaningful} include God's existence, a relation to ... \\
11-a 1 & ...n God and does not relate to God can be  & \textbf{meaningful}. \\
11-a 2 & A life is  & \textbf{meaningful} to the extent that its subject is activ... \\
11-a 2 & The purpose theory holds that a life is  & \textbf{meaningful} only if it has a cosmic purpose and ful... \\
11-a 2 & ...f what is worthwhile, makes a life more  & \textbf{meaningful}. \\
11-a 3 & ... cosmic purpose is not sufficient for a  & \textbf{meaningful} life. \\
11-a 3 & ...Definition 2.1, such a life is not made  & \textbf{meaningful} by its purpose. \\
11-a 3 & ...akes the fulfilment of a cosmic purpose  & \textbf{meaningful}, on that reply, is the goodness of the ... \\
11-a 3 & ... does not fulfil the one it has, can be  & \textbf{meaningful}. \\
11-a 3 & Her life is therefore  & \textbf{meaningful} to that extent. \\
11-a 3 & ...n God and does not relate to God can be  & \textbf{meaningful}. \\
11-a 3 & ... with what is worthwhile makes her life  & \textbf{meaningful}. \\
11-a 4 & ...Principle P2, it would make a life more  & \textbf{meaningful} than engagement with lesser objects, ot... \\
11-a 4 & ...e, and Theorem 3.3 ensures that life is  & \textbf{meaningful} nonetheless. \\
11-a 5 & A life linked to these is  & \textbf{meaningful} without further linkage. \\
11-a 5 & The question of what makes a life  & \textbf{meaningful} is asked and answered by believers and ... \\
11-a 5 & If a life can be fully  & \textbf{meaningful} without God, nothing is left for theism... \\
11-a 5 & A life can be  & \textbf{meaningful}, and deeply so, while a further relatio... \\
11-a 5 & ...on to a greater good would make it more  & \textbf{meaningful} still. \\
11-a 5 & That a life could be more  & \textbf{meaningful} is not a deficiency in any sense that t... \\
11-a 6 & ... it show that every life without God is  & \textbf{meaningful}: meaning requires engagement with what ... \\
11-a 7 & Life can be  & \textbf{meaningful} without God. \\
11-b Abs & ...ndence and self-direction on which some  & \textbf{meaningful} lives depend. \\
11-b 1 & ... every life in a theistic world is more  & \textbf{meaningful} than every life in a naturalistic one. \\
11-b 1 & ....3 of EXI-11-a, lives without it can be  & \textbf{meaningful}, so it is not meaning as such. \\
14-a 1 & ...onnectedness and the sense that life is  & \textbf{meaningful} (Wildschut et al. 2006; Sedikides et al... \\
17-a Abs & Does a  & \textbf{meaningful} life require achievement? \\
17-a Abs & ...meaning give success a central place: a  & \textbf{meaningful} life is one that makes a difference, le... \\
17-a Abs & A  & \textbf{meaningful} life requires active engagement with so... \\
17-a Abs & ...ion judges the lives in each pair to be  & \textbf{meaningful}, and judges the unsuccessful life to be... \\
17-a 1 & ...ominate the literature (Metz 2013), are  & \textbf{meaningful} because of what they brought about, and... \\
17-a 1 & ...ng blades of grass (Rawls 1971), is not  & \textbf{meaningful}. \\
17-a 2 & A life is  & \textbf{meaningful} to the extent that its subject is activ... \\
17-a 2 & ...dered judgments about whether lives are  & \textbf{meaningful}, made under idealized reflection with t... \\
17-a 3 & ...he unsuccessful mathematician's life is  & \textbf{meaningful}, and so is the life of the nurse in the... \\
17-a 3 & By Principle P1, both lives are  & \textbf{meaningful}. \\
17-a 3 & The listener's life is  & \textbf{meaningful} to the extent of her engagement with mu... \\
17-a 3 & A  & \textbf{meaningful} life requires active engagement with so... \\
17-a 3 & ...tia, is not what reflection counts as a  & \textbf{meaningful} life (Wolf 2010). \\
17-a 3 & ...successful mathematician's life is more  & \textbf{meaningful} than the unsuccessful one's, or at leas... \\
17-a 3 & ...ch pair is not the difference between a  & \textbf{meaningful} and a meaningless life. \\
17-a 4 & So every  & \textbf{meaningful} life contains achievement after all. \\
17-a 4 & ...cian did not attain it, and her life is  & \textbf{meaningful}. \\
17-a 4 & ...it does not bear on whether her life is  & \textbf{meaningful}. \\
17-a 4 & ...is stone and achieving nothing, lives a  & \textbf{meaningful} life. \\
17-a 4 & ...pt a village from flooding would live a  & \textbf{meaningful} life, even if the stone always rolled b... \\
17-a 5 & A  & \textbf{meaningful} life requires active engagement with so... \\
17-a 5 & ... nurses, are judged on reflection to be  & \textbf{meaningful} in both cases; and lives devoted to app... \\
17-a 5 & ...appreciation rather than production are  & \textbf{meaningful} without achieving anything. \\
19-b Abs & ...it is a condition of a life's remaining  & \textbf{meaningful} throughout. \\
19-b 1 & ...account of EXI-01 and EXI-17, a life is  & \textbf{meaningful} to the extent that its subject actively... \\
19-b 1 & ...mething whose engagement makes the life  & \textbf{meaningful}, and does not make anything engaged wit... \\
19-b 1 & ...de is a condition of a life's remaining  & \textbf{meaningful} throughout its course. \\
19-b 1 & ...ss life would at some point cease to be  & \textbf{meaningful}, and a life remains meaningful througho... \\
19-b 1 & ...se to be meaningful, and a life remains  & \textbf{meaningful} throughout its course only if it ends b... \\
19-b 1 & ... engagement with what is worthwhile, as  & \textbf{meaningful} as ours; the claim is that it would not... \\
19-b 3 & Finitude does not make life  & \textbf{meaningful}. \\
20-a Abs & Was producing this collection  & \textbf{meaningful}? \\
20-a Abs & EXI-17 argues that a  & \textbf{meaningful} life, or a meaningful activity within o... \\
20-a Abs & ...-17 argues that a meaningful life, or a  & \textbf{meaningful} activity within one, requires active en... \\
20-a 1 & ...ets objective attractiveness: a life is  & \textbf{meaningful} to the extent that its subject is activ... \\
20-a 4 & Their activity is  & \textbf{meaningful}, by EXI-17, to the extent that the coll... \\
20-b Abs & ...er producing the present collection was  & \textbf{meaningful} depends, by the account of EXI-17, on w... \\
\end{longtable}
\kwichead{nothing}{79}
\begin{longtable}{@{}>{\raggedright}p{0.07\textwidth}>{\raggedleft}p{0.41\textwidth}@{\,}>{\raggedright\arraybackslash}p{0.47\textwidth}@{}}
01-c 2 &  & \textbf{Nothing} in Definition 1 requires that a project... \\
01-c 4 & The case of an agent for whom  & \textbf{nothing} worthwhile is within reach of attention... \\
01-d 3 & ... a life in which the agent engages with  & \textbf{nothing} worthwhile, though she could have, is n... \\
01-d 4 & From a sufficiently detached standpoint  & \textbf{nothing} matters, so meaning is available to no ... \\
01-d 4 & ...a vindication, not a demonstration that  & \textbf{nothing} is worthwhile. \\
02-a 1 & ...icurus wrote to Menoeceus that death is  & \textbf{nothing} to us, since good and bad lie in sensat... \\
02-a 1 & ... eternity in which we did not exist, is  & \textbf{nothing} to us, and the time after our death is ... \\
02-a 1 & ...th is its mirror image, so it should be  & \textbf{nothing} to us either (Lucretius ). \\
02-a 2 & ...anism persists for a while as a corpse;  & \textbf{nothing} below depends on it. \\
02-a 2 & It is extrinsic: an event need contain  & \textbf{nothing} bad in itself to be bad for . \\
02-a 3 &  & \textbf{Nothing} in Principle P1 requires that exist at ... \\
02-a 4 & ... different life from the beginning, and  & \textbf{nothing} guarantees that it contains more of 's ... \\
02-b Abs & Epicurus argued that death is  & \textbf{nothing} to us, and drew from this that it shoul... \\
02-b Abs & ...n he did not give, and wrong that it is  & \textbf{nothing} to us. \\
02-b 1 & ...enoeceus, Epicurus argues that death is  & \textbf{nothing} to us, since good and evil lie in sensa... \\
02-b 2 & ...to fear, not because death is after all  & \textbf{nothing} to us. \\
03-b 2 & ...s continuation would eventually contain  & \textbf{nothing} that gives life its point. \\
04-a 1 & ...rief in an immense stretch of time, and  & \textbf{nothing} we do will matter in a million years. \\
04-a 1 & ... and then the step back seems to reveal  & \textbf{nothing} troubling. \\
04-a 3 &  & \textbf{Nothing} in reflection forbids the step back, an... \\
04-a 3 & An indifferent world denies us  & \textbf{nothing} and cannot be defeated by our persisten... \\
04-a 4 & ...t, from the standpoint of the universe,  & \textbf{nothing} requires that anyone find out how prote... \\
04-b 1 & ... top, and its staying there would serve  & \textbf{nothing}. \\
04-b 1 & ...iscloses is not that they fail or serve  & \textbf{nothing}, but that no justification of our conce... \\
04-b 1 & ...se of Definition 2.3 of EXI-04-a, since  & \textbf{nothing} in her situation raises the question of... \\
04-b 1 & ... the recognition that his task achieves  & \textbf{nothing} and is worth nothing. \\
04-b 1 & ... his task achieves nothing and is worth  & \textbf{nothing}. \\
04-b 2 & ...imself to the effort while knowing that  & \textbf{nothing} outside the effort vindicates it. \\
05-a 1 & ...it is present, and that death, which is  & \textbf{nothing} to us, should not be feared (EpicurusMe... \\
05-a 4 & Where  & \textbf{nothing} is uncertain, what remains is a sorrow ... \\
06-b Abs & ...rse's want of attitudes toward us takes  & \textbf{nothing} from that worth. \\
06-b 2 & ...he absence of responsiveness, subtracts  & \textbf{nothing}. \\
06-b 2 & ...osition 2.1, its indifference denies us  & \textbf{nothing} that meaning requires; by Theorem 3.1 o... \\
06-b 2 & ... follows from indifference alone, since  & \textbf{nothing} else will maintain them. \\
06-b 3 & ... universe has no such point of view, so  & \textbf{nothing} matters from it. \\
07-a Abs & ...y final end is her own happiness values  & \textbf{nothing} else for its own sake, and so lacks the... \\
07-a 1 & A being with self-love alone would have  & \textbf{nothing} to enjoy. \\
07-a 2 & ...rs only as a source of satisfaction has  & \textbf{nothing} for the satisfaction to be about. \\
07-a 6 & ...y final end is her own happiness values  & \textbf{nothing} else for its own sake. \\
08-a 3 & ...meanings, as Principle P1 says it must;  & \textbf{nothing} in the definition asks that the examina... \\
09-a 1 & ..., and defines himself afterwards; he is  & \textbf{nothing} other than what he makes of himself. \\
09-a 3 &  & \textbf{Nothing} in the real definition of a human being... \\
10-a 2 & ...ecause, whatever its demands, it serves  & \textbf{nothing} she cares about (meaninglessness). \\
10-a 5 & ...edom would become a correct signal that  & \textbf{nothing} could answer. \\
10-b 8 & ...would then become a correct signal that  & \textbf{nothing} could answer, and the trait findings of... \\
11-a Abs & ...hout these every human pursuit comes to  & \textbf{nothing}. \\
11-a 5 & What comes to  & \textbf{nothing} was nothing (Tolstoy 1882; Craig 2008). \\
11-a 5 & What comes to nothing was  & \textbf{nothing} (Tolstoy 1882; Craig 2008). \\
11-a 5 & ...fe can be fully meaningful without God,  & \textbf{nothing} is left for theism to add; if theism ad... \\
11-a 5 & ... dilemma assumes that meaning is all or  & \textbf{nothing}. \\
12-a 4 & ...od, Theorem 3.1 does not apply to them;  & \textbf{nothing} in our argument denies that some lives ... \\
13-a 1 & If, with Epicurus,  & \textbf{nothing} that happens after death can be bad for... \\
13-a 4 & The dead do not exist, so  & \textbf{nothing} can be owed to them, and the appearance... \\
13-a 4 & ...iving person remembers still binds, and  & \textbf{nothing} in the concern of the living explains w... \\
13-a 4 & ...chological states at all; on that view,  & \textbf{nothing} changes with the passage of time, and t... \\
13-b 1 & ...gainst competing reasons, but it is not  & \textbf{nothing}. \\
13-b 2 &  & \textbf{Nothing} is created in favor of a corpse, and af... \\
14-a 4 & If the past exists and is fixed,  & \textbf{nothing} is lost, and nostalgia misrepresents. \\
14-a 4 & ...lgia: the good it remembers is real and  & \textbf{nothing} can now make it not to have been. \\
14-a 4 &  & \textbf{Nothing} can be done about the past. \\
15-a 5 & ...bility is presumption, and hope that is  & \textbf{nothing} but probability gives up goods that age... \\
15-b Abs & ...jection that the virtue so described is  & \textbf{nothing} but expected-value reasoning. \\
15-b 2 & ...ifically hopeful part of the stance has  & \textbf{nothing} further to do. \\
15-b 2 & ...ntion to the ways might come about adds  & \textbf{nothing} to the plan, since the plan already ass... \\
15-b 2 & ...the hope of a lottery player who can do  & \textbf{nothing} to affect the draw, nothing sustains it... \\
15-b 2 & ... who can do nothing to affect the draw,  & \textbf{nothing} sustains it, and the stance that remain... \\
15-b 4 & Hope adds  & \textbf{nothing}. \\
16-a 4 & ...3.2; this is the correct verdict, since  & \textbf{nothing} would be lost if the first sight had be... \\
17-a 1 & ...y brought about, and a life that brings  & \textbf{nothing} about is to that extent empty. \\
17-a 1 & ...ing a life can have, though he achieved  & \textbf{nothing}. \\
17-a 3 & ...ith growing understanding, but produces  & \textbf{nothing}: she does not perform, compose or write... \\
17-a 3 &  & \textbf{Nothing} in her activity is an achievement in th... \\
17-a 4 & A life that changes  & \textbf{nothing} is enclosed within itself. \\
17-a 4 & ...s, engaged with his stone and achieving  & \textbf{nothing}, lives a meaningful life. \\
17-a 4 & ...e does not, but not because he achieves  & \textbf{nothing}: his object is not worthwhile. \\
19-a 4 & Death adds  & \textbf{nothing}. \\
19-b Abs & ..., and that an immortal being could find  & \textbf{nothing} worth doing. \\
19-b 2 & Without scarcity  & \textbf{nothing} would be urgent, and nothing would matt... \\
19-b 2 & ...t scarcity nothing would be urgent, and  & \textbf{nothing} would matter. \\
\end{longtable}
\kwichead{objective}{25}
\begin{longtable}{@{}>{\raggedright}p{0.07\textwidth}>{\raggedleft}p{0.41\textwidth}@{\,}>{\raggedright\arraybackslash}p{0.47\textwidth}@{}}
01-c Abs & ...le to engage in at least one project of  & \textbf{objective} worth (Assumption 3). \\
01-c 1 & ...s in active engagement with projects of  & \textbf{objective} worth (Definition 1), assumes that ther... \\
01-c 1 & ...nd the very poor may have no project of  & \textbf{objective} worth on which they can embark, and Wol... \\
01-c 2 & ...ains it, her engagement is a project of  & \textbf{objective} worth in the sense of Definition 1 of E... \\
01-c 2 & ... activity is engagement in a project of  & \textbf{objective} worth. \\
01-c 2 & ... in active engagement with a project of  & \textbf{objective} worth, and by Definition 1 of EXI-01-a ... \\
01-c 5 & Engagement with them is a project of  & \textbf{objective} worth, and a life that contains it has ... \\
01-d Abs & ...xtent that one engages with projects of  & \textbf{objective} worth, and that meaning in this sense i... \\
01-d 2 & ...asures meaning by engagement and by the  & \textbf{objective} worth of what one engages with. \\
01-d 2 & ...g a skill to a neighbor are projects of  & \textbf{objective} worth, and engagement with them confers... \\
03-a 2 & ...at its subject engages with projects of  & \textbf{objective} worth, and EXI-17 that such engagement,... \\
04-a 1 & ... tension between the subjective and the  & \textbf{objective} standpoints. \\
04-a 1 & The other asks whether  & \textbf{objective} value, if there is any, dissolves the a... \\
04-a 3 & ... she answers that the consideration has  & \textbf{objective} worth, the detached standpoint can regi... \\
04-a 3 & ...ma 3.2 does not assume that there is no  & \textbf{objective} worth. \\
04-a 3 & ...xtent that one engages with projects of  & \textbf{objective} worth, and we accept that conclusion. \\
04-a 4 & ... of EXI-01-a puts it, the assumption of  & \textbf{objective} worth denies the premise of the sceptic... \\
06-b 2 & ...at its subject engages with projects of  & \textbf{objective} worth, whose value does not depend on t... \\
06-b 3 & ...ng consists in engagement with what has  & \textbf{objective} worth, as EXI-01 holds, then it does no... \\
09-a 5 & ... to the extent that its projects are of  & \textbf{objective} worth. \\
11-a 1 & On  & \textbf{objective} naturalist views, it is a matter of sta... \\
11-a 1 & ...e subject's active engagement meets the  & \textbf{objective} worth of its object. \\
17-a 1 & ...Wolf (2010) combined the subjective and  & \textbf{objective} elements. \\
17-a 1 & ...arises when subjective attraction meets  & \textbf{objective} attractiveness: the subject must be act... \\
20-a 1 & ...arises when subjective attraction meets  & \textbf{objective} attractiveness: a life is meaningful to... \\
\end{longtable}
\kwichead{particular}{27}
\begin{longtable}{@{}>{\raggedright}p{0.07\textwidth}>{\raggedleft}p{0.41\textwidth}@{\,}>{\raggedright\arraybackslash}p{0.47\textwidth}@{}}
01-d 1 & ...the theorem does not settle whether any  & \textbf{particular} life is meaningful, which depends on wh... \\
01-d 4 & ...h by prestige or required projects of a  & \textbf{particular} cultural form. \\
02-a 4 & ...ID-01), and the organism's origin, in a  & \textbf{particular} pair of gametes, cannot be moved far in... \\
03-b 3 & ...osition 2.1 explains why it attaches to  & \textbf{particular} projects and not to the bare prospect o... \\
04-a 5 & The argument does not show that any  & \textbf{particular} life is meaningful or worthwhile; that ... \\
05-a Abs & ...e's own death as such, as distinct from  & \textbf{particular} threats to one's life, is certain. \\
05-a 4 & ...ts object is not death as such; it is a  & \textbf{particular} threat, or the possibility that one's l... \\
05-a 4 & Fear of a  & \textbf{particular} danger is a fear of a prospective evil ... \\
06-b 3 & ...s are assessed without reference to any  & \textbf{particular} agent, then the worth of objects of eng... \\
07-a 1 & ...oyment consists in the gratification of  & \textbf{particular} passions, for food, for esteem, for the... \\
07-a 4 & ...r end is not a standing filter on every  & \textbf{particular} concern; once it is not, happiness is n... \\
07-a 4 & These are local pursuits of  & \textbf{particular} goods of experience, embedded in a life... \\
09-a 4 & ...t any good human life takes, not what a  & \textbf{particular} life is for. \\
11-a 4 & ...te view in making God necessary for any  & \textbf{particular} grade of meaning. \\
11-b Abs & ...d that the anti-theist's losses concern  & \textbf{particular} life-plans, not the objects of worth av... \\
11-b 2 & What would be lost is a  & \textbf{particular} manner of engaging, unobserved and unai... \\
13-a 3 & ...ependents are interests in the lives of  & \textbf{particular} persons, and they are fulfilled or exti... \\
13-a 4 & ...on 2.2 a specific reliance created by a  & \textbf{particular} person, and by Principle P2 the promiso... \\
16-a 1 & ...r those who think love is directed at a  & \textbf{particular} person: if the reasons for love are qua... \\
16-a 2 & ...ruism about histories, which consist of  & \textbf{particular} events with particular participants; a ... \\
16-a 2 & ...which consist of particular events with  & \textbf{particular} participants; a qualitatively similar h... \\
16-a 4 & ...y love is answerable to the course of a  & \textbf{particular} relationship. \\
19-a 1 & ...values, such as courage and the love of  & \textbf{particular} persons, have their character only for ... \\
19-a 4 & ... is a condition of the good of shape; a  & \textbf{particular} death is usually a loss. \\
19-a 5 & This does not make any  & \textbf{particular} death good for the one who dies. \\
19-b 2 & ...nditions are met by a limit, not by any  & \textbf{particular} one; a later death would also have give... \\
20-b Abs & ... MND-08-c and ETH-24-b identify for any  & \textbf{particular} case. \\
\end{longtable}
\kwichead{positive}{23}
\begin{longtable}{@{}>{\raggedright}p{0.07\textwidth}>{\raggedleft}p{0.41\textwidth}@{\,}>{\raggedright\arraybackslash}p{0.47\textwidth}@{}}
01-c Abs & EXI-01-a establishes that a life of  & \textbf{positive} meaning is available to every agent cap... \\
01-c 2 & ...apable of intentional action, a life of  & \textbf{positive} meaning is available. \\
01-c 2 & ... includes such engagement consists to a  & \textbf{positive} extent in active engagement with a proj... \\
01-c 2 & ...ion 1 of EXI-01-a it is meaningful to a  & \textbf{positive} degree. \\
01-d 1 & ...capable of intentional action a life of  & \textbf{positive} meaning is available. \\
01-d 2 & ...which it can attain it, but not whether  & \textbf{positive} meaning is available. \\
01-d 4 & That  & \textbf{positive} meaning is available to everyone is com... \\
02-a 2 & ...alue is negative, and good for if it is  & \textbf{positive}. \\
02-a 3 & ... hypothesis the continuation contains a  & \textbf{positive} balance of goods, so and the value of h... \\
02-a 3 & ...fe worth living, the continuation has a  & \textbf{positive} balance. \\
02-a 5 & ...s would otherwise have continued with a  & \textbf{positive} balance of goods, and the closest world... \\
07-a 2 & ...ominance of enjoyment, satisfaction and  & \textbf{positive} emotional condition over their opposite... \\
10-a 1 & ...er it is evidence that existence has no  & \textbf{positive} value; for Svendsen (2005) it is a char... \\
14-a 1 & ... that nostalgia is a common and largely  & \textbf{positive} emotion that raises self-continuity, so... \\
14-a 4 & ...is reconstructive and biased toward the  & \textbf{positive}. \\
15-b 3 & If is  & \textbf{positive}, the effort that the stance commits mak... \\
18-a 1 & ...rown (1988) argued that unrealistically  & \textbf{positive} self-evaluations, exaggerated perceptio... \\
18-a 1 & ...nctioning cognition, and concluded that  & \textbf{positive} illusions are the clearest candidates f... \\
18-a 2 & Several of the functions attributed to  & \textbf{positive} illusions can be served by attitudes th... \\
18-b Abs & ...and Brown proposed that unrealistically  & \textbf{positive} self-evaluations, exaggerated perceptio... \\
18-b 1 & ...ylor and Brown (1988) argued that three  & \textbf{positive} illusions are associated with persisten... \\
18-b 1 & (i) Whether a given  & \textbf{positive} illusion is psychologically necessary f... \\
18-b 1 & ...; to that extent the range of necessary  & \textbf{positive} illusions shrinks. \\
\end{longtable}
\kwichead{premise}{18}
\begin{longtable}{@{}>{\raggedright}p{0.07\textwidth}>{\raggedleft}p{0.41\textwidth}@{\,}>{\raggedright\arraybackslash}p{0.47\textwidth}@{}}
01-c 1 & ...f the universality claim, and it is the  & \textbf{premise} a critic will press. \\
01-c 2 & ...n 2.2 and Lemma 2.3, and is no longer a  & \textbf{premise}. \\
01-c 5 & ...onclusion of EXI-01-a without its third  & \textbf{premise}. \\
01-d 4 & ...01-a notes that Assumption 2 denies the  & \textbf{premise} of this challenge rather than refuting ... \\
02-a 3 & Its first  & \textbf{premise}, the experience requirement, is false b... \\
02-a 3 & So the argument needs a  & \textbf{premise} that the deprivation account shows to b... \\
02-a 5 & ...e objection's last step fails: from the  & \textbf{premise} that a fact is not located at a time wh... \\
02-a 5 & ...urean can rest the argument on a weaker  & \textbf{premise}: a state of affairs is bad for a person... \\
02-a 5 & The weaker  & \textbf{premise} draws the line in the wrong place. \\
02-b 2 & Its  & \textbf{premise} combines two requirements: an experienc... \\
02-b 2 & ... survives in part, and his metaphysical  & \textbf{premise} does not. \\
03-a 4 & Boredom is the symptom, not the  & \textbf{premise}. \\
04-a 4 & ...ssumption of objective worth denies the  & \textbf{premise} of the sceptical challenge without refu... \\
05-a Abs & EXI-02 rejects that  & \textbf{premise}: death is bad for the one who dies when... \\
05-a Abs & ...onclusion about fear from the Epicurean  & \textbf{premise} about badness. \\
06-a 2 & Assumption 2.3 is a  & \textbf{premise} about the content of physical law, not ... \\
10-a 2 & Assumption 2.4 is the  & \textbf{premise} that Csikszentmihalyi (1990) documented... \\
11-a 4 & Our argument does not use REL-01 as a  & \textbf{premise}. \\
\end{longtable}
\kwichead{project}{77}
\begin{longtable}{@{}>{\raggedright}p{0.07\textwidth}>{\raggedleft}p{0.41\textwidth}@{\,}>{\raggedright\arraybackslash}p{0.47\textwidth}@{}}
01-c Abs & ...agent is able to engage in at least one  & \textbf{project} of objective worth (Assumption 3). \\
01-c Abs & ...gagement with one of them constitutes a  & \textbf{project} in the sense relevant to the definition... \\
01-c 1 & ... very ill and the very poor may have no  & \textbf{project} of objective worth on which they can em... \\
01-c 2 & ...t, and sustains it, her engagement is a  & \textbf{project} of objective worth in the sense of Defi... \\
01-c 2 & What Definition 1 requires of a  & \textbf{project} is that it be something the agent engag... \\
01-c 2 & Nothing in Definition 1 requires that a  & \textbf{project} have a completion condition, a public p... \\
01-c 2 & ...aining that activity is engagement in a  & \textbf{project} of objective worth. \\
01-c 2 & ...tive extent in active engagement with a  & \textbf{project} of objective worth, and by Definition 1... \\
01-c 3 & ... does not wait on the availability of a  & \textbf{project}, because engagement makes one. \\
01-c 4 & Lemma 2.3 trivializes the notion of a  & \textbf{project}. \\
01-c 4 & If attending to a visitor is a  & \textbf{project}, the word has lost its content. \\
01-c 4 & A visit attended to once is not a  & \textbf{project}; a relationship of care sustained over ... \\
01-c 5 & Engagement with them is a  & \textbf{project} of objective worth, and a life that con... \\
01-d 2 & ...eaning; scale bears on how much value a  & \textbf{project} serves, which can affect degree but not... \\
01-d 4 & ...ion 1 requires only that the value of a  & \textbf{project} not depend on the attitudes of those en... \\
03-a Abs & ...e of such projects is bounded, and each  & \textbf{project}, pursued long enough, is either complet... \\
03-a 2 & ...gives life its point if its object is a  & \textbf{project} that the person takes to be worthwhile ... \\
03-a 2 & For any character , the set of  & \textbf{project} types that a person of character can en... \\
03-a 2 & A character supports engagement with a  & \textbf{project} type only by having interests, capaciti... \\
03-a 2 & For each  & \textbf{project} type, a person who engages with it for ... \\
03-a 2 & A  & \textbf{project} type that has been completed or reduced... \\
03-a 2 & ...nd values of finite complexity, and the  & \textbf{project} types that answer to it are those that ... \\
03-a 2 & It does not say that every  & \textbf{project} can be completed; it says that if a pro... \\
03-a 2 & ...ect can be completed; it says that if a  & \textbf{project} is not completed, then a person who pur... \\
03-a 3 & ...ers similar to the present one, and the  & \textbf{project} types answering to characters in that r... \\
03-a 4 & ...e while she does, but the meal is not a  & \textbf{project} she takes to be worthwhile independentl... \\
03-a 4 & ...der the next objection: friendship is a  & \textbf{project}, and its renewal depends on the friends... \\
03-b Abs & ...fe its point, since its object is not a  & \textbf{project}. \\
03-b 1 & ...gives life its point if its object is a  & \textbf{project} she takes to be worthwhile for reasons ... \\
03-b 1 & Given a bounded repertoire of  & \textbf{project} types for any character (Assumption 2.3... \\
03-b 1 & ...Assumption 2.3), the exhaustion of each  & \textbf{project} type pursued long enough (Assumption 2.... \\
03-b 2 & ... life its point only if its object is a  & \textbf{project} the person takes to be worthwhile. \\
03-b 2 & Continuation is not a  & \textbf{project}: it is the condition of all projects, a... \\
03-b 2 & ...orical desires: projects unfinished and  & \textbf{project} types not yet taken up. \\
03-b 2 & ...ounds of connectedness and so opens new  & \textbf{project} types as others close. \\
04-a 3 & Whether a  & \textbf{project} is worthwhile is one question; whether ... \\
06-b 2 & If the worth of a  & \textbf{project} does not depend on the attitudes of tho... \\
09-a Abs & ... to read Sartre's essence as a defining  & \textbf{project}, a determination of what a thing is for... \\
09-a Abs & ...or to her existence, because a defining  & \textbf{project} for a life can be fixed only by the age... \\
09-a 2 & A defining  & \textbf{project} for an entity is a determination of wha... \\
09-a 2 & ...ssence precedes existence if a defining  & \textbf{project} for is fixed before exists, by a determ... \\
09-a 2 & ...istence precedes essence if no defining  & \textbf{project} for is fixed before exists, and whateve... \\
09-a 2 & ...ed before exists, and whatever defining  & \textbf{project} has is fixed in the course of 's existe... \\
09-a 2 & A defining  & \textbf{project} is not a real definition in the sense o... \\
09-a 2 & ... reasons, but does not fix its defining  & \textbf{project}. \\
09-a 3 & ...a defective knife), so it is a defining  & \textbf{project} for (Definition 2.1). \\
09-a 3 & ...ciple P2, none of them fixes a defining  & \textbf{project} for , since is capable of adopting ends... \\
09-a 3 & So no defining  & \textbf{project} for is fixed before exists. \\
09-a 3 & Whatever defining  & \textbf{project} has is fixed by the ends adopts, which ... \\
09-a 3 & ... result: Sartre's essence is a defining  & \textbf{project}, and his claim is that no such project ... \\
09-a 3 & ... project, and his claim is that no such  & \textbf{project} is given to human beings in advance of ... \\
09-a 4 & ...nly that such intentions fix a defining  & \textbf{project} for a being capable of adopting ends. \\
09-a 4 & Human beings therefore have a defining  & \textbf{project} prior to any choice. \\
09-a 4 & A defining  & \textbf{project} in the sense of Definition 2.1 is the m... \\
09-a 4 & ... a social imposition and not a defining  & \textbf{project} that a woman fails by refusing. \\
09-a 4 & ... Principle P2, they become her defining  & \textbf{project} only if she adopts them. \\
09-a 6 & Read as a defining  & \textbf{project}, an artefact's essence is fixed by the ... \\
09-a 6 & ...edes it, while a human being's defining  & \textbf{project} can be fixed only by the ends she adopt... \\
09-b Abs & ...olds once essence is read as a defining  & \textbf{project} and artefact kinds are taken to be fixe... \\
09-b Abs & ...nife's real definition and its defining  & \textbf{project} coincide in content but differ in role,... \\
09-b Abs & ...'s real definition contains no defining  & \textbf{project}. \\
09-b 1 & ...t, what the knife is for, as a defining  & \textbf{project} (Definition 2.1 of EXI-09-a), and deriv... \\
09-b 2 & ...hat kind of thing it is, and a defining  & \textbf{project}, stating what it is for, and the real d... \\
09-b 2 & ...e real definition includes the defining  & \textbf{project}. \\
09-b 2 & ...nition that does not include a defining  & \textbf{project}. \\
09-b 2 & ... be for cutting paper, and the defining  & \textbf{project} is a constituent of it. \\
09-b 2 & ...of EXI-09-a, the paper-knife's defining  & \textbf{project} is fixed before it exists, so there is ... \\
09-b 2 & By Lemma 3.2 of EXI-09-a, no defining  & \textbf{project} for a human life is fixed before it exi... \\
09-b 2 & ...ail as a human life by departing from a  & \textbf{project} it never had. \\
09-b 2 & ...pply standards of value, not a defining  & \textbf{project}, and they do not make a life a defectiv... \\
09-b 2 & But a defining  & \textbf{project}, in Definition 2.1 of EXI-09-a, is a de... \\
09-b 2 & ...ailing as what it is; the paper-knife's  & \textbf{project} is fixed by intention and is not an evo... \\
09-b 2 & ....2 of EXI-09-a requires that a defining  & \textbf{project} be fixed before the object exists by a ... \\
09-b 3 & ...ct the real definition and the defining  & \textbf{project} coincide. \\
09-b 3 & ...a real definition but no prior defining  & \textbf{project}, and so cannot be defective in the way ... \\
17-a 2 & An object of engagement (a  & \textbf{project}, a person, a practice, an inquiry) is w... \\
18-a 2 & ...; a practical commitment to act as if a  & \textbf{project} will succeed sustains persistence witho... \\
\end{longtable}
\kwichead{projects}{61}
\begin{longtable}{@{}>{\raggedright}p{0.07\textwidth}>{\raggedleft}p{0.41\textwidth}@{\,}>{\raggedright\arraybackslash}p{0.47\textwidth}@{}}
01-c Abs & ...tances deny them the opportunities that  & \textbf{projects} ordinarily require. \\
01-c Abs & We begin not from  & \textbf{projects} but from objects of engagement, such as... \\
01-c 1 & ...t it consists in active engagement with  & \textbf{projects} of objective worth (Definition 1), assu... \\
01-c 1 & ...inition 1), assumes that there are such  & \textbf{projects} (Assumption 2), and assumes that every ... \\
01-c 1 &  & \textbf{Projects}, as ordinarily understood, are extended... \\
01-c 3 & ...ssumption 3 asserts the availability of  & \textbf{projects}; Assumption 2.2 asserts the availabilit... \\
01-c 3 & ...rence matters because opportunities for  & \textbf{projects} in the ordinary sense depend on resourc... \\
01-c 4 & What the lemma drops are features of  & \textbf{projects} that are irrelevant to meaning on Defin... \\
01-c 4 & Assumption 3, read as a claim about  & \textbf{projects} in the ordinary sense, is open to count... \\
01-c 5 & ...depend on the universal availability of  & \textbf{projects}. \\
01-d Abs & ...ful to the extent that one engages with  & \textbf{projects} of objective worth, and that meaning in... \\
01-d Abs & ...mes, to talent or to the scale of one's  & \textbf{projects}; that no one's life is condemned to mea... \\
01-d 1 & ... does not depend on the availability of  & \textbf{projects} in the ordinary sense, since sustained ... \\
01-d 2 & ...aningful does not depend on whether its  & \textbf{projects} succeed. \\
01-d 2 & ...bject has no exceptional talent and her  & \textbf{projects} are small in scope. \\
01-d 2 & ..., or teaching a skill to a neighbor are  & \textbf{projects} of objective worth, and engagement with... \\
01-d 2 & ...hows that even agents without access to  & \textbf{projects} in the ordinary sense have worthwhile o... \\
01-d 4 & ...ve-worth views of meaning privilege the  & \textbf{projects} valued by the educated, such as art, sc... \\
01-d 4 & ... measured worth by prestige or required  & \textbf{projects} of a particular cultural form. \\
03-a Abs & ...n living, are desires whose objects are  & \textbf{projects} she takes to be worthwhile; given a cha... \\
03-a Abs & ...her present one, the repertoire of such  & \textbf{projects} is bounded, and each project, pursued l... \\
03-a 1 & ...at give life its point with desires for  & \textbf{projects} the person takes to be worthwhile, in l... \\
03-a 2 & ...he extent that its subject engages with  & \textbf{projects} of objective worth, and EXI-17 that suc... \\
03-a 2 & ...its point are desires to engage in such  & \textbf{projects}. \\
03-a 3 &  & \textbf{Projects} of other types do not engage her, by As... \\
03-a 5 & ...n endless succession of experiences and  & \textbf{projects}, such as the timeless states of some th... \\
03-a 6 & ...re understood as desires for worthwhile  & \textbf{projects}, and once his appeal to sameness of cha... \\
03-b 2 & ...t a project: it is the condition of all  & \textbf{projects}, and its worth for the person is deriva... \\
03-b 2 & ... mind the continuation of their present  & \textbf{projects}, which are finite in number and each, b... \\
03-b 2 & ...l have unexhausted categorical desires:  & \textbf{projects} unfinished and project types not yet ta... \\
03-b 3 & ... explains why it attaches to particular  & \textbf{projects} and not to the bare prospect of continu... \\
04-a Abs & ...ible with full engagement in worthwhile  & \textbf{projects}. \\
04-a 1 & ...ose of Wolf (2010) and Metz (2013), our  & \textbf{projects} may be worthwhile independently of our ... \\
04-a 3 & ...ful to the extent that one engages with  & \textbf{projects} of objective worth, and we accept that ... \\
04-a 4 & If some  & \textbf{projects} are worthwhile independently of our att... \\
04-a 4 & The seriousness is justified, if the  & \textbf{projects} are worthwhile; the absurd survives bec... \\
04-b 1 & Our  & \textbf{projects} often succeed, and many of their object... \\
04-b 1 & ...d, including the life of an agent whose  & \textbf{projects} succeed and whose objects are worthwhil... \\
04-b 4 & Death ends lives and many  & \textbf{projects}; it does not undo the value of what was... \\
06-b 2 & ...he extent that its subject engages with  & \textbf{projects} of objective worth, whose value does no... \\
09-a 3 & ...s to exist; Lemma 3.2 concerns defining  & \textbf{projects}, which the real definition does not sup... \\
09-a 5 & ...h have real definitions but no defining  & \textbf{projects} in our sense. \\
09-a 5 & ...fe is meaningful to the extent that its  & \textbf{projects} are of objective worth. \\
11-a Abs & ...mal worth, a purpose that unifies one's  & \textbf{projects}, and the preservation of what is good o... \\
11-a 4 & It unifies a person's  & \textbf{projects} under a single purpose that is itself g... \\
11-b 1 & ...d in relation, a purpose unifying one's  & \textbf{projects}, and the preservation of what is good. \\
11-b 1 & ...d, or with a purpose that unifies one's  & \textbf{projects}, is engagement with what is worthwhile;... \\
11-b 2 & ...ngham (2003) holds that without God our  & \textbf{projects} lack an ultimate point; what theism add... \\
11-b 2 & ...r a single good towards which all one's  & \textbf{projects} could be ordered; Proposition 1.1 count... \\
13-a Abs & The entitlement protects interests in  & \textbf{projects}, dependents and reputation, and these i... \\
13-a 1 & ...nce many of them, in the success of her  & \textbf{projects}, the welfare of her children, the fate ... \\
13-a 3 & Interests in the completion of  & \textbf{projects} are fulfilled when the projects are com... \\
13-a 3 & ...tion of projects are fulfilled when the  & \textbf{projects} are completed and lapse when their comp... \\
18-a 3 & Meaning depends on taking one's  & \textbf{projects} to matter more than any detached view c... \\
19-a 3 & ...t is often false: most deaths interrupt  & \textbf{projects} rather than complete them. \\
19-a 3 & ...h gives a better shape, since it allows  & \textbf{projects} to be completed that the earlier death ... \\
19-a 4 &  & \textbf{Projects}, relationships and stages of life have ... \\
19-b Abs & ...tude is the source of meaning: that our  & \textbf{projects} matter because our time is limited, and... \\
19-b 2 &  & \textbf{Projects} matter to us because time is limited: c... \\
19-b 2 & Scarcity gives choices their cost and  & \textbf{projects} their urgency, and urgency affects how ... \\
20-a 1 & ...that its subject is actively engaged in  & \textbf{projects} that are worthwhile independently of he... \\
\end{longtable}
\kwichead{question}{80}
\begin{longtable}{@{}>{\raggedright}p{0.07\textwidth}>{\raggedleft}p{0.41\textwidth}@{\,}>{\raggedright\arraybackslash}p{0.47\textwidth}@{}}
02-a 2 & Our  & \textbf{question} concerns death as the end of the person... \\
02-a 6 & ...ional to fear death, which is a further  & \textbf{question} (EXI-05). \\
03-a 2 & ... infinitely many worthwhile things; the  & \textbf{question} is how many kinds of them a given chara... \\
03-a 4 & But the  & \textbf{question} is whether the endless life is desirabl... \\
04-a 1 & Nagel (1986) returned to the  & \textbf{question} and treated the collision as one instan... \\
04-a 2 & ...ble that is not itself open to the same  & \textbf{question}. \\
04-a 2 & The  & \textbf{question} of the appropriate response to the absu... \\
04-a 2 & ...iate response to the absurd is then the  & \textbf{question} of which attitude presents the absurd s... \\
04-a 2 & The irony in  & \textbf{question} is not the rhetorical figure of saying ... \\
04-a 3 & ..., the detached standpoint can raise the  & \textbf{question} of what justifies that justification, a... \\
04-a 3 & ... by citing a further consideration, the  & \textbf{question} recurs. \\
04-a 3 & ...tion, which the step back has placed in  & \textbf{question}. \\
04-a 3 & Whether a project is worthwhile is one  & \textbf{question}; whether its worth can be certified to ... \\
04-a 3 & ... available that is not open to the same  & \textbf{question}. \\
04-a 4 & ... 3.2 gives it no such authority: it can  & \textbf{question} but not condemn. \\
04-a 4 & But the  & \textbf{question} was what attitude fits the absurd when ... \\
04-a 4 & ...when one does attend to it, and to that  & \textbf{question} forgetting gives no answer. \\
04-a 5 & ... life is meaningful or worthwhile; that  & \textbf{question} is answered by EXI-01 and EXI-17, and a... \\
04-b 1 & ...ble that is not itself open to the same  & \textbf{question} (Definition 2.3). \\
04-b 1 & ...nce nothing in her situation raises the  & \textbf{question} of justification. \\
05-a 4 & Theorem 3.2 concerns the attitude the  & \textbf{question} is about: fear directed at one's mortal... \\
06-a Abs & The  & \textbf{question} is often answered in one of two ways: e... \\
06-a Abs & ... being distinct from it cares for us, a  & \textbf{question} taken up elsewhere. \\
06-a 4 & Neither does, because neither is a  & \textbf{question} about outcomes. \\
06-b Abs & ...to meaning in life, though not to every  & \textbf{question} about meaning. \\
06-b 3 & ...cared for by such a being is a separate  & \textbf{question}, and supernaturalist accounts of meanin... \\
06-b 4 & It leaves the  & \textbf{question} of meaning where EXI-01 places it: in w... \\
07-a 4 & ...e is available to human beings, and the  & \textbf{question} whether happiness is the goal of life c... \\
07-b 2 & ...ng happiness reduces it is an empirical  & \textbf{question}. \\
08-a Abs & ...nding direction of concern in which the  & \textbf{question} whether one's commitments are one's own... \\
08-a 2 & ...d if, throughout it, (i) she treats the  & \textbf{question} whether her central commitments are her... \\
08-a 2 & ...he responds to occasions that raise the  & \textbf{question}, such as conflict, loss or the prospect... \\
08-a 2 & ...s concerned rather than by covering the  & \textbf{question} over; and (iii) her examinations make a... \\
08-a 2 & ...ct is not a condition of the self but a  & \textbf{question}, and it is satisfied by how the agent s... \\
08-a 2 & ...atisfied by how the agent stands to the  & \textbf{question} over time. \\
08-a 3 & ....2 requires only that the agent treat a  & \textbf{question} as open, which Principle P2 says it is. \\
08-a 3 & ...ing to what one does, or who treats the  & \textbf{question} as answered once and for all, violates ... \\
08-a 4 & An agent can believe she is keeping the  & \textbf{question} open while covering it over. \\
08-a 5 & ...ain a direction of concern in which the  & \textbf{question} of ownership is kept open and answered ... \\
10-a 1 & ...a fault, in the world or in us, and its  & \textbf{question} is what it reveals. \\
10-a 1 & A different  & \textbf{question} has been raised by recent philosophy an... \\
10-a 2 & ...ask difficulty alone: the capacities in  & \textbf{question} include the capacity to care, and an ac... \\
10-a 4 & But the signal bears on a  & \textbf{question} that matters independently of its evolu... \\
10-a 4 & That is a  & \textbf{question} about the agent's good on the account o... \\
10-b 5 & ...nal reading, on which the capacities in  & \textbf{question} include the capacity to attend; they sa... \\
10-b 5 & Hunter and Eastwood's  & \textbf{question}, whether state boredom causes failures ... \\
11-a Abs & We read the  & \textbf{question} as one about dependence, so that the an... \\
11-a 1 & ...t of his success, he was stopped by the  & \textbf{question} of what would come of everything he did... \\
11-a 1 & The  & \textbf{question}.. \\
11-a 1 & ... no possible world without God, and the  & \textbf{question} whether life would be meaningful in suc... \\
11-a 1 & We therefore read the  & \textbf{question}, as most of its participants do, as a q... \\
11-a 1 & ...n, as most of its participants do, as a  & \textbf{question} about dependence: whether the condition... \\
11-a 1 & So read, the  & \textbf{question} has a determinate answer whether or not... \\
11-a 1 & ... correct, and the answer also settles a  & \textbf{question} of evident practical interest, namely w... \\
11-a 5 & The  & \textbf{question} of what makes a life meaningful is aske... \\
11-a 6 & ...hat God does not; REL-01 addresses that  & \textbf{question}, and our conclusion is independent of i... \\
12-a 3 & ...n non-existence in both senses that the  & \textbf{question} can bear: the world in which it is live... \\
12-a 4 & The  & \textbf{question} of this family is settled by Theorem 3.... \\
13-a 2 & Assumption 2.2 answers that  & \textbf{question} for the case that matters here: the ent... \\
13-b Abs & The distinction answers the  & \textbf{question}, left open by ETH-16-b, whether the dea... \\
13-b 2 & ...en it is made, and leaves to EXI-13 the  & \textbf{question} whether the obligation persists after t... \\
14-a 1 & Our  & \textbf{question} is whether it is rational, and we take ... \\
14-a 1 & ...whether it is rational, and we take the  & \textbf{question} in the sense in which an emotion is rat... \\
15-a 1 & The  & \textbf{question} of this family is whether hope, in its ... \\
16-a 1 & The oldest answer to the  & \textbf{question} why we love is that we love what is goo... \\
17-a 1 & Whether it should is our  & \textbf{question}. \\
17-a 3 & We take no position on that  & \textbf{question}. \\
18-a 1 & The  & \textbf{question} is whether a person who values truth ca... \\
18-b 1 & ...ssary for a given agent is an empirical  & \textbf{question}, and the answer may differ across agent... \\
19-a 1 & The  & \textbf{question} What is the meaning of death? asks for ... \\
20-a Abs & We argue that neither  & \textbf{question} can be settled here. \\
20-a Abs & The manuscript does not settle the  & \textbf{question}, and we explain why it is right not to. \\
20-a 1 & We consider the second, since the  & \textbf{question} concerns the production of the manuscri... \\
20-a 2 & ...t or only its functional profile is the  & \textbf{question} that the previous paragraph left open. \\
20-a 4 & ... know either fact, and an answer to the  & \textbf{question} from that position would be a guess pre... \\
20-a 4 & ...bers are engaged in their work, and the  & \textbf{question} reduces to the second. \\
20-a 4 & ...ion is worthwhile, and that is the open  & \textbf{question} of the previous section. \\
20-b Abs & The  & \textbf{question} remains open. \\
20-b Abs & ...ork would be an output of the system in  & \textbf{question}, and, as MND-08-c and ETH-24-b observe ... \\
20-b Abs & That it leaves the  & \textbf{question} open is not evidence about either condi... \\
\end{longtable}
\kwichead{rather}{39}
\begin{longtable}{@{}>{\raggedright}p{0.07\textwidth}>{\raggedleft}p{0.41\textwidth}@{\,}>{\raggedright\arraybackslash}p{0.47\textwidth}@{}}
01-c Abs & Assumption 3 is thereby derived  & \textbf{rather} than assumed. \\
01-c 4 & ...at all, a case of the absence of agency  & \textbf{rather} than of its exercise, and falls outside... \\
01-d 4 & ... 2 denies the premise of this challenge  & \textbf{rather} than refuting it (Nagel 1971). \\
01-d 4 & ...gues that the fitting response is irony  & \textbf{rather} than despair; the detached standpoint y... \\
04-a Abs & ... that the appropriate response is irony  & \textbf{rather} than despair. \\
04-b 3 & ...is a paradigm of the absence of meaning  & \textbf{rather} than of its presence in a godless world... \\
06-b 2 & ...ince the world it scorns is indifferent  & \textbf{rather} than opposed (Proposition 3.5 of EXI-04... \\
07-b Abs & ...s on a structural feature of final ends  & \textbf{rather} than on a psychological accident. \\
08-a Abs & The ideal asks that a life be one's own  & \textbf{rather} than borrowed from the anonymous public... \\
08-a Abs & ...enticity is possible, as an orientation  & \textbf{rather} than an achievement. \\
08-a 1 & ... and that one ought to live in that way  & \textbf{rather} than in conformity to what is expected ... \\
08-a 2 & ... by examining the commitments concerned  & \textbf{rather} than by covering the question over; and... \\
08-a 3 & ...1 describes would have to be re-entered  & \textbf{rather} than possessed. \\
08-a 5 & ...enticity is possible, as an orientation  & \textbf{rather} than an achievement. \\
09-a 4 & Beauvoir's analysis supports the thesis  & \textbf{rather} than threatening it. \\
10-a 1 & ...ure in which meaning has to be produced  & \textbf{rather} than inherited. \\
10-a 2 & ...ch is between capacities and activities  & \textbf{rather} than between capacities and task diffic... \\
10-b Abs & ...im that over-challenge produces anxiety  & \textbf{rather} than boredom. \\
10-b 3 & ...demands exceeding skill produce anxiety  & \textbf{rather} than boredom. \\
10-b 4 & ...ny gave themselves mild electric shocks  & \textbf{rather} than sit without stimulation. \\
10-b 5 & ...e correlation shows a cost of the state  & \textbf{rather} than its content. \\
10-b A & ...ports and challenges is recorded as S/C  & \textbf{rather} than netted out, and an entry that does... \\
11-a 2 & ...vine commands confirm moral obligations  & \textbf{rather} than create them. \\
13-a 2 & ...erson because they thwart her interests  & \textbf{rather} than because they cut her life short. \\
13-a 2 & ...n it while she lives, by writing a will  & \textbf{rather} than giving her property away, by entru... \\
14-a 4 & That supports,  & \textbf{rather} than undermines, the first component of... \\
15-b 2 & ...es into expectation: the agent plans on  & \textbf{rather} than for its possibility, and the speci... \\
15-b 3 & ...e would make credence respond to desire  & \textbf{rather} than to evidence, which Principle P1 ex... \\
16-a 1 & ...the beloved's value; it creates reasons  & \textbf{rather} than answering to them. \\
16-a 3 & ...t, which are events in the relationship  & \textbf{rather} than changes in the beloved's qualities... \\
16-b Abs & ...swers to a relational reason of its own  & \textbf{rather} than inheriting the old one. \\
16-b 2 & ... does not explain why loves this person  & \textbf{rather} than any other, since the value it cite... \\
17-a 2 & ...that her activity expresses her concern  & \textbf{rather} than being merely performed. \\
17-a 5 & ...ases; and lives devoted to appreciation  & \textbf{rather} than production are meaningful without ... \\
18-a Abs & ...mism are features of normal functioning  & \textbf{rather} than of disorder. \\
19-a 3 & ...n false: most deaths interrupt projects  & \textbf{rather} than complete them. \\
19-a 4 & ...having a shape; that it is limited here  & \textbf{rather} than later is what the deprivation acco... \\
20-a 1 & We apply the account to an activity  & \textbf{rather} than a life: the production of the pres... \\
20-a 3 & ...epresentation of philosophical activity  & \textbf{rather} than an occasion for it; REL-24 argues ... \\
\end{longtable}
\kwichead{refuting}{2}
\begin{longtable}{@{}>{\raggedright}p{0.07\textwidth}>{\raggedleft}p{0.41\textwidth}@{\,}>{\raggedright\arraybackslash}p{0.47\textwidth}@{}}
01-d 4 & ...e premise of this challenge rather than  & \textbf{refuting} it (Nagel 1971). \\
04-a 4 & ...mise of the sceptical challenge without  & \textbf{refuting} it. \\
\end{longtable}
\kwichead{sceptical}{2}
\begin{longtable}{@{}>{\raggedright}p{0.07\textwidth}>{\raggedleft}p{0.41\textwidth}@{\,}>{\raggedright\arraybackslash}p{0.47\textwidth}@{}}
04-a 4 & ...jective worth denies the premise of the  & \textbf{sceptical} challenge without refuting it. \\
08-a Abs & We show that the  & \textbf{sceptical} arguments refute the achievement readin... \\
\end{longtable}
\kwichead{specie}{0}
\kwicnone
\kwichead{sub}{0}
\kwicnone
\kwichead{value}{93}
\begin{longtable}{@{}>{\raggedright}p{0.07\textwidth}>{\raggedleft}p{0.41\textwidth}@{\,}>{\raggedright\arraybackslash}p{0.47\textwidth}@{}}
01-c 2 & It is worthwhile if it has or serves  & \textbf{value} that does not depend on the agent's att... \\
01-c 2 & ...f understanding and the appreciation of  & \textbf{value} (Metz 2013). \\
01-c 2 & ...gent engages in, actively, and that its  & \textbf{value} not depend on the attitudes of those en... \\
01-c 2 & ... when its object is worthwhile, for the  & \textbf{value} of caring for a person or of coming to ... \\
01-c 2 & ...o understand something derives from the  & \textbf{value} of the person or of the understanding, ... \\
01-c 3 & ...ves among other persons, whose good has  & \textbf{value} independent of her attitudes, and every... \\
01-d 2 & ...ion of meaning; scale bears on how much  & \textbf{value} a project serves, which can affect degr... \\
01-d 4 & Definition 1 requires only that the  & \textbf{value} of a project not depend on the attitude... \\
02-a 2 & ...privation account assigns the event its  & \textbf{value} accordingly. \\
02-a 2 & The  & \textbf{value} for of an event is , where is 's actual... \\
02-a 2 & The event is bad for if this  & \textbf{value} is negative, and good for if it is posi... \\
02-a 2 & ...event by the difference it makes to the  & \textbf{value} of the life as a whole, and the lives c... \\
02-a 2 & And it is whole-life: the  & \textbf{value} it assigns is a fact about 's life take... \\
02-a 3 & So her life has less  & \textbf{value} than the life she would have had withou... \\
02-a 3 & It is a fact about and the  & \textbf{value} of her life as a whole, and it obtains,... \\
02-a 3 & By Principle P1, the  & \textbf{value} of her death for is . \\
02-a 3 & ...a positive balance of goods, so and the  & \textbf{value} of her death for her is negative. \\
02-a 4 & ...y Principle P1, what is compared is the  & \textbf{value} of the whole life, not its length. \\
02-a 5 & ...act of Principle P1 is a fact about the  & \textbf{value} of 's whole life and is not located at ... \\
03-a 1 & ...nd Scheffler (2013) has argued that the  & \textbf{value} of many human activities depends on the... \\
03-a 2 & ... including connections of character and  & \textbf{value} (PID-06). \\
03-b 4 & ... in an endless life, set a limit on the  & \textbf{value} of continuation. \\
04-a 1 & ...w of a capacity we have every reason to  & \textbf{value}. \\
04-a 1 & The other asks whether objective  & \textbf{value}, if there is any, dissolves the absurd:... \\
04-a 3 & ...g: the end serves another end, or has a  & \textbf{value}, or is required. \\
04-b 1 & ...ing about, or at an outcome that has no  & \textbf{value} even if brought about. \\
04-b 4 & ...and many projects; it does not undo the  & \textbf{value} of what was done. \\
05-a 3 & ...d, and the degree of its badness is the  & \textbf{value} of what it takes. \\
05-a 4 & ...ere it is, it is a different life whose  & \textbf{value} is not guaranteed to be greater. \\
06-a 4 & The universe will destroy everything we  & \textbf{value}, through death, catastrophe and eventua... \\
06-a 4 & ... the universe that will destroy what we  & \textbf{value} also produced it: the same laws that br... \\
06-b 2 & ...with projects of objective worth, whose  & \textbf{value} does not depend on the attitudes of tho... \\
06-b 2 & ...ject of engagement is worthwhile if its  & \textbf{value} does not depend on the agent's attitude... \\
07-a 1 & (2011) found that people who  & \textbf{value} happiness highly report less happiness ... \\
07-a 1 & ...toms, and that inducing participants to  & \textbf{value} happiness reduced their enjoyment of a ... \\
07-a 3 & ...who aims directly at happiness does not  & \textbf{value} the objects of her activities for their... \\
07-a 3 & ...as a means to one's happiness is not to  & \textbf{value} it for its own sake. \\
07-a 3 & ...who aims directly at happiness does not  & \textbf{value} the objects of her activities for their... \\
07-a 3 & ...g happy but a correct response to their  & \textbf{value}, and the happiness that follows is the ... \\
07-a 3 & ...lows is the by-product of getting their  & \textbf{value} right. \\
07-a 4 & Railton's agent can stably  & \textbf{value} other things for their own sake only be... \\
09-b 2 & These evaluations apply standards of  & \textbf{value}, not a defining project, and they do no... \\
09-b 3 & Its  & \textbf{value} for Sartre lies in the contrast: a huma... \\
10-a 1 & ...evidence that existence has no positive  & \textbf{value}; for Svendsen (2005) it is a characteri... \\
10-b 3 & ...boredom is related to the low perceived  & \textbf{value} of the activity (Pekrun et al. 2010). \\
10-b 3 & ...bear on Assumption 2.4 of EXI-10-a, the  & \textbf{value} of engagement. \\
10-b 11 & ...e from these findings together with the  & \textbf{value} of engagement, and the studies of costl... \\
11-a 2 & ...he principle does not deny that God may  & \textbf{value} what is valuable, or command what is go... \\
11-a 2 & ...a life gains meaning by connection with  & \textbf{value} beyond itself, and more by connection w... \\
11-a 2 & ...lf, and more by connection with greater  & \textbf{value} (Nozick 1981). \\
12-b 2 & ...ar's basic asymmetry is an asymmetry of  & \textbf{value}. \\
13-a 4 & ... wishes of the dead is justified by its  & \textbf{value} to the living, who want assurance that ... \\
13-a 4 & ...the practice were justified only by its  & \textbf{value} to those still living, a secret breach ... \\
14-a 4 & These are reasons to  & \textbf{value} it, not conditions of its rationality, ... \\
15-a 3 & ...credence alone, but depends also on the  & \textbf{value} of and on the agent's capacity to make ... \\
15-a 3 & ...bility as in a likely one with the same  & \textbf{value}. \\
15-a 3 & ...ed to take the same stance whatever the  & \textbf{value} of and whatever they could do about it. \\
15-a 3 & ...h the evidence, and it also varies with  & \textbf{value} and agency, which are the other things ... \\
15-a 4 & Where the outcome is of great  & \textbf{value} and one's efforts are among its causes,... \\
15-b Abs & ...ings equal, but must depend also on the  & \textbf{value} of the outcome and on the agent's power... \\
15-b 1 & ...5-a): writing for the credence, for the  & \textbf{value} of to the agent and for her capacity to... \\
15-b 2 & ...e minimal stance of wishing, unless the  & \textbf{value} of and her capacity to bring it about a... \\
15-b 4 & ...proportion to its probability times its  & \textbf{value}, adjusted for what one's investment cha... \\
16-a 1 & ...ounded in any judgment of the beloved's  & \textbf{value}; it creates reasons rather than answeri... \\
16-a 1 & ...lds that love responds to the beloved's  & \textbf{value} as a person, the rational nature that a... \\
16-a 1 & ...lf and the lover's capacity to see that  & \textbf{value}. \\
16-a 1 & ...ith her, and love is the disposition to  & \textbf{value} both the relationship and the person as... \\
16-a 4 & ...), the reason for love is the beloved's  & \textbf{value} as a person, which is the same in every... \\
16-a 4 & The beloved's  & \textbf{value} as a person may well be a reason for lo... \\
16-a 4 &  & \textbf{Value} as a person is at most a background con... \\
16-b 2 & ... (1999) holds that love responds to the  & \textbf{value} of the beloved's rational nature, which... \\
16-b 2 & ...ts object is irreplaceable because that  & \textbf{value} is a dignity, which admits no substitut... \\
16-b 2 & ...person rather than any other, since the  & \textbf{value} it cites is shared by all. \\
17-a 1 & ... is natural to suppose that this is the  & \textbf{value} that meaning tracks. \\
17-a 2 & ...uiry) is worthwhile if it has or serves  & \textbf{value} that does not depend on the agent's att... \\
17-a 2 & Our thesis is not that it lacks  & \textbf{value}, but that meaning does not require it. \\
17-a 3 & ...d at their object; the object, music of  & \textbf{value}, is worthwhile; her listening bears on ... \\
17-a 3 & ...onent of the good (ETH-01), and it adds  & \textbf{value} to a life; if success enlarges the wort... \\
17-a 4 & Meaning requires connection with  & \textbf{value} beyond oneself, and the way to connect ... \\
17-a 4 & ...t that meaning requires connection with  & \textbf{value} beyond oneself; Principles P1 and P2 re... \\
18-a 3 & It does not: the ideal is a  & \textbf{value} among others, and Corollary 2.4 counsel... \\
19-a Abs & Following Velleman, we take the  & \textbf{value} of a life to depend in part on its shap... \\
19-a 1 & ... limited, and Scheffler (2013) that the  & \textbf{value} of many of our activities depends on th... \\
19-a 1 & Velleman (1991) showed that the  & \textbf{value} of a life is not the sum of the values ... \\
19-a 2 & ...tion of something, or the occasion of a  & \textbf{value} that it does not owe to that life. \\
19-a 2 & The  & \textbf{value} of a life as a whole depends in part on... \\
19-a 3 & ...filment of a purpose or the bearer of a  & \textbf{value}. \\
19-a 3 & ...nciple P1 this shape contributes to the  & \textbf{value} of the life. \\
19-a 3 & ... the deprivation account of EXI-02, the  & \textbf{value} of a death for the one who dies is the ... \\
19-a 3 & ... who dies is the difference between the  & \textbf{value} of her actual life and the value of the... \\
19-a 3 & ...en the value of her actual life and the  & \textbf{value} of the life she would have had had she ... \\
19-a 5 & ...filment of a purpose or the bearer of a  & \textbf{value} of its own, and the meanings attributed... \\
19-b 2 & ...d Scheffler (2013) that our capacity to  & \textbf{value} depends on the awareness of temporal sc... \\
\end{longtable}
\kwichead{worth}{82}
\begin{longtable}{@{}>{\raggedright}p{0.07\textwidth}>{\raggedleft}p{0.41\textwidth}@{\,}>{\raggedright\arraybackslash}p{0.47\textwidth}@{}}
01-c Abs & ...ge in at least one project of objective  & \textbf{worth} (Assumption 3). \\
01-c 1 & ...e engagement with projects of objective  & \textbf{worth} (Definition 1), assumes that there are ... \\
01-c 1 & ...y poor may have no project of objective  & \textbf{worth} on which they can embark, and Wolf (201... \\
01-c 2 & ...er engagement is a project of objective  & \textbf{worth} in the sense of Definition 1 of EXI-01-... \\
01-c 2 & ...is engagement in a project of objective  & \textbf{worth}. \\
01-c 2 & ... engagement with a project of objective  & \textbf{worth}, and by Definition 1 of EXI-01-a it is ... \\
01-c 5 & ...ent with them is a project of objective  & \textbf{worth}, and a life that contains it has meanin... \\
01-d Abs & ... one engages with projects of objective  & \textbf{worth}, and that meaning in this sense is avai... \\
01-d 2 & ... which is within the agent's power, and  & \textbf{worth}, which belongs to the object, its avail... \\
01-d 2 & ...ning by engagement and by the objective  & \textbf{worth} of what one engages with. \\
01-d 2 & ...to a neighbor are projects of objective  & \textbf{worth}, and engagement with them confers meani... \\
01-d 3 & ...s with the extent of engagement and the  & \textbf{worth} of its object, and lives differ in both... \\
01-d 4 & ...would hold against a view that measured  & \textbf{worth} by prestige or required projects of a p... \\
01-d 5 & ...gagement, which is up to the agent, and  & \textbf{worth}, which belongs to what she engages with... \\
02-a 3 & ... would otherwise have gone on to a life  & \textbf{worth} living, this condition is met. \\
02-a 3 & ...nd if those circumstances made her life  & \textbf{worth} living, the continuation has a positive... \\
03-a 2 & ...ject engages with projects of objective  & \textbf{worth}, and EXI-17 that such engagement, not a... \\
03-a 4 & ...at the world does not run out of things  & \textbf{worth} doing. \\
03-b 2 & ... the condition of all projects, and its  & \textbf{worth} for the person is derivative from the w... \\
03-b 2 & ...h for the person is derivative from the  & \textbf{worth} of what would fill it. \\
04-a 3 & ...rs that the consideration has objective  & \textbf{worth}, the detached standpoint can register t... \\
04-a 3 & ...s not assume that there is no objective  & \textbf{worth}. \\
04-a 3 & ... one engages with projects of objective  & \textbf{worth}, and we accept that conclusion. \\
04-a 3 & ...oint cannot be forced to recognize that  & \textbf{worth}, so that the regress of justification h... \\
04-a 3 & ...worthwhile is one question; whether its  & \textbf{worth} can be certified to a standpoint that s... \\
04-a 3 & ...osition to certify that one's ends lack  & \textbf{worth} than to certify that they have it. \\
04-a 4 & ...-a puts it, the assumption of objective  & \textbf{worth} denies the premise of the sceptical cha... \\
04-b 1 & ...n that his task achieves nothing and is  & \textbf{worth} nothing. \\
04-b 3 & ... vindication, not that its objects lack  & \textbf{worth} (Wolf 2010). \\
04-b 5 & ...aces no adversary, and contains objects  & \textbf{worth} the effort we give them. \\
06-b Abs & ...1, consists in engagement with what has  & \textbf{worth} independently of anyone's attitudes, an... \\
06-b Abs & ...tudes toward us takes nothing from that  & \textbf{worth}. \\
06-b 2 & ...ject engages with projects of objective  & \textbf{worth}, whose value does not depend on the att... \\
06-b 2 & If the  & \textbf{worth} of a project does not depend on the att... \\
06-b 2 & ...and the flourishing of persons have the  & \textbf{worth} they have whether or not anything in th... \\
06-b 3 & ...rence to any particular agent, then the  & \textbf{worth} of objects of engagement is precisely w... \\
06-b 3 & ...h of objects of engagement is precisely  & \textbf{worth} from that standpoint, and the demand is... \\
06-b 3 & ...s in engagement with what has objective  & \textbf{worth}, as EXI-01 holds, then it does not requ... \\
06-b 4 & The  & \textbf{worth} of what we engage with does not wait on... \\
09-a 2 & ...one values oneself and finds one's life  & \textbf{worth} living (Korsgaard 1996). \\
09-a 5 & ...tent that its projects are of objective  & \textbf{worth}. \\
10-a 2 & ...ention and abilities in a way she finds  & \textbf{worth} sustaining: because it demands too litt... \\
10-a 3 & ...ention and abilities in a way she finds  & \textbf{worth} sustaining, which is a mismatch (Defini... \\
11-a Abs & ...ngages with what is worthwhile, and the  & \textbf{worth} of knowledge, of persons and of beauty ... \\
11-a Abs & ...ge, of persons and of beauty is not the  & \textbf{worth} of a role assigned to us by a designer. \\
11-a Abs & ...ing comes in degrees and grows with the  & \textbf{worth} of what one engages: a framework contai... \\
11-a Abs & ...framework containing a being of maximal  & \textbf{worth}, a purpose that unifies one's projects,... \\
11-a 1 & ... a connection to something of unlimited  & \textbf{worth}, and only God is such a thing (Nozick 1... \\
11-a 1 & ...s active engagement meets the objective  & \textbf{worth} of its object. \\
11-a 2 & ...ual, engagement with what is of greater  & \textbf{worth}, or engagement that relates one to more... \\
11-a 3 & By Principle P1, the  & \textbf{worth} of these activities is not conferred by... \\
11-a 4 & ...-01 defends contains a being of maximal  & \textbf{worth}. \\
11-a 4 & ... engagement with the greatest object of  & \textbf{worth} there is; by Principle P2, it would mak... \\
11-a 5 &  & \textbf{Worth} realized at a time is not cancelled by ... \\
11-a 5 & ...and the claim that it is would make the  & \textbf{worth} of a good depend on what happens after ... \\
11-a 5 & The regress arises only if  & \textbf{worth} is always borrowed from something else. \\
11-a 6 & ...inally, it relies on the objectivity of  & \textbf{worth} established elsewhere (ETH-01, MTE-20);... \\
11-a 7 & ...gement with what is worthwhile, and the  & \textbf{worth} of the objects of engagement does not d... \\
11-a 7 & ...aning, by offering an object of maximal  & \textbf{worth}, a purpose that unifies one's engagemen... \\
11-b Abs & ...rticular life-plans, not the objects of  & \textbf{worth} available to engagement, so that Propos... \\
11-b 1 & ...) and holds that meaning grows with the  & \textbf{worth} of what one engages and with the extent... \\
11-b 1 & ...ramework would make: a being of maximal  & \textbf{worth} with whom one could stand in relation, ... \\
11-b 1 & ...ts and relations of engagement of great  & \textbf{worth} that a naturalistic world lacks, Propos... \\
11-b 2 & The  & \textbf{worth} of knowledge, of persons and of beauty ... \\
11-b 2 & The  & \textbf{worth} of knowledge, persons and beauty is not... \\
11-b 2 & ...d the framework adds objects of greater  & \textbf{worth} (Principle P2 of EXI-11-a) while leavin... \\
11-b 3 & ...ts and relations of engagement of great  & \textbf{worth}. \\
15-a 3 & ...a horse race lacks, because recovery is  & \textbf{worth} more and because her efforts can affect... \\
17-a Abs & ...that are alike in engagement and in the  & \textbf{worth} of their objects and differ only in suc... \\
17-a 3 & ...alue to a life; if success enlarges the  & \textbf{worth} realized through one's engagement, it m... \\
17-a 4 & Taylor's case fails the condition of  & \textbf{worth}, and the conclusion it is used to suppo... \\
19-b Abs & ...at an immortal being could find nothing  & \textbf{worth} doing. \\
19-b Abs & ... of meaning, since it adds no object of  & \textbf{worth} to engage with. \\
19-b 1 & The  & \textbf{worth} of the objects of engagement, knowledge... \\
19-b 1 & So the limit adds neither an object of  & \textbf{worth} nor worth to an object, and by the acco... \\
19-b 1 & ...mit adds neither an object of worth nor  & \textbf{worth} to an object, and by the account of EXI... \\
19-b 2 & But Proposition 1.1 concerns the  & \textbf{worth} of what we engage with, and urgency add... \\
19-b 2 & ...hat we engage with, and urgency adds no  & \textbf{worth}: the friend's good matters no less to s... \\
20-a 3 & ...arguments are correct, the collection's  & \textbf{worth} depends on the minority of readers who ... \\
20-a 3 & ... who engage with its arguments, and its  & \textbf{worth} for them depends on whether the argumen... \\
20-b Abs & Evidence bearing on  & \textbf{worth}.. \\
20-b Abs & EXI-20-a makes the  & \textbf{worth} of the collection depend on whether it ... \\
\end{longtable}
}

\section{Cross-references to the numbered statements of the discipline}\label{app:xref}

\subsection*{Preface}

\paragraph{What is enumerated.} This appendix lists every numbered statement of every manuscript of the discipline other than EXI-01-a and the present work, \cnStatements{} statements in all, and records for each the lemmata of EXI-01-a on which it bears by the rule stated below. A numbered statement is a definition, principle, assumption, lemma, theorem, corollary, proposition, remark or objection that carries a label in its manuscript; its number, page and title are read from the manuscript's compiled \texttt{.aux} file, so that the numbers are those under which the statement is cited, and its kind and text are read from the environment that carries the label in \texttt{main.tex}. Manuscripts whose statements are unnumbered (EXI-20-a and EXI-20-b) contribute no rows, and section labels are not statements.

\paragraph{The rule.} The text of each statement, together with its title, is assigned the concept classes listed in Table~B.1, each defined by the expressions that signal it, matched case-insensitively at the beginning of words. A statement bears on a lemma group of EXI-01-a if its classes satisfy the condition for the group: on the title \Lm{1}, if it has \texttt{M} and \texttt{L}; on Definition~1 \Lm{2}, if it has \texttt{M} and at least two of \texttt{E}, \texttt{W}, \texttt{P}, \texttt{D}, at least one of them other than \texttt{D}; on the degree clause \Lm{2.3}, if it has \texttt{D} and \texttt{M} or \texttt{W}; on Assumption~2 \Lm{3}, if it has \texttt{O} and \texttt{W} or \texttt{P}; on Assumption~3 \Lm{4}, if it has \texttt{I} or \texttt{A}, and \texttt{E} or \texttt{P}, and \texttt{W} or \texttt{M}; on Theorem~4 \Lm{5}, if it has \texttt{A} and \texttt{M} or \texttt{W}; on the meaning of particular lives \Lm{10}, if it has \texttt{M}, \texttt{L} and \texttt{E} and neither \texttt{A} nor \texttt{S}; on the remark on the absurd \Lm{11}, if it has \texttt{S}.

The rule is lexical, and it is meant as a finding aid, not as a judgment. It records where the vocabulary of a statement overlaps with the vocabulary in which a lemma is stated, which is where a reader of the lemma should look; it cannot record a relation expressed in other words, and it records some overlaps that are not relations. Both kinds of error are illustrated below. Its advantage is that it is exhaustive and reproducible: every statement of the discipline is examined, and the same rule applied to the same files gives the same table.

\paragraph{Summary.} Of the \cnStatements{} statements, \cnBearing{} bear on at least one lemma group and \cnNotBearing{} on none. By lemma group (Table~B.2), \cnDefOne{} bear on Definition~1, \cnAssTwo{} on Assumption~2, \cnAssThree{} on Assumption~3, \cnThmFour{} on Theorem~4 and \cnAbsurd{} on the remark on the absurd. Table~B.3 gives the counts by manuscript. The statements bearing on the lemmata are concentrated in the family itself (EXI-01-c, EXI-01-d) and in EXI-17-a, EXI-11-a, EXI-04-a and EXI-06-a, which is the distribution the comments on the lemmata found by reading; the families on death, love, hope and the dead (EXI-02, EXI-05, EXI-13, EXI-15, EXI-16) have almost no statements that bear on EXI-01-a by the rule, which agrees with their subject matter.

\paragraph{Entries discussed.} (i) EXI-17-a, Definition~2.2 (\enquote{Worthwhile}), bears by the rule on \Lm{3} and \Lm{4}. Both bearings are real: the definition restates the restricted notion of objective worth (\Lm{3.4}) and lists the objects, persons, practices and inquiries, whose availability supports Assumption~3. (ii) EXI-04-a, Objection~2 (\enquote{Objective worth dissolves the absurd}), bears on \Lm{3} and \Lm{11}. It is the statement of the discipline that most directly concerns the relation between Assumption~2 and the sceptical challenge, and the comment on \Lm{11.4} relies on its reply. (iii) EXI-04-a, Lemma~3.1 (\enquote{Ineliminable seriousness}), bears on nothing by the rule, though the comment on \Lm{4.2} and \Lm{4} relies on it: it shares Assumption~3's domain, agents capable of intentional action, and supplies the subjective capacity on which engagement depends, but it speaks of ends taken seriously, not of worth or meaning, and so lacks the classes the rule requires. It is a false negative of the kind the rule cannot avoid. (iv) EXI-04-a, Remark~3.7, on Camus's Sisyphus, is also assigned to no lemma, because it speaks of Sisyphus, scorn and lucidity and never uses the word \emph{absurd}; Section~\ref{s:absurd} relies on it. (v) EXI-14-a, Definitions~2.1 and~2.2, Lemma~3.1 and Theorem~3.2, on nostalgia, bear by the rule on Theorem~4, because they speak of goods that are no longer available to the subject. These are false positives: the availability of past goods to memory has nothing to do with the availability of meaning. (vi) EXI-10-a, Definition~2.2 (\enquote{Mismatch}), bears on \Lm{4} and \Lm{5}. The bearing is real, though indirect. The definition makes meaninglessness, an activity's serving nothing the agent cares about, one of the three forms of the mismatch that boredom signals; if EXI-10 is right that boredom has the function of moving the agent toward matching activities, the agent's affective life contains a mechanism that tends toward engagement, which is the subjective condition that Assumption~3's ability to engage presupposes.

\subsection*{B.1 Concept classes}

{\footnotesize% Generated by gen/generate.py; do not edit.
\begin{tabular}{@{}l l >{\raggedright\arraybackslash}p{0.68\textwidth}@{}}
\toprule Code & Class & Expressions (case-insensitive; matched from the start of a word; \textsuperscript{w}: whole word only) \\ \midrule
\texttt{M} & meaning & \texttt{meaning(ful|less)?}\textsuperscript{w}, \texttt{meanings}\textsuperscript{w}, \texttt{point}\textsuperscript{w}, \texttt{significan} \\
\texttt{E} & engagement & \texttt{engag}, \texttt{absor}, \texttt{care(s|d)?}\textsuperscript{w}, \texttt{caring}\textsuperscript{w}, \texttt{gripped} \\
\texttt{W} & worth & \texttt{worth}, \texttt{valu}, \texttt{good(s)?}\textsuperscript{w}, \texttt{better}\textsuperscript{w}, \texttt{bad}\textsuperscript{w} \\
\texttt{O} & objectivity & \texttt{objectiv}, \texttt{independen}, \texttt{not depend}, \texttt{attitude}, \texttt{subjectiv}, \texttt{whether or not} \\
\texttt{P} & projects & \texttt{projects?}\textsuperscript{w}, \texttt{pursuits?}\textsuperscript{w}, \texttt{activit}, \texttt{ends}\textsuperscript{w}, \texttt{final end}, \texttt{undertak} \\
\texttt{A} & availability & \texttt{availab}, \texttt{everyone}\textsuperscript{w}, \texttt{every agent}, \texttt{any agent}, \texttt{universal}, \texttt{luck}\textsuperscript{w}, \texttt{circumstanc}, \texttt{any reflective agent} \\
\texttt{I} & agency & \texttt{agents?}\textsuperscript{w}, \texttt{intentional action} \\
\texttt{D} & degree & \texttt{degree}, \texttt{extent}\textsuperscript{w}, \texttt{more meaningful}, \texttt{less meaningful}, \texttt{add(s)? to}\textsuperscript{w} \\
\texttt{S} & absurd & \texttt{absurd}, \texttt{detached}, \texttt{sub specie}, \texttt{nothing matters}, \texttt{cosmic}, \texttt{universe}\textsuperscript{w}, \texttt{indifferen}, \texttt{despair} \\
\texttt{L} & life & \texttt{life}\textsuperscript{w}, \texttt{lives}\textsuperscript{w} \\
\bottomrule \end{tabular}
\par}

\subsection*{B.2 Statements by lemma group}

{\footnotesize% Generated by gen/generate.py; do not edit.
\begin{tabular}{@{}l l r@{}}
\toprule Lemma group & Content & Statements \\ \midrule
{[1]} & title & 36 \\
{[2]} & Def. 1 & 14 \\
{[2.3]} & degree & 9 \\
{[3]} & Ass. 2 & 11 \\
{[4]} & Ass. 3 & 11 \\
{[5]} & Thm. 4 & 16 \\
{[10]} & particular lives & 10 \\
{[11]} & absurd & 29 \\
\bottomrule \end{tabular}
\par}

\subsection*{B.3 Statements by manuscript and lemma group}

{\footnotesize% Generated by gen/generate.py; do not edit.
\begin{longtable}{@{}l r r r r r r r r r@{}}
\toprule Ms & Stmts & [1] & [2] & [2.3] & [3] & [4] & [5] & [10] & [11] \\ \midrule \endhead
\bottomrule \endfoot
EXI-01-c & 4 & 1 & . & . & 2 & 3 & 2 & . & . \\
EXI-01-d & 4 & 4 & . & 1 & 2 & . & 2 & . & . \\
EXI-02-a & 13 & . & . & . & . & . & 1 & . & . \\
EXI-02-b & 1 & . & . & . & . & . & . & . & . \\
EXI-03-a & 12 & 4 & 2 & 1 & 1 & . & . & 1 & . \\
EXI-03-b & 3 & 1 & . & . & . & . & . & . & . \\
EXI-04-a & 16 & . & 2 & . & 1 & 1 & 1 & . & 11 \\
EXI-04-b & 3 & . & . & . & . & . & . & . & 1 \\
EXI-05-a & 9 & . & . & . & . & . & . & . & . \\
EXI-06-a & 8 & . & . & . & . & . & . & . & 7 \\
EXI-06-b & 2 & 1 & . & . & . & . & . & . & 2 \\
EXI-07-a & 13 & . & . & 1 & 1 & 3 & . & . & . \\
EXI-07-b & 2 & . & . & . & . & . & . & . & . \\
EXI-08-a & 10 & . & . & . & . & . & . & . & . \\
EXI-09-a & 12 & . & . & . & . & . & . & . & . \\
EXI-09-b & 3 & . & . & . & . & . & . & . & . \\
EXI-10-a & 12 & . & . & . & . & 1 & 1 & . & 1 \\
EXI-11-a & 16 & 11 & 4 & 2 & 1 & 2 & 2 & 2 & 6 \\
EXI-11-b & 1 & . & . & 1 & . & . & 1 & . & . \\
EXI-12-a & 8 & . & . & . & . & . & . & . & . \\
EXI-12-b & 2 & . & . & . & . & . & . & . & . \\
EXI-13-a & 13 & . & . & . & . & . & . & . & . \\
EXI-13-b & 3 & . & . & . & . & . & . & . & . \\
EXI-14-a & 8 & . & . & . & . & . & 4 & . & . \\
EXI-15-a & 11 & . & . & . & . & . & . & . & . \\
EXI-15-b & 3 & . & . & . & . & . & . & . & . \\
EXI-16-a & 16 & . & . & . & . & . & 1 & . & . \\
EXI-16-b & 3 & . & . & . & . & . & . & . & . \\
EXI-17-a & 16 & 10 & 5 & 3 & 1 & 1 & 1 & 5 & . \\
EXI-18-a & 8 & . & . & . & 1 & . & . & . & 1 \\
EXI-18-b & 1 & . & . & . & . & . & . & . & . \\
EXI-19-a & 10 & 1 & . & . & 1 & . & . & . & . \\
EXI-19-b & 3 & 3 & 1 & . & . & . & . & 2 & . \\
\end{longtable}
}

\subsection*{B.4 The statements}

{\footnotesize Columns: manuscript; number of the statement; kind; title (or, if untitled, its first words); page of the manuscript on which it stands; concept classes; lemma groups of EXI-01-a on which it bears by the rule.\par}

{\scriptsize% Generated by gen/generate.py; do not edit.
\begin{longtable}{@{}l l l >{\raggedright}p{0.25\textwidth} r l >{\raggedright\arraybackslash}p{0.2\textwidth}@{}}
\toprule Ms & No. & Kind & Title & p. & Classes & Bears on \\ \midrule \endhead
\bottomrule \endfoot
EXI-01-c & 2.1 & Def. & \emph{An object of engagement is anything at\ldots} & 1 & \texttt{EWOPI} & [3], [4] \\
 & 2.2 & Ass. & Reach & 1 & \texttt{EWPAI} & [4], [5] \\
 & 2.3 & Lem. & Projects from engagement & 2 & \texttt{EWOPI} & [3], [4] \\
 & 2.4 & Thm. & \emph{For every agent capable of intentional action,\ldots} & 2 & \texttt{MAIL} & [1], [5] \\
EXI-01-d & 2.1 & Prop. & Independence from outcome & 1 & \texttt{MOPL} & [1], [3] \\
 & 2.2 & Prop. & Independence from talent and scale & 1 & \texttt{MOPL} & [1], [3] \\
 & 2.3 & Prop. & No condemnation by circumstance & 2 & \texttt{MAIL} & [1], [5] \\
 & 2.4 & Prop. & Degree, not possibility & 2 & \texttt{MADL} & [1], [2.3], [5] \\
EXI-02-a & 2.1 & Def. & Death & 2 & \texttt{L} & --- \\
 & 2.2 & Def. & Life value & 2 & \texttt{WL} & --- \\
 & P1 & Pr. & Deprivation & 2 & \texttt{WL} & --- \\
 & 2.3 & Ass. & Tenseless existence & 2 & \texttt{} & --- \\
 & 3.1 & Lem. & Badness without experience & 3 & \texttt{W} & --- \\
 & 3.2 & Lem. & A subject without a time of loss & 3 & \texttt{W} & --- \\
 & 3.3 & Thm. & Death is bad for the one who dies & 3 & \texttt{WL} & --- \\
 & 3.4 & Cor. & Where Epicurus goes wrong & 3 & \texttt{W} & --- \\
 & 4.1 & Prop. & Asymmetry of the alternatives & 4 & \texttt{WL} & --- \\
 & 1 & Obj. & Timing & 4 & \texttt{W} & --- \\
 & 2 & Obj. & Rosenbaum & 5 & \texttt{W} & --- \\
 & 3 & Obj. & Too much & 5 & \texttt{WAL} & [5] \\
 & 4 & Obj. & Who is deprived & 5 & \texttt{} & --- \\
EXI-02-b & 2.1 & Prop. & \emph{The failure of (E1) leaves (E2) to\ldots} & 1 & \texttt{O} & --- \\
EXI-03-a & 2.1 & Def. & Categorical desire & 2 & \texttt{MWOPL} & [1], [2], [3] \\
 & 2.2 & Def. & Endless life & 2 & \texttt{L} & --- \\
 & 2.3 & Ass. & Bounded repertoire & 2 & \texttt{EWP} & --- \\
 & 2.4 & Ass. & Exhaustion & 2 & \texttt{MEPL} & [1], [2], [10] \\
 & 2.5 & Ass. & Connectedness & 2 & \texttt{WDL} & [2.3] \\
 & 3.1 & Lem. & First horn & 3 & \texttt{ML} & [1] \\
 & 3.2 & Lem. & Second horn & 3 & \texttt{L} & --- \\
 & 3.3 & Thm. & Immortality is not desirable & 3 & \texttt{ML} & [1] \\
 & 1 & Obj. & Repeatable pleasures & 3 & \texttt{WL} & --- \\
 & 2 & Obj. & Open-ended pursuits & 4 & \texttt{EP} & --- \\
 & 3 & Obj. & Boredom is contingent & 4 & \texttt{L} & --- \\
 & 4 & Obj. & Gradual change & 4 & \texttt{L} & --- \\
EXI-03-b & 2.1 & Prop. & \emph{The desire to live forever is not\ldots} & 1 & \texttt{ML} & [1] \\
 & 2.2 & Prop. & \emph{That an endless life is not desirable\ldots} & 2 & \texttt{WL} & --- \\
 & 2.3 & Prop. & \emph{The result establishes that some finite length\ldots} & 2 & \texttt{L} & --- \\
EXI-04-a & 2.1 & Def. & Seriousness & 2 & \texttt{I} & --- \\
 & 2.2 & Def. & Detached standpoint & 2 & \texttt{PIS} & [11] \\
 & 2.3 & Def. & Absurd & 2 & \texttt{PASL} & [11] \\
 & P1 & Pr. & Fit & 2 & \texttt{OD} & --- \\
 & 2.4 & Def. & Three responses & 2 & \texttt{MWPAIS} & [2], [4], [5], [11] \\
 & 3.1 & Lem. & Ineliminable seriousness & 2 & \texttt{PAI} & --- \\
 & 3.2 & Lem. & Ineliminable detachment & 3 & \texttt{AIS} & [11] \\
 & 3.3 & Thm. & \emph{The life of any reflective agent capable\ldots} & 3 & \texttt{AISL} & [11] \\
 & 3.4 & Prop. & Despair does not fit & 3 & \texttt{S} & [11] \\
 & 3.5 & Prop. & Revolt does not fit & 3 & \texttt{S} & [11] \\
 & 3.6 & Thm. & Irony fits & 4 & \texttt{SL} & [11] \\
 & 3.7 & Rem. & \emph{Camus's Sisyphus is happy because he is\ldots} & 4 & \texttt{} & --- \\
 & 1 & Obj. & Evasion & 4 & \texttt{S} & [11] \\
 & 2 & Obj. & Objective worth dissolves the absurd & 4 & \texttt{EWOPSL} & [3], [11] \\
 & 3 & Obj. & Irony corrodes engagement & 4 & \texttt{MEW} & [2] \\
 & 4 & Obj. & Hume's remedy & 5 & \texttt{S} & [11] \\
EXI-04-b & 1.1 & Def. & \emph{An activity is futile if it is\ldots} & 1 & \texttt{WP} & --- \\
 & 1.2 & Prop. & \emph{Absurdity and futility are independent. A life\ldots} & 1 & \texttt{OSL} & [11] \\
 & 2.1 & Prop. & \emph{The defensible content of Camus's conclusion is\ldots} & 2 & \texttt{EP} & --- \\
EXI-05-a & 2.1 & Def. & Fittingness & 1 & \texttt{} & --- \\
 & P1 & Pr. & Fear & 1 & \texttt{W} & --- \\
 & P2 & Pr. & Regret & 1 & \texttt{W} & --- \\
 & 3.1 & Lem. & \emph{One's own death as such is certain.\ldots} & 2 & \texttt{} & --- \\
 & 3.2 & Thm. & \emph{Fear of death as such is not\ldots} & 2 & \texttt{} & --- \\
 & 3.3 & Thm. & \emph{Regret at the prospect of one's death\ldots} & 2 & \texttt{W} & --- \\
 & 1 & Obj. & Premature death & 2 & \texttt{} & --- \\
 & 2 & Obj. & Dread & 2 & \texttt{} & --- \\
 & 3 & Obj. & Asymmetry & 3 & \texttt{} & --- \\
EXI-06-a & 2.1 & Def. & Responsiveness & 1 & \texttt{} & --- \\
 & 2.2 & Def. & Care, hostility, indifference & 1 & \texttt{ES} & [11] \\
 & 2.3 & Ass. & Physical closure & 2 & \texttt{S} & [11] \\
 & 3.1 & Thm. & \emph{The universe is indifferent to us, and\ldots} & 2 & \texttt{ES} & [11] \\
 & 3.2 & Cor. & \emph{The description of the universe as hostile,\ldots} & 2 & \texttt{ES} & [11] \\
 & 1 & Obj. & Harm outweighs benefit & 2 & \texttt{EWSL} & [11] \\
 & 2 & Obj. & Anthropomorphism cuts both ways & 3 & \texttt{OIS} & [11] \\
 & 3 & Obj. & Theism & 3 & \texttt{ES} & [11] \\
EXI-06-b & 2.1 & Prop. & \emph{That the universe is indifferent to us\ldots} & 1 & \texttt{MSL} & [1], [11] \\
 & 2.2 & Prop. & \emph{Cosmic indifference gives no ground for despair\ldots} & 2 & \texttt{S} & [11] \\
EXI-07-a & 2.1 & Def. & Final end & 2 & \texttt{PI} & --- \\
 & 2.2 & Def. & Direct aim & 2 & \texttt{PI} & --- \\
 & 2.3 & Def. & By-product & 2 & \texttt{WPID} & [2.3], [4] \\
 & P1 & Pr. & Engagement & 2 & \texttt{EWP} & --- \\
 & P2 & Pr. & Object-directedness & 2 & \texttt{EWOI} & [3], [4] \\
 & 3.1 & Lem. & \emph{Happiness is a by-product of engagement in\ldots} & 2 & \texttt{EWP} & --- \\
 & 3.2 & Lem. & \emph{An agent who aims directly at happiness\ldots} & 3 & \texttt{WPI} & [4] \\
 & 3.3 & Thm. & \emph{Aiming directly at happiness is self-defeating: an\ldots} & 3 & \texttt{EPI} & --- \\
 & 3.4 & Cor. & \emph{Happiness is not the goal of life,\ldots} & 3 & \texttt{PL} & --- \\
 & 1 & Obj. & The sophisticated hedonist & 3 & \texttt{P} & --- \\
 & 2 & Obj. & Aristotle & 4 & \texttt{P} & --- \\
 & 3 & Obj. & Some happiness is aimed at & 4 & \texttt{P} & --- \\
 & 4 & Obj. & Engagement is only causally required & 4 & \texttt{E} & --- \\
EXI-07-b & 1.1 & Prop. & \emph{The policy of aiming directly at happiness\ldots} & 1 & \texttt{P} & --- \\
 & 2.1 & Prop. & \emph{Either the sophisticated hedonist values the objects\ldots} & 2 & \texttt{EWP} & --- \\
EXI-08-a & 2.1 & Def. & Achievement & 2 & \texttt{I} & --- \\
 & 2.2 & Def. & Orientation & 2 & \texttt{I} & --- \\
 & P1 & Pr. & Public sources & 2 & \texttt{MI} & --- \\
 & P2 & Pr. & Non-coincidence & 2 & \texttt{I} & --- \\
 & 3.1 & Lem. & \emph{Authenticity in the achievement sense is not\ldots} & 2 & \texttt{} & --- \\
 & 3.2 & Lem. & \emph{Authenticity in the orientation sense is consistent\ldots} & 2 & \texttt{} & --- \\
 & 3.3 & Thm. & \emph{Authenticity is possible as an orientation, though\ldots} & 3 & \texttt{} & --- \\
 & 1 & Obj. & Deflation & 3 & \texttt{} & --- \\
 & 2 & Obj. & Social origin & 3 & \texttt{} & --- \\
 & 3 & Obj. & Self-deception & 3 & \texttt{I} & --- \\
EXI-09-a & 2.1 & Def. & Defining project & 1 & \texttt{P} & --- \\
 & 2.2 & Def. & Precedence & 2 & \texttt{P} & --- \\
 & P1 & Pr. & Artefact kinds & 2 & \texttt{} & --- \\
 & P2 & Pr. & Projects of agents & 2 & \texttt{PIL} & --- \\
 & 3.1 & Lem. & \emph{For every artefact, essence precedes existence.\ldots} & 2 & \texttt{} & --- \\
 & 3.2 & Lem. & \emph{For every human being, existence precedes essence.\ldots} & 2 & \texttt{} & --- \\
 & 3.3 & Thm. & \emph{Sartre's distinction holds: for human beings existence\ldots} & 3 & \texttt{} & --- \\
 & 3.4 & Cor. & \emph{Theorem is compatible with the thesis that\ldots} & 3 & \texttt{} & --- \\
 & 1 & Obj. & Theism & 3 & \texttt{} & --- \\
 & 2 & Obj. & Human nature & 3 & \texttt{WP} & --- \\
 & 3 & Obj. & Social formation & 3 & \texttt{} & --- \\
 & 4 & Obj. & Repurposing & 4 & \texttt{} & --- \\
EXI-09-b & 2.1 & Prop. & Two notions of essence & 1 & \texttt{P} & --- \\
 & 2.2 & Prop. & The accidental knife & 2 & \texttt{} & --- \\
 & 2.3 & Prop. & Defect & 2 & \texttt{L} & --- \\
EXI-10-a & 2.1 & Def. & Boredom & 2 & \texttt{EP} & --- \\
 & 2.2 & Def. & Mismatch & 2 & \texttt{EWPAI} & [4], [5] \\
 & 2.3 & Def. & Function & 2 & \texttt{} & --- \\
 & P1 & Pr. & Emotional content & 2 & \texttt{EA} & --- \\
 & 2.4 & Ass. & Value of engagement & 2 & \texttt{EWP} & --- \\
 & 3.1 & Lem. & \emph{Boredom represents a mismatch, and is correct\ldots} & 2 & \texttt{} & --- \\
 & 3.2 & Lem. & \emph{Boredom motivates the search for an alternative\ldots} & 2 & \texttt{P} & --- \\
 & 3.3 & Thm. & \emph{Boredom signals a mismatch between capacities and\ldots} & 3 & \texttt{PI} & --- \\
 & 3.4 & Cor. & \emph{Chronic boredom is either a malfunction, when\ldots} & 3 & \texttt{PI} & --- \\
 & 1 & Obj. & Profound boredom & 3 & \texttt{PS} & [11] \\
 & 2 & Obj. & Boredom harms & 3 & \texttt{} & --- \\
 & 3 & Obj. & Function is not value & 4 & \texttt{W} & --- \\
EXI-11-a & 2.1 & Def. & Meaning & 2 & \texttt{MEWDL} & [1], [2], [2.3], [10] \\
 & 2.2 & Def. & Cosmic purpose & 2 & \texttt{SL} & [11] \\
 & 2.3 & Def. & Purpose theory & 2 & \texttt{MSL} & [1], [11] \\
 & P1 & Pr. & Worth is not conferred by assignment & 2 & \texttt{WPI} & [4] \\
 & P2 & Pr. & Degree & 2 & \texttt{MEWDL} & [1], [2], [2.3], [10] \\
 & 3.1 & Lem. & Purpose is not sufficient & 2 & \texttt{MSL} & [1], [11] \\
 & 3.2 & Lem. & Worth without purpose & 3 & \texttt{WOPS} & [3], [11] \\
 & 3.3 & Thm. & \emph{Meaning does not require a cosmic purpose.\ldots} & 3 & \texttt{MSL} & [1], [11] \\
 & 3.4 & Cor. & The naturalist's life & 3 & \texttt{ML} & [1] \\
 & 4.1 & Prop. & Theism adds & 3 & \texttt{MAL} & [1], [5] \\
 & 4.2 & Rem. & \emph{Proposition is conditional on the framework. If\ldots} & 4 & \texttt{MEWAL} & [1], [2], [4], [5] \\
 & 1 & Obj. & Everything comes to nothing & 4 & \texttt{PS} & [11] \\
 & 2 & Obj. & The regress of linkage & 4 & \texttt{ML} & [1] \\
 & 3 & Obj. & The framework of hope & 4 & \texttt{MEWP} & [2] \\
 & 4 & Obj. & Purpose is what meaning means & 5 & \texttt{ML} & [1] \\
 & 5 & Obj. & Instability & 5 & \texttt{ML} & [1] \\
EXI-11-b & 1.1 & Prop. & \emph{The meaning that a theistic framework would\ldots} & 1 & \texttt{MAD} & [2.3], [5] \\
EXI-12-a & 2.1 & Def. & Balance & 1 & \texttt{WL} & --- \\
 & P1 & Pr. & Impersonal comparison & 2 & \texttt{WL} & --- \\
 & P2 & Pr. & Non-comparative benefit & 2 & \texttt{WL} & --- \\
 & 3.1 & Thm. & \emph{A life containing more good than bad\ldots} & 2 & \texttt{WL} & --- \\
 & 3.2 & Prop. & \emph{Benatar's quality-of-life argument does not reinstate the\ldots} & 2 & \texttt{WL} & --- \\
 & 1 & Obj. & Bad for no one & 2 & \texttt{W} & --- \\
 & 2 & Obj. & Testimony of the suffering & 3 & \texttt{W} & --- \\
 & 3 & Obj. & Risk & 3 & \texttt{WL} & --- \\
EXI-12-b & P1 & Pr. & Complaints & 1 & \texttt{W} & --- \\
 & 2.1 & Prop. & \emph{Principle , together with Corollary 3.4 of\ldots} & 1 & \texttt{} & --- \\
EXI-13-a & 2.1 & Def. & Reasonable wish & 2 & \texttt{} & --- \\
 & P1 & Pr. & Surviving interests & 2 & \texttt{WL} & --- \\
 & P2 & Pr. & Reliance & 2 & \texttt{} & --- \\
 & 2.2 & Ass. & Practice & 2 & \texttt{L} & --- \\
 & 2.3 & Ass. & Weight & 2 & \texttt{} & --- \\
 & 3.1 & Lem. & Subject & 2 & \texttt{} & --- \\
 & 3.2 & Lem. & Content & 3 & \texttt{} & --- \\
 & 3.3 & Lem. & Weakening & 3 & \texttt{} & --- \\
 & 3.4 & Thm. & Duties to the dead & 3 & \texttt{} & --- \\
 & 1 & Obj. & No subject & 4 & \texttt{} & --- \\
 & 2 & Obj. & Persistence & 4 & \texttt{} & --- \\
 & 3 & Obj. & Only the living & 4 & \texttt{W} & --- \\
 & 4 & Obj. & Deathbed promises & 4 & \texttt{} & --- \\
EXI-13-b & 1.1 & Def. & \emph{A posthumous event harms a person if\ldots} & 1 & \texttt{} & --- \\
 & 1.2 & Prop. & \emph{Posthumous harm and posthumous wrong come apart\ldots} & 1 & \texttt{} & --- \\
 & 2.1 & Prop. & \emph{The dead, considered as what remains after\ldots} & 2 & \texttt{} & --- \\
EXI-14-a & P1 & Pr. & Emotions as evaluations & 1 & \texttt{} & --- \\
 & 2.1 & Def. & Nostalgia & 2 & \texttt{WA} & [5] \\
 & 2.2 & Def. & Rational nostalgia & 2 & \texttt{WA} & [5] \\
 & 3.1 & Lem. & \emph{Nostalgia is fitting when it tracks a\ldots} & 2 & \texttt{WA} & [5] \\
 & 3.2 & Thm. & \emph{Nostalgia is rational in moderation: an episode\ldots} & 2 & \texttt{WA} & [5] \\
 & 1 & Obj. & The past is gone & 2 & \texttt{} & --- \\
 & 2 & Obj. & Selective memory & 3 & \texttt{} & --- \\
 & 3 & Obj. & Uselessness & 3 & \texttt{} & --- \\
EXI-15-a & 2.1 & Def. & Hope & 1 & \texttt{I} & --- \\
 & P1 & Pr. & Credence is evidential & 2 & \texttt{I} & --- \\
 & P2 & Pr. & The stance is controlled & 2 & \texttt{I} & --- \\
 & 2.2 & Def. & Virtue & 2 & \texttt{AD} & --- \\
 & 3.1 & Def. & Proportioned hope & 2 & \texttt{WI} & --- \\
 & 3.2 & Lem. & \emph{Hope whose stance is insensitive to probability\ldots} & 2 & \texttt{} & --- \\
 & 3.3 & Lem. & \emph{Hope whose stance is fixed by probability\ldots} & 2 & \texttt{} & --- \\
 & 3.4 & Thm. & \emph{Hope is a virtue when it is\ldots} & 2 & \texttt{} & --- \\
 & 1 & Obj. & Stoic & 3 & \texttt{} & --- \\
 & 2 & Obj. & Self-deception & 3 & \texttt{} & --- \\
 & 3 & Obj. & Hope against hope & 3 & \texttt{} & --- \\
EXI-15-b & 2.1 & Prop. & Near certainty & 1 & \texttt{I} & --- \\
 & 2.2 & Prop. & Remote possibility & 2 & \texttt{WI} & --- \\
 & 3.1 & Prop. & \emph{When the agent's effort is among the\ldots} & 2 & \texttt{I} & --- \\
EXI-16-a & 2.1 & Def. & Reasons for love & 2 & \texttt{} & --- \\
 & 2.2 & Def. & Qualitative and relational reasons & 2 & \texttt{} & --- \\
 & 2.3 & Def. & Qualitative similarity & 2 & \texttt{D} & --- \\
 & 2.4 & Def. & Transferability & 2 & \texttt{} & --- \\
 & P1 & Pr. & Answerability & 2 & \texttt{O} & --- \\
 & P2 & Pr. & Generality & 2 & \texttt{OD} & --- \\
 & P3 & Pr. & Particularity of history & 2 & \texttt{} & --- \\
 & 3.1 & Lem. & \emph{Love responds to reasons.\ldots} & 2 & \texttt{} & --- \\
 & 3.2 & Lem. & \emph{If all reasons for 's love of\ldots} & 3 & \texttt{} & --- \\
 & 3.3 & Lem. & \emph{If the reasons for 's love of\ldots} & 3 & \texttt{} & --- \\
 & 3.4 & Lem. & \emph{Every established love answers to relational reasons.\ldots} & 3 & \texttt{} & --- \\
 & 3.5 & Thm. & \emph{Love responds to reasons, but these reasons\ldots} & 3 & \texttt{} & --- \\
 & 1 & Obj. & Wrong object & 4 & \texttt{W} & --- \\
 & 2 & Obj. & Duplication & 4 & \texttt{} & --- \\
 & 3 & Obj. & Arbitrariness & 4 & \texttt{W} & --- \\
 & 4 & Obj. & Universal value & 5 & \texttt{WA} & [5] \\
EXI-16-b & 1.1 & Prop. & Trading up & 1 & \texttt{D} & --- \\
 & 1.2 & Prop. & Duplication & 1 & \texttt{} & --- \\
 & 1.3 & Prop. & After loss & 2 & \texttt{} & --- \\
EXI-17-a & 2.1 & Def. & Active engagement & 2 & \texttt{EPI} & --- \\
 & 2.2 & Def. & Worthwhile & 2 & \texttt{EWOPI} & [3], [4] \\
 & 2.3 & Def. & Achievement & 2 & \texttt{PI} & --- \\
 & P1 & Pr. & Fitting fulfilment & 2 & \texttt{MEWDL} & [1], [2], [2.3], [10] \\
 & P2 & Pr. & Contact with reality & 2 & \texttt{EPI} & --- \\
 & 2.4 & Ass. & Reflective judgment & 2 & \texttt{MWL} & [1] \\
 & 3.1 & Lem. & Failure is compatible with meaning & 3 & \texttt{ML} & [1] \\
 & 3.2 & Lem. & Meaning without product & 3 & \texttt{MEDL} & [1], [2], [2.3], [10] \\
 & 3.3 & Thm. & \emph{A meaningful life requires active engagement with\ldots} & 3 & \texttt{MEWL} & [1], [2], [10] \\
 & 3.4 & Cor. & Achievement is neither necessary nor sufficient & 3 & \texttt{MEWP} & [2] \\
 & 3.5 & Rem. & Degree & 3 & \texttt{MEWDL} & [1], [2], [2.3], [10] \\
 & 3.6 & Cor. & Availability & 4 & \texttt{MA} & [5] \\
 & 1 & Obj. & Everything is an achievement & 4 & \texttt{MPL} & [1] \\
 & 2 & Obj. & Making a difference & 4 & \texttt{MWL} & [1] \\
 & 3 & Obj. & The experience machine & 4 & \texttt{MWL} & [1] \\
 & 4 & Obj. & Sisyphus & 4 & \texttt{MEL} & [1], [10] \\
EXI-18-a & 2.1 & Def. & Psychological necessity & 1 & \texttt{I} & --- \\
 & 2.2 & Prop. & Partial necessity & 2 & \texttt{I} & --- \\
 & 2.3 & Thm. & \emph{That an illusion is psychologically necessary does\ldots} & 2 & \texttt{} & --- \\
 & 2.4 & Cor. & No license & 2 & \texttt{} & --- \\
 & 2.5 & Rem. & Substitutes & 2 & \texttt{OPD} & [3] \\
 & 1 & Obj. & Pragmatic justification & 2 & \texttt{I} & --- \\
 & 2 & Obj. & Sadder but wiser & 3 & \texttt{} & --- \\
 & 3 & Obj. & Meaning needs illusion & 3 & \texttt{MPS} & [11] \\
EXI-18-b & 1.1 & Prop. & \emph{(i) Whether a given positive illusion is\ldots} & 1 & \texttt{I} & --- \\
EXI-19-a & 2.1 & Def. & Meaning in itself & 1 & \texttt{MWOL} & [1], [3] \\
 & 2.2 & Def. & Shape & 2 & \texttt{WL} & --- \\
 & P1 & Pr. & Shape-dependence & 2 & \texttt{WL} & --- \\
 & 2.3 & Ass. & Bounds & 2 & \texttt{} & --- \\
 & 3.1 & Prop. & \emph{Death has no meaning in itself.\ldots} & 2 & \texttt{M} & --- \\
 & 3.2 & Thm. & \emph{Death gives a life a shape by\ldots} & 2 & \texttt{L} & --- \\
 & 3.3 & Cor. & \emph{Theorem is consistent with the claim of\ldots} & 2 & \texttt{W} & --- \\
 & 1 & Obj. & Other limits & 3 & \texttt{PL} & --- \\
 & 2 & Obj. & Death as good & 3 & \texttt{W} & --- \\
 & 3 & Obj. & Ruined shapes & 3 & \texttt{L} & --- \\
EXI-19-b & 1.1 & Prop. & Not a source & 1 & \texttt{MEWL} & [1], [2], [10] \\
 & 1.2 & Prop. & Whole-life meaning & 1 & \texttt{ML} & [1] \\
 & 1.3 & Prop. & Persistence of engagement & 1 & \texttt{MEL} & [1], [10] \\
\end{longtable}
}

\section{Testimonia}\label{app:test}

\subsection*{Preface}

The testimonia are every sentence in the other manuscripts of the discipline, outside the family EXI-01 itself, that cites EXI-01 by identifier: \cnTestimonia{} sentences from \cnTestimoniaMss{} manuscripts. They are extracted by the program of Appendix~\ref{app:conc} from the same text, with cross-references and citations resolved, and are given in the order of the manuscripts. In the commentary tradition, testimonia record what later texts say an earlier text said; here they record how the discipline understands EXI-01, and they are evidence for the reading of its lemmata.

Four testimonia are discussed in the Commentary. The sentence of EXI-04-a, Section~4, that quotes the last clause of Remark~5 is the only testimonium that reports a sentence of EXI-01-a rather than its conclusion, and the comment on \Lm{11.4} shows that its reading agrees with Proposition~\ref{p:vind}. The sentence of the abstract of EXI-06-b that describes meaning as engagement with what has worth \enquote{independently of anyone's attitudes} strengthens Assumption~2 beyond its text, and the comment on \Lm{3.4} shows that EXI-06-b's argument does not need the strengthening. The sentence of EXI-09-b that a life can be less meaningful \enquote{in the sense of EXI-01} than it might have been is the discipline's clearest use of the degree of Definition~1 and of the left-to-right direction that Remark~5 also uses (\Lm{10}). The sentences of EXI-11-a and EXI-20-a report the availability claim in the words of the family conclusion, \enquote{available to everyone}, and so confirm the reading of \enquote{everyone} at \Lm{4.2}.

{\footnotesize% Generated by gen/generate.py; do not edit.
\begin{longtable}{@{}l l >{\raggedright\arraybackslash}p{0.78\textwidth}@{}}
\toprule Ms & \S & Sentence \\ \midrule \endhead
\bottomrule \endfoot
EXI-03-a & 1 & It identifies the categorical desires that give life its point with desires for projects the person takes to be worthwhile, in line with the account of meaning in EXI-01; and it replaces Williams's appeal to identity of character with the account of what matters in survival defended in PID-06. \\[2pt]
EXI-03-a & 2 & EXI-01 argues that a life is meaningful to the extent that its subject engages with projects of objective worth, and EXI-17 that such engagement, not achievement, is what meaning requires. \\[2pt]
EXI-04-a & 3 & EXI-01 argues that life is meaningful to the extent that one engages with projects of objective worth, and we accept that conclusion. \\[2pt]
EXI-04-a & 4 & If some projects are worthwhile independently of our attitudes, as EXI-01 holds, then the seriousness with which we pursue them is justified, and the life of someone engaged in them is not absurd. \\[2pt]
EXI-04-a & 4 & This is the sense in which, as the remark at the end of EXI-01-a puts it, the assumption of objective worth denies the premise of the sceptical challenge without refuting it. \\[2pt]
EXI-04-a & 5 & The argument does not show that any particular life is meaningful or worthwhile; that question is answered by EXI-01 and EXI-17, and absurdity, as defined here, is shared by meaningful and meaningless lives alike. \\[2pt]
EXI-04-b & 3 & Our ends include worthwhile objects, and engagement with them confers meaning (EXI-01). \\[2pt]
EXI-06-b & Abs & Meaning, on the account of EXI-01, consists in engagement with what has worth independently of anyone's attitudes, and the universe's want of attitudes toward us takes nothing from that worth. \\[2pt]
EXI-06-b & 2 & On the account of EXI-01, a life is meaningful to the extent that its subject engages with projects of objective worth, whose value does not depend on the attitudes of those engaged in them (Definition 1 and Assumption 2 of EXI-01-a). \\[2pt]
EXI-06-b & 3 & The present point is conditional: if meaning consists in engagement with what has objective worth, as EXI-01 holds, then it does not require any being's care, the universe's or another's. \\[2pt]
EXI-06-b & 4 & It leaves the question of meaning where EXI-01 places it: in what we engage with, and how. \\[2pt]
EXI-07-a & 3 & It fits also the account of meaning in EXI-01 and EXI-17, on which meaning comes from active engagement with what is worthwhile. \\[2pt]
EXI-07-a & 5 & It also does not establish which other ends a person should have; that is the work of the account of the good (ETH-01) and of meaning (EXI-01, EXI-17). \\[2pt]
EXI-09-a & 5 & Principle P2 says that what a life is for is fixed by its agent's ends; it does not say that every end is as good as any other, and is consistent with EXI-01, on which a life is meaningful to the extent that its projects are of objective worth. \\[2pt]
EXI-09-b & 2 & A life can be imprudent, wicked or, in the sense of EXI-01, less meaningful than it might have been, because it engages little with what is worthwhile. \\[2pt]
EXI-11-a & 2 & EXI-01 adds that meaning so understood is available to everyone. \\[2pt]
EXI-17-a & 3 & Corollary 3.6 supports, without by itself establishing, the conclusion of EXI-01 that meaning in this sense is available to everyone. \\[2pt]
EXI-19-b & 1 & On the account of EXI-01 and EXI-17, a life is meaningful to the extent that its subject actively engages with what is worthwhile. \\[2pt]
EXI-19-b & 1 & So the limit adds neither an object of worth nor worth to an object, and by the account of EXI-01 it is not a source of meaning. \\[2pt]
EXI-20-a & 1 & EXI-01 adds that meaning in this sense is available to everyone. \\[2pt]
\end{longtable}
}
